\chapter[Propret\'e cohomologique des faisceaux d'ensembles...]{Propret\'e cohomologique\\ des faisceaux d'ensembles et\\ des faisceaux de groupes non commutatifs}
\label{XIII}

\begin{center}
{par Mme M.\ \textsc{Raynaud}\footnote{D'apr\`es des notes in\'edites de
A.\, Grothendieck.}}
\end{center}
\vspace*{1cm}

Cet expos\'e se propose d'utiliser la cohomologie \'etale pour
g\'en\'eraliser certains r\'esultats de \ref{IX} et \ref{X}. Il montre aussi
comment on peut \'etendre aux faisceaux en groupes non
n\'ecessairement commutatifs les r\'esultats de~SGA~5~II qui ont
encore un sens pour de tels faisceaux. On suppose connues les notions
de cohomologie \'etale expos\'ees dans~SGA~4.

Le r\'esultat principal \eqref{XIII.2.4} donne un exemple important
de morphisme non propre $f\colon~U\to~S$, qui soit \og cohomologiquement
propre en dimension~$\leq1$\fg, c'est-\`a-dire tel que, pour
certains faisceaux en groupes $F$ sur~$U$ (au sens de la topologie
\'etale), la formation de~$f_*F$ et $\R^1f_*F$ commute \`a tout
changement de base~$S'\to~S$. Cette propri\'et\'e est en effet
satisfaite par l'ouvert $U$ d'un sch\'ema $X$ propre sur~$S$,
compl\'ementaire d'un diviseur $D$ \`a croisements normaux
relativement \`a~$S$, du moins si l'on impose \`a~$F$ d'\^etre
constant fini, d'ordre premier aux caract\'eristiques
r\'esiduelles de~$S$. Si l'on ne suppose plus $F$ d'ordre premier
aux caract\'eristiques r\'esiduelles de~$S$, on a un r\'esultat
analogue en rempla\c cant $\R^1f_*F$ par le sous-faisceau
$\R^1_{\tame}f_*F$ obtenu en se bornant \`a consid\'erer les torseurs
sous~$F$ \og mod\'er\'ement ramifi\'es sur~$X$ relativement
\`a~$S$\fg. En particulier cela permet de montrer que le groupe
fondamental mod\'er\'ement ramifi\'e d'une courbe alg\'ebrique
propre et lisse sur un corps s\'eparablement clos, priv\'ee d'un
nombre fini de points ferm\'es, est topologiquement de type fini
\eqref{XIII.2.12}.

Le
\no \ref{XIII.4} est consacr\'e \`a la suite exacte d'homotopie et \`a
la formule de K\"unneth.

Enfin un appendice donne des variantes utiles du lemme d'Abhyankar
d\'emontr\'e dans X.\ref{X.3.6}.

\setcounter{section}{-1}

\section{Rappels sur la th\'eorie des champs}
\label{XIII.0}

Nous utiliserons dans ce qui suit la th\'eorie des champs
expos\'ee dans \cite{XIII.1} et \cite{XIII.2}. Nous nous bornons au cas du site
\'etale d'un sch\'ema. \'Etant donn\'e un sch\'ema $X$, notons
$X_\et$ le site \'etale de~$X$. Rappelons qu'un
champ $F$
\index{champ|hyperpage}sur~$X$ est une cat\'egorie fibr\'ee au-dessus de~$X_\et$ telle
que, pour tout sch\'ema $X'$ \'etale sur~$X$ et pour tout couple
d'objets
$x$, $y$ de la fibre $F_{X'}$, le pr\'efaisceau
$\SheafHom_{X'}(x,y)$ soit un faisceau, et telle que, pour tout
morphisme \'etale surjectif $X'' \to X'$, tout objet de~$F_{X''}$
muni d'une donn\'ee de descente relativement \`a $X'' \to X'$ soit
image inverse d'un objet de~$F_{X'}$.

On note $F(X')$ la cat\'egorie des sections cart\'esiennes de
$F/X'$. Plus g\'en\'eralement, si $\Sch_X$ est la cat\'egorie
des sch\'emas au-dessus de~$X$ munie de la topologie \'etale, le
champ $F$ peut s'\'etendre en un champ $\cal{F}$ sur~$\Sch_X$ et,
pour tout morphisme $f\colon X' \to X$, on note encore $F(X')$ la
cat\'egorie des sections cart\'esiennes de ce champ $\cal{F}$
au-dessus de~$X'$.

Une gerbe
\index{gerbe|hyperpage}est un champ tel que, pour tout sch\'ema $X'$ \'etale
sur~$X$ et pour tout couple d'objets $x$, $y$ de~$F_{X'}$, tout
morphisme de~$x$ dans $y$ soit un isomorphisme, que $x$ et $y$ soient
localement isomorphes, et tel que l'ensemble des objets $X'$ de
$X_\et$ tels que $F_{X'}$ soit non vide est un raffinement de
$X_\et$. Par exemple le champ des torseurs sous un faisceau en groupes
est une gerbe qui, de plus, a une section cart\'esienne.
R\'eciproquement une gerbe qui a une section, \ie telle qu'il
existe un objet $x$ de~$F_X$, est \'equivalente au champ des
torseurs sous le faisceau
en groupes $\SheafAut_X(x)$.

On a une notion \'evidente de sous-gerbe et de sous-gerbe maximale
d'un champ~$F$. \'Etant donn\'e une section cart\'esienne $x$ de
$F(X)$, il existe une unique sous-gerbe maximale $G_x$ de~$F$ telle
que $x$ se factorise \`a travers~$G_x$. On appelle $G_x$ la
sous-gerbe engendr\'ee par~$x$; c'est par d\'efinition une gerbe
triviale.
\index{gerbe triviale|hyperpage}Le pr\'efaisceau $SF$
\label{indnot:mb}d\'efini par
$$
SF(X') = \{ \text{sous-gerbes maximales de }F|X' \}
$$
est un faisceau appel\'e, faisceau des sous-gerbes maximales
de~$F$.
\index{faisceau des sous-gerbes maximales|hyperpage}Soit $O$ le pr\'efaisceau d\'efini par
$$
O(X') = \{ \text{classes d'objets de }F_{X'}\text{ mod. isomorphisme}
\} \text{.}
$$
En associant \`a tout objet $x$ de~$F_{X'}$ la sous-gerbe maximale
de~$F|X'$, engendr\'ee par~$x$, on obtient un morphisme
$$
O \to SF \text{;}
$$
d'apr\`es \cite[III 2.1.4]{XIII.2}, ce morphisme fait de~$SF$ un faisceau
associ\'e \`a~$O$.

Un champ $F$ est dit \emph{constructible}
\index{constructible (champ)|hyperpage}(\emph{\resp ind-$\LL$-fini},
\index{champ ind-fini (\resp ind-$l$-fini)|hyperpage}$\LL$ \'etant un ensemble de nombres premiers) si, pour tout
sch\'ema $X'$ \'etale sur~$X$ et pour tout objet $x$ de~$F_{X'}$,
il en est ainsi du faisceau $\SheafAut_{X'}(x)$ \cite[VII 2.2.1]{XIII.2}. On dit
qu'un champ est $1$-constructible
\index{champ $1$-constructible (\resp constructible)|hyperpage}s'il est constructible et si le faisceau des sous-gerbes maximales est
constructible.

\section{Propret\'e cohomologique}
\label{XIII.1}
\setcounter{subsection}{-1}
\subsection{}
\label{XIII.1.0}
Soient $S$ un sch\'ema, $f \colon X \to Y$ un morphisme de
$S$-sch\'emas. Si $S'$ est un $S$-sch\'ema, on consid\`ere le
diagramme suivant, dont tous les carr\'es sont cart\'esiens:
\begin{equation*}
\label{eq:XIII.1.0.1}
\tag{\thesubsection.1}
\begin{array}{c}
\xymatrix{{X}\ar[d]_f&{X'}\ar[l]_h\ar[d]^{f'}\\{Y}\ar[d]&{Y'}\ar[l]_g\ar[d]\\{S}&{\,S'.}\ar[l]}
\end{array}
\end{equation*}

Si $Y_1$ est un sch\'ema \'etale au-dessus de~$Y$, on pose $X_1 =
X \times_Y Y_1$, $Y_1' = Y' \times_Y Y_1$, et on consid\`ere le
carr\'e cart\'esien
\begin{equation*}
\label{eq:XIII.1.0.2}
\tag{\thesubsection.2}
\begin{array}{c}
\xymatrix{{X_1}\ar[d]_{f_1}&{X_1'}\ar[l]_{h_1}\ar[d]^{f_1'}\\{Y_1}&{\,Y_1'.}\ar[l]^{g_1}}
\end{array}
\end{equation*}

\begin{definition}
\label{XIII.1.1}
Soit $F$ un champ sur~$X$. On dit que $(F,f)$ est
\emph{cohomologiquement propre relativement à~$S$ en dimension
$\leqslant -1$}
\index{cohomologiquement propre|hyperpage}\index{propret\'e cohomologique|hyperpage}(\resp en dimension $\leqslant 0$, \resp en dimension $\leqslant 1$)
si, pour tout $S$-sch\'ema $S'$, le foncteur canonique (d\'efini
de fa\c con \'evidente par la propri\'et\'e universelle de
l'image inverse de champs):
$$
g^* f_* F \to f'_* h^* F \qquad\text{ (\cf 1.0.1)}
$$
est fid\`ele (\resp pleinement fid\`ele, \resp une \'equivalence
de cat\'egories).
\end{definition}

S'il n'y a pas de confusion possible sur~$S$, en particulier si $S=Y$,
on dit cohomologiquement propre au lieu de cohomologiquement propre
relativement \`a~$S$.

\subsection{}
\label{XIII.1.2}
Soit $F$ un faisceau d'ensembles sur~$X$; soit $\Phi$ le champ en
cat\'egories discr\`etes associ\'e \`a $F$, \ie le champ dont
la fibre au-dessus de tout sch\'ema $X_1$ \'etale sur~$X$ est la
cat\'egorie discr\`ete ayant pour ensemble d'objets $F(X_1)$. On
dit que $(F,f)$ est cohomologiquement propre relativement \`a~$S$ en
dimension $\leqslant -1$ (\resp en dimension $\leqslant 0$) si $(\Phi,
f)$ est cohomologiquement propre relativement \`a~$S$ en dimension
$\leqslant 0$ (\resp en dimension $\leqslant 1$).

Le
morphisme canonique
\begin{equation*}
\label{eq:XIII.1.2.1}
\tag{\thesubsection.1} { g^* f_* F \to f'_* h^* F }
\end{equation*}
donne par passage aux champs en cat\'egories discr\`etes
associ\'ees le morphisme canonique
$$
g^* f_* \Phi \to f'_* h^* \Phi \text{.}
$$
Par suite dire que $(F,f)$ est cohomologiquement propre relativement
\`a~$S$ en dimension $\leqslant -1$ (\resp en dimension $\leqslant
0$) \'equivaut \`a dire que, pour tout $S$-sch\'ema $S'$, le
morphisme \eqref{eq:XIII.1.2.1} est injectif (\resp bijectif).

\subsection{}
\label{XIII.1.3}
Soit $F$ un faisceau en groupes sur~$X$ et $\Phi$ le champ des
torseurs sur~$X$ de groupe $F$ \cite[II 2.3.2]{XIII.1}. On dit que $(F,f)$ est
cohomologiquement propre relativement \`a~$S$ en dimension
$\leqslant -1$ (\resp $\leqslant 0$, \resp $\leqslant 1$) si
$(\Phi,f)$ est cohomologiquement propre relativement \`a~$S$ en
dimension $\leqslant -1$ (\resp $\leqslant 0$, \resp $\leqslant
1$). La condition de propret\'e cohomologique peut s'expliciter
comme suit.

\begin{subproposition}
\label{XIII.1.3.1}
Les notations sont celles de \eqref{eq:XIII.1.0.1} et
\eqref{eq:XIII.1.0.2}. Soit $F$ un faisceau en groupes sur~$X$. On
d\'esigne par $F'$ (\resp $F_1$, \resp $F_1'$, etc.), l'image
inverse de~$F$ sur~$X'$ (\resp sur~$X_1$, \resp sur~$X_1'$,
etc.). Alors les conditions suivantes sont \'equivalentes:
\begin{enumerate}
\item[(i)] $(F,f)$ est cohomologiquement propre relativement \`a~$S$
en dimension $\leqslant -1$ (\resp $\leqslant 0$, \resp $\leqslant
1$).
\item[(ii)] Pour tout morphisme $S' \to S$, pour tout sch\'ema $Y_1$
\'etale au-dessus de~$Y$, et pour tout torseur $P$ sur~$X_1$ de
groupe $F_1$, si $\leftexp{P\mkern-4mu}{F}_1$ d\'esigne le groupe tordu de
$F_1$ par $P$ \textup{\cite[II 4.1.2.3]{XIII.1}}, le morphisme canonique
$$
a_0 \colon g_1^*(f_{1 *}(\leftexp{P\mkern-4mu}{F}_1)) \to f'_{1
*}(\leftexp{P'\mkern-6mu}{F}'_1)
$$
est injectif (\resp $a_0$ est bijectif et le morphisme canonique
$$
a_1 \colon g^*(\R^1 f_* F) \to \R^1 f'_* F'
$$
est
injectif, \resp $a_0$ et $a_1$ sont bijectifs).
\item[(ii~bis)] Pour tout morphisme $S' \to S$, pour tout sch\'ema
$Y_1$ \'etale au-dessus de~$Y$, pour tout torseur $P$ sur~$X_1$ de
groupe $F_1$, et pour tout torseur $R$ sous $\leftexp{P\mkern-4mu}{F}_1$, le
morphisme canonique
$$
\alpha_0 \colon g_1^*(f_{1 *}R) \to f'_{1 *} R'
$$
est injectif (\resp $\alpha_0$ est bijectif, \resp les morphismes
$\alpha_0$ et
$$
\alpha_1 \colon g_1^*(\R^1 f_{1 *}(\leftexp{P\mkern-4mu}{F}_1)) \to \R^1
f'_{1 *}(\leftexp{P'\mkern-6mu}{F}'_1)
$$
sont bijectifs).
\end{enumerate}
\end{subproposition}

\subsubsection*{D\'emonstration}
(i)$\To$(ii~bis). D'apr\`es \cite[II 4.2.5]{XIII.1} tout torseur $R$
de groupe $\leftexp{P\mkern-4mu}{F}_1$ est de la forme $R = Q
\overset{F_1}{\wedge} P^\circ$, o\`u $Q$ est un torseur de groupe
$F_1$ et $P^\circ$ l'oppos\'e de~$P$. On a alors $R' \simeq Q'
\overset{F'_1}{\wedge} {P'}^\circ$. Soit $\Phi$ le champ des torseurs
sous $F$ et soient $x$, $y$ (\resp~$x'$,~$y'$) les objets de la
cat\'egorie fibre $(g^* f_* \Phi)_{Y_1'}$ (\resp $(f'_*
\Phi')_{Y_1'}$) associ\'es \`a $P$, $Q$ (\resp $P'$,~$Q'$). On a
la relation
$$
Q\overset{F_1}{\wedge} P^\circ \simeq \SheafHom_{F_1}(P,Q) \text{,}
$$
et il en r\'esulte que l'on a des isomorphismes canoniques
$$
\SheafHom_{Y'_1}(x,y) \simeq g_1^* f_{1 *}(Q
\overset{F_1}{\wedge} P^\circ) \text{,}\quad \SheafHom_{Y_1'}(x',y')
\simeq f'_{1 *}(Q' \overset{F'_1}{\wedge}{P'}^\circ)\text{.}
$$
Par suite le morphisme $\alpha_0$ s'identifie au morphisme
$$
\SheafHom_{Y'_1}(x,y) \to \SheafHom_{Y'_1}(x',y') \text{,}
$$
d'o\`u il r\'esulte que, si $(F,f)$ est cohomologiquement propre
relativement \`a~$S$ en dimension $\leqslant -1$, $\alpha_0$ est
injectif et que, si $(F,f)$ est cohomologiquement propre en dimension
$\leqslant 0$, $\alpha_0$ est bijectif.

Supposons maintenant que $(F,f)$ soit cohomologiquement propre
relativement \`a~$S$ en dimension $\leqslant 1$, \ie que le
morphisme canonique
$$
\varphi \colon g^* f_* \Phi \to f'_* \Phi'
$$
soit une \'equivalence. Soit $G$ le faisceau des sous-gerbes
maximales du champ $f_* \Phi$ \hbox{\cite[III~2.1.8]{XIII.1}}; on a alors un
isomorphisme $G \simeq \R^1 f_* F$. Comme $g^* G$ est le
faisceau des sous-gerbes maximales de~$g^* f_* \Phi$ \cite[III
2.1.5.5]{XIII.2}, le morphisme $\alpha_1$ est obtenu \`a partir de
$\varphi|Y'_1$ en prenant les faisceaux des sous-gerbes maximales,
donc est un isomorphisme.

(ii~bis)$\To$(ii). Il suffit de montrer que, si les
morphismes $\alpha_0$ sont bijectifs, alors les morphismes $a_1$ sont
injectifs. Soient $Y'_1$ un sch\'ema \'etale au-dessus de~$Y'$,
$s$ et $t$ deux \'el\'ements de~$g^*(\R^1 f_* F)(Y'_1)$
ayant m\^eme image dans $\R^1 f'_* F'(Y'_1)$ et montrons que
l'on a $s=t$. L'assertion est locale pour la topologie \'etale de
$Y'_1$ et, compte tenu de la d\'efinition de l'image inverse
$g^*(\R^1 f_* F)$, on peut supposer que $Y'_1$ est l'image
inverse d'un sch\'ema $Y_1$ \'etale au-dessus de~$Y$ et que $s$ et
$t$ proviennent de torseurs $P$ et $Q$ sur~$X_1$. L'hypoth\`ese
faite sur~$s$ et $t$ signifie alors que les images inverses $P'$ et
$Q'$ de~$P$ et $Q$ sur~$X'_1$ sont isomorphes localement pour la
topologie \'etale de~$Y'_1$. Si l'on pose $R =
\SheafHom_{F_1}(P,Q)$, le fait que le morphisme
$$
g_1^* f_{1 *} R \to f'_{1 *} R'
$$
soit bijectif prouve que $P$ et $Q$ sont isomorphes localement pour la
topologie \'etale de~$Y_1$, donc que l'on a $s=t$.

(ii)$\To$(i). Pour prouver que $\varphi$ est fid\`ele
(\resp pleinement fid\`ele), il suffit de montrer que, si $Y_1$ est
un sch\'ema \'etale sur~$Y$, si $P$, $Q$ sont deux torseurs sur~$X_1$ de groupe~$F_1$, si $x$, $y$ (\resp $x'$, $y'$) sont les objets
de~$(g^* f_* \Phi)_{Y'_1}$ (\resp $(f'_* \Phi')_{Y'_1}$)
associ\'es \`a $P$, $Q$ (\resp $P'$, $Q'$), le morphisme
$$
a \colon \Hom(x,y) \to \Hom(x',y')
$$
est injectif (\resp bijectif). Or $a$ s'identifie au morphisme
canonique
$$
\H^0(Y'_1,g^*_1 f_{1 *} (Q \overset{F_1}{\wedge} P^\circ)) \to
\H^0(Y'_1, f'_{1 *}(Q' \overset{F'_1}{\wedge} {P'}^\circ))\text{.}
$$
Si
l'on a $\Hom(x,y) \neq \emptyset$, alors $Q \overset{F_1}{\wedge}
P^\circ$ est un torseur sous $\leftexp{P\mkern-4mu}{F}_1$ localement trivial sur~$Y_1$; il en r\'esulte que $f_{1 *}(Q \overset{F_1}{\wedge}
P^\circ)$ est un torseur sous $f_{1*}(\leftexp{P\mkern-4mu}{F}_1)$, et que
$g^*_1 f_{1 *} (Q \overset{F_1}{\wedge} P^\circ)$ est un torseur
trivial. Le morphisme $a$ s'identifie alors au morphisme canonique
$$
\H^0(Y_1, g^*_1 f_{1 *}(\leftexp{P\mkern-4mu}{F}_1)) \to \H^0 (Y'_1,
f'_{1 *}(\leftexp{P'\mkern-6mu}{F}'_1)) \text{.}
$$
Il en est de m\^eme si l'on a $\Hom(x',y') \neq \emptyset$ et si
$a_1$ est injectif car alors $Q' \overset{F'_1}{\wedge} {P'}^\circ$
est trivial, et il r\'esulte de l'injectivit\'e de~$a_1$ que $P$
et $Q$ sont localement isomorphes sur~$Y_1$. On en conclut que, si
$a_0$ est injectif (\resp si $a_0$ est bijectif et $a_1$ injectif),
$\varphi$ est fid\`ele (\resp pleinement fid\`ele).

Il reste \`a montrer que, si $a_0$ et $a_1$ sont bijectifs, le
foncteur $\varphi$ est essentiellement surjectif. Soient $Y''$ un
sch\'ema \'etale au-dessus de~$Y'$, $X'' = X' \times_{Y'} Y''$ et
soit $P''$ un torseur sur~$X''$ de groupe $F'' = F'|X''$. On va
montrer qu'il existe un \'el\'ement~$x$ de~$(g^* f_*
\Phi)_{Y''}$ dont l'image dans $(f'_* \Phi')_{Y''}$ est isomorphe
\`a $P''$. Soit $p''$ la classe de~$P''$. Du fait que $a_1$ est
surjectif r\'esulte que l'on peut trouver un morphisme \'etale
surjectif $Y''_1 \to Y''$, un morphisme \'etale $Y_1 \to Y$ tel que
l'on ait un morphisme $Y''_1 \to Y'_1$ et un torseur $P_1$ sur~$X_1$
de groupe $F_1$ dont l'image inverse $P''_1$ sur~$X''_1$ soit
isomorphe \`a l'image inverse de~$P''$. Utilisant le fait que
$\varphi$ est pleinement fid\`ele, on voit que l'objet $x_1$ de
$(g^* f_* \Phi)_{Y''_1}$ qui correspond \`a $P''_1$ est muni
d'une donn\'ee de descente relativement \`a $Y''_1 \to Y''$, donc
provient d'un \'el\'ement $x$ de~$(g^* f_*
\Phi)_{Y''}$. Comme l'image de~$x$ dans $(f'_* \Phi')_{Y''}$ est
$P''$, ceci prouve que $\varphi$ est essentiellement surjectif et
ach\`eve la d\'emonstration.

\begin{exemple}
\label{XIII.1.4}
Soit $f \colon X \to Y$ un morphisme propre. Il r\'esulte de \cite[VII
2.2.2]{XIII.2} que, pour tout champ ind-fini $F$ sur~$X$, le couple $(F,f)$
est cohomologiquement propre (relativement \`a $Y$) en dimension
$\leqslant 1$. En particulier, pour tout faisceau d'ensembles
(\resp tout faisceau de groupes, \resp tout faisceau en groupes
ind-fini) $F$ sur~$X$, $(F,f)$ est cohomologiquement
propre en dimension $\leqslant 0$ (\resp en dimension $\leqslant 0$,
\resp en dimension $\leqslant 1$).
\end{exemple}

\begin{remarques}
\label{XIII.1.5}
a) Soit $F$ un faisceau en groupes sur~$X$ tel que $(F,f)$ soit
cohomologiquement propre relativement \`a~$S$ en dimension
$\leqslant -1$ (\resp $\leqslant 0$). Si l'on consid\`ere $F$ comme
faisceau d'ensembles, $(F,f)$ est cohomologiquement propre
relativement \`a~$S$ en dimension $\leqslant -1$ (\resp $\leqslant
0$), mais la r\'eciproque est fausse.

Soit par exemple $Y$ le spectre d'un anneau de valuation discr\`ete
strictement local de point ferm\'e $t$, de point g\'en\'erique
$s$, $f \colon X \to Y$ un sch\'ema non vide sur~$Y$ dont la fibre
ferm\'ee est vide, $F$ un faisceau en groupes constant non trivial
sur~$X$ et $P$ un torseur sous $F$ tel que l'on ait $\H^0(X_s,
\leftexp{P\mkern-4mu}{F}|X_s)=1$. Alors $(\leftexp{P\mkern-4mu}{F},f)$ est
cohomologiquement propre relativement \`a $Y$ en dimension
$\leqslant -1$ lorsque l'on consid\`ere $\leftexp{P\mkern-4mu}{F}$ comme
faisceau d'ensembles. Si l'on consid\`ere $\leftexp{P\mkern-4mu}{F}$ comme faisceau
en groupes, on a un isomorphisme ${\vphantom{F}}^{P^\circ\mkern-7mu}(\leftexp{P\mkern-4mu}{F})
\simeq F$; comme le morphisme canonique
$$
\H^0(X,F)\to\H^0(X_t,F|X_t)=1
$$
n'est pas injectif, ceci prouve que $(\leftexp{P\mkern-4mu}{F},f)$ n'est pas
cohomologiquement propre relativement \`a $Y$ en dimension
$\leqslant -1$.


b) Supposons $f$ coh\'erent (\ie quasi-compact et
quasi-s\'epar\'e). Soit $F$ un champ sur~$X$. Pour tout point
g\'eom\'etrique $\overline{y}$ de~$Y'$, on note $\overline{Y}$
(\resp $\overline{Y'}$) le localis\'e strict de~$Y$ (\resp $Y'$) en
$\overline{y}$, et on pose $\overline{X} = X \times_Y \overline{Y}$,
$\overline{X'} = X' \times_{Y'} \overline{Y'}$, etc. Pour que $(F,f)$
soit cohomologiquement propre relativement \`a~$S$ en dimension
$\leqslant -1$ (\resp $\leqslant 0$, \resp $\leqslant 1$), il faut et
il suffit que, pour tout $S$-sch\'ema $S'$ et pour tout point
g\'eom\'etrique $\overline{y}$ de~$Y'$, le foncteur canonique
$$
\overline{F}(\overline{X})\to\overline{F'}(\overline{X'})
$$
soit fid\`ele (\resp pleinement fid\`ele, \resp une
\'equivalence).

En
effet, si $S'$ est un $S$-sch\'ema, pour que le foncteur
$$
g^* f_* F \to f'_* F'
$$
soit fid\`ele (\resp pleinement fid\`ele, \resp une
\'equivalence), il faut et il suffit qu'il en soit ainsi du foncteur
induit sur les fibres aux diff\'erents points g\'eom\'etriques
$\overline{y}'$ de~$\overline{Y'}$ \cite[III 2.1.5.9]{XIII.2}. L'assertion
r\'esulte donc du calcul des fibres g\'eom\'etriques de l'image
directe d'un champ par un morphisme coh\'erent \cite[VII 2.1.5]{XIII.2}.


c) Soit $F$ un champ sur~$X$. Le fait que $(F,f)$ soit
cohomologiquement propre relativement \`a~$S$ en dimension
$\leqslant -1$ (\resp $\leqslant 0$, \resp $\leqslant 1$) est local
sur~$Y$ pour la topologie \'etale.

Soit $S'$ un $S$-sch\'ema, $F'$ l'image inverse de~$F$ sur~$X'$
(\cf \eqref{eq:XIII.1.0.1}). Si $(F,f)$ est cohomologiquement propre
relativement \`a~$S$ en dimension $\leqslant 1$, il en est de
m\^eme de~$(F', f')$. Mais, si $(F,f)$ est cohomologiquement propre
relativement \`a~$S$ en dimension $\leqslant -1$ (\resp $\leqslant
0$), il n'en est pas n\'ecessairement de m\^eme de~$(F', f')$.

Soit par exemple $S'$ un anneau de valuation discr\`ete, $f' \colon
E_{S'} \to S'$ l'espace affine au-dessus de~$S'$, $x$ un point
ferm\'e de~$E_{S'}$ au-dessus du point g\'en\'erique de~$S'$ et
$F'$ le faisceau d'ensembles sur~$E_{S'}$ dont la restriction \`a
$E_{S'}-\{x\}$ est le faisceau constant \`a un \'el\'ement et
dont la fibre en un point g\'eom\'etrique au-dessus de~$x$ a deux
\'el\'ements. Alors $(F',f')$ n'est pas cohomologiquement propre
relativement \`a~$S'$ en dimension $\leqslant -1$. Soient $S =
S'[Z]$, $f \colon E_S \to S$ l'espace affine sur~$S$ et $T$ une partie
ferm\'ee de~$X = E_S$ qui ne rencontre par le ferm\'e $Z=0$ et
telle que $f(T)$ contienne le point g\'en\'erique de~$S$. Soient
$G$ l'image inverse de~$F'$ sur~$X$ et $F$ le faisceau sur~$X$ obtenu
en prolongeant $G|X-T$ par le vide. Alors $(F,f)$ est
cohomologiquement propre relativement \`a~$S$ en dimension
$\leqslant -1$, mais il n'en est plus de m\^eme apr\`es le
changement de base $S' \to S$ d\'efini par $Z=0$.


d)
Soit $F$ un champ sur~$X$ tel que $(F,f)$ soit cohomologiquement
propre relativement \`a $Y$ en dimension $\leqslant -1$ (\resp $\leqslant 0$, \resp $\leqslant 1$). Alors, pour tout point
g\'eom\'etrique $\overline{y}$ de~$Y$, le foncteur canonique
$$
(f_* F)_{\overline{y}} \to F(X_{\overline{y}})
$$
est fid\`ele (\resp pleinement fid\`ele, \resp une \'equivalence
de cat\'egories).
\end{remarques}

\begin{proposition}
\label{XIII.1.6}
Soient $f \colon X \to Y$ et $g \colon Y \to Z$ deux $S$-morphismes,
$\Phi$ un champ sur~$X$.
\begin{enumerate}
\item[1)] Supposons que $(\Phi,f)$ et $(f_* \Phi, g)$ soient
cohomologiquement propres relativement \`a~$S$ en dimension
$\leqslant -1$ (\resp $\leqslant 0$, \resp $\leqslant 1$). Alors il en
est de m\^eme de~$(\Phi, gf)$.
\item[2)] Supposons que $(\Phi, gf)$ soit cohomologiquement propre
relativement \`a~$S$ en dimension $\leqslant -1$ (\resp que $(\Phi,
gf)$ soit cohomologiquement propre relativement \`a~$S$ en dimension
$\leqslant 0$ et $(\Phi,f)$ cohomologiquement propre relativement
\`a~$S$ en dimension $\leqslant -1$, \resp que $(\Phi, gf)$ soit
cohomologiquement propre relativement \`a~$S$ en dimension
$\leqslant 1$ et $(\Phi, f)$ cohomologiquement propre relativement
\`a~$S$ en dimension $\leqslant 0$). Alors $(f_* \Phi, g)$ est
cohomologiquement propre relativement \`a~$S$ en dimension
$\leqslant -1$ (\resp en dimension $\leqslant 0$, \resp en dimension
$\leqslant 1$).
\end{enumerate}
\end{proposition}

Pour tout $S$-sch\'ema $S'$, on consid\`ere le diagramme suivant,
dont tous les carr\'es sont cart\'esiens:
\begin{equation*}
\label{eq:XIII.1.6.1}
\tag{\thesubsection.1}
\begin{array}{c}
\xymatrix{{X'}\ar[r]^{f'}\ar[d]_h&{Y'}\ar[r]^{g'}\ar[d]_k&{Z'}\ar[r]\ar[d]_{m}&{S'}\ar[d]\\{X}\ar[r]^f&{Y}\ar[r]^g&{Z}\ar[r]&{\,S.}}
\end{array}
\end{equation*}
Le
morphisme canonique
$$
m^* (g_* f_* \Phi) \to g'_* f'_* (h^* \Phi)
$$
s'identifie
au compos\'e des morphismes canoniques
$$
m^*(g_* f_* \Phi) \lto{i} g'_*(k^*f_* \Phi) \lto{j} g'_* f'_*(h^* \Phi).
$$
\begin{enumerate}
\item[1)] L'hypoth\`ese entra\^ine que $i$ et $j$ sont
fid\`eles (\resp pleinement fid\`eles, \resp des
\'equi\-va\-lences); il en est donc de m\^eme de~$ji$.
\item[2)] L'hypoth\`ese entra\^ine que $ji$ est fid\`ele
(\resp que $ji$ est
pleinement
fid\`ele et $j$ fid\`ele, \resp que $ji$ est une \'equivalence
et $j$ pleinement fid\`ele); il en r\'esulte que $i$ est
fid\`ele (\resp pleinement fid\`ele, \resp une \'equivalence).
\end{enumerate}

\begin{corollaire}
\label{XIII.1.7}
Soient $f \colon X \to Y$ et $g \colon Y \to Z$ deux $S$-morphismes,
et soit $F$ un faisceau en groupes sur~$X$. Supposons que $(F,gf)$
soit cohomologiquement propre relativement \`a~$S$ en dimension
$\leqslant -1$ (\resp que $(F,gf)$ soit cohomologiquement propre
relativement \`a~$S$ en dimension $\leqslant 0$ et que $(F,f)$ soit
cohomologiquement propre relativement \`a~$S$ en dimension
$\leqslant -1$). Alors $(f_* F, g)$ est cohomologiquement propre
relativement \`a~$S$ en dimension $\leqslant -1$ (\resp en dimension
$\leqslant 0$).
\end{corollaire}

Reprenons les notations de \eqref{eq:XIII.1.6.1} et, pour tout
sch\'ema $Y_1$ \'etale au-dessus de~$Y$, notons $f_1$, $F_1$ les
images inverses respectives de~$f$, $F$ par le morphisme $Y_1 \to
Y$. Soient $\Phi$ le champ des torseurs sous $F$ et $\Psi$ le champ
des torseurs sous $f_* F$. On a un foncteur canonique
$$
\varphi \colon \Psi \to f_* \Phi \text{,}
$$
obtenu en associant \`a tout sch\'ema $Y_1$ \'etale sur~$Y$ et
\`a tout torseur $P$ sur~$Y_1$ de groupe $f_{1 *}F_1$ le torseur
$\tilde{P}$ sur~$X_1$ d\'eduit de~$f_1^* P$ par l'extension du
groupe structural $f_1^* f_{1 *}F_1 \to F_1$. Le foncteur
$\varphi$ est pleinement fid\`ele. En effet, si $P$ et $Q$ sont deux
torseurs sur~$Y_1$ de groupe $f_{1 *} F_1$, on a un morphisme
canonique
$$
\SheafIsom_{f_{1 *}F_1}(P,Q) \to
f_{1*}(\SheafIsom_{F_1}(\tilde{P},\tilde{Q}))
$$
qui
est un isomorphisme car il en est ainsi localement. On en d\'eduit
que le morphisme canonique
$$
\Isom_{f_{1 *}F_1}(P,Q) \to \Isom_{F_1}(\tilde{P},\tilde{Q})
$$
est un isomorphisme, donc que $\varphi$ est pleinement fid\`ele.

On a un diagramme commutatif
$$
\xymatrix{{g'_* k^* \Psi}\ar[r]^i\ar[d]_{g'_*
k^*(\varphi)}&{m^*(g_* \Psi)}\ar[d]^{m^*
g_*(\varphi)}\\{g'_* k^*(f_* \Phi)}\ar[r]^j&{m^*(g_* f_* \Phi)\text{,}}}
$$
o\`u $i$ et $j$ sont les morphismes de changement de base. Il
r\'esulte de~\ref{XIII.1.6} 2) que $j$ est fid\`ele
(\resp pleinement fid\`ele). Comme $g'_* k^* (\varphi)$ et
$m^* g_* (\varphi)$ sont pleinement fid\`eles, on d\'eduit
du diagramme ci-dessus que $i$ est fid\`ele (\resp pleinement
fid\`ele).

\begin{corollaire}
\label{XIII.1.8}
Soient $f \colon X \to Y$ un $S$-morphisme coh\'erent, $g \colon Y
\to Z$ un $S$-morphisme propre, $\Phi$ un champ ind-fini sur~$X$ \textup{\cite[VII 2.2.1]{XIII.2}}. Supposons que $(\Phi,f)$ soit cohomologiquement propre
relativement \`a~$S$ en dimension $\leqslant -1$ (\resp $\leqslant
0$, \resp $\leqslant 1$), alors il en est de m\^eme de~$(\Phi, gf)$.
\end{corollaire}

Comme $f$ est coh\'erent, $f_* \Phi$ est un champ ind-fini (SGA~4 IX
1.6~(ii)). Le corollaire r\'esulte donc de \ref{XIII.1.6}~1)
et~\ref{XIII.1.4}.

\begin{corollaire}
\label{XIII.1.9}
Soient $f \colon X \to Y$ un $S$-morphisme entier, $g \colon Y \to Z$
un $S$-morphisme. Si $F$ est un faisceau d'ensembles sur~$X$, pour que
$(f_* F, g)$ soit cohomologiquement propre relativement \`a~$S$
en dimension $\leqslant -1$ (\resp $\leqslant 0$), il faut et il
suffit qu'il en soit ainsi de~$(F, gf)$. Si $F$ est un faisceau en
groupes sur~$X$, pour que $(f_* F, g)$ soit cohomologiquement
propre relativement \`a~$S$ en dimension $\leqslant -1$
(\resp $\leqslant 0$, \resp $\leqslant 1$), il faut et il suffit qu'il
en soit ainsi de~$(F, gf)$.
\end{corollaire}

L'assertion
relative au cas d'un faisceau d'ensembles r\'esulte
de~\ref{XIII.1.6} et du fait que $(F,f)$ est cohomologiquement propre
relativement \`a~$S$ en dimension $\leqslant 0$. Soient $F$ un
faisceau en groupes sur~$X$ et $\Phi$ le champ des torseurs sous $F$.
D'apr\`es SGA~4 VIII~5.8, tout torseur sous $F$ est localement
trivial sur~$Y$. Il en r\'esulte que le champ $f_* \Phi$ est
\'equivalent au champ des torseurs sous $f_* F$, l'\'equivalence
\'etant obtenue en associant \`a tout sch\'ema $Y_1$ \'etale
sur~$Y$ et \`a tout torseur $P$ sur~$X_1 = X \times_Y Y_1$ de groupe
$F|X_1$ le torseur $f_* P$ de groupe $f_* F | Y_1$. Comme $(F,f)$ est
cohomologiquement propre relativement \`a~$S$ en dimension
$\leqslant 1$, le corollaire r\'esulte donc de~\ref{XIII.1.6}.

\enlargethispage{.5cm}
\skpt
\begin{definitions}
\label{XIII.1.10}

\subsubsection{}
\label{XIII.1.10.1}
Soit $E$ une cat\'egorie et consid\'erons un diagramme
$$
\xymatrix@C=.6cm{{\Phi}\ar[r]^p&{\Phi_1}\ar@<2pt>[r]^{p_1}\ar@<-2pt>[r]_{p_2}&{\Phi_2}} \text{,}
$$
o\`u $\Phi$, $\Phi_1$, $\Phi_2$ sont des cat\'egories fibr\'ees
au-dessus de~$E$ et les fl\`eches des morphismes de cat\'egories
fibr\'ees, et soit
$$
a \colon p_1 p \isomto p_2 p
$$
un isomorphisme de foncteurs.

On dit que le diagramme ci-dessus est exact si la condition suivante
est satisfaite

a) Pour tout couple d'objets $x$, $y$ de~$\Phi$ et tout morphisme $f
\colon p(x) \to p(y)$ tel que l'on ait $p_1(f) = p_2(f)$ ($p_1p$ et
$p_2p$ \'etant identifi\'es gr\^ace \`a $a$), il existe un
unique morphisme $g \colon x \to y$ tel que l'on ait $p(g)=f$.

\subsubsection{}
\label{XIII.1.10.2}
Consid\'erons le diagramme
$$
\xymatrix@C=.5cm{{\Phi}\ar[r]^p&{\Phi_1}\ar@<2pt>[rr]^{p_1,p_2}\ar@<-2pt>[rr]&&{\Phi_2}\ar@<4pt>[rrr]^{p_{12},p_{23},p_{31}}\ar[rrr]\ar@<-4pt>[rrr]&&&{\Phi_3}} \text{,}
$$
o\`u
$\Phi$, $\Phi_i$, $1 \leqslant i \leqslant 3$, sont des cat\'egories
fibr\'ees sur~$E$ et les fl\`eches des morphismes de
cat\'egories fibr\'ees. Supposons donn\'es des isomorphismes de
foncteurs
$$
a \colon p_1 p \isomto p_2 p
$$
$$
a_1 \colon p_{31} p_2 \isomto p_{12} p_1 \text{,}\quad a_2 \colon p_{12}
p_2 \isomto p_{23} p_1 \text{,}\quad a_3 \colon p_{23} p_2 \isomto
p_{31}p_1
$$
tels que le diagramme suivant soit commutatif:
$$
\xymatrix@C=1.2cm{{p_{23} p_1 p}\ar[r]^{\id.a}&{p_{23} p_2 p}\ar[r]^{a_3.\id}&{p_{31} p_1 p}\ar[d]^{\id.a}\\{p_{12} p_2 p}\ar[u]^{a_2.\id}&{p_{12} p_1 p}\ar[l]^{\id.a}&{p_{31} p_2 p}\ar[l]^{a_1.\id}} \text{.}
$$
On identifie $p_1 p$ et $p_2 p$, $p_{31} p_2$ et $p_{12} p_1$, etc.

On dit que le diagramme ci-dessus est exact si les conditions
suivantes sont satisfaites:
\begin{enumerate}
\item[a)] Analogue \`a la condition a) de~\ref{XIII.1.10.1}.
\item[b)] Pour tout objet $x_1$ de~$\Phi_1$ et pour tout isomorphisme
$u \colon p_1(x_1) \isomto p_2(x_1)$ tel que l'on ait
\begin{equation*}
\label{eq:XIII.1.10.2.1}
\tag{\thesubsubsection.1} { p_{23}(u) p_{31}(u) = p_{12}(u)^{-1}
\text{,} }
\end{equation*}
il existe un objet $x$ de~$\Phi$ tel que l'on ait un isomorphisme $i
\colon p(x) \isomto x_1$ rendant commutatif le diagramme
\begin{equation*}
\label{eq:XIII.1.10.2.2}
\tag{\thesubsubsection.2}
\begin{array}{c}
{ \xymatrix{{p_1p(x)}\ar@{=}[r]\ar[d]_{p_1(i)}&{p_2p(x)}\ar[d]^{p_2(i)}\\{p_1(x_1)}\ar[r]^{\sim}_u&{p_2(x_1)}} }
\end{array}
\end{equation*}
\end{enumerate}

\subsubsection{}
\label{XIII.1.10.3}
On d\'efinit de fa\c con \'evidente la notion de morphisme de
diagrammes exacts de cat\'egories fibr\'ees au-dessus d'une
cat\'egorie $E$.
\end{definitions}

\subsubsection{}
\label{XIII.1.10.4}
Nous utiliserons
en particulier la notion de diagramme exact dans le
cas o\`u $E$ est un site et $\Phi$, $\Phi_i$, $1 \leqslant i
\leqslant 3$, des champs sur~$E$.

Soient $f \colon E \to E'$ un morphisme de sites et
\begin{equation*}
\label{eq:XIII.1.10.4.1}
\tag{\thesubsubsection.1} { \xymatrix@C=.5cm{{\Phi}\ar[r]^p&{\Phi_1}\ar@<2pt>[rr]^{p_1,p_2}\ar@<-2pt>[rr]&&{\Phi_2}\ar@<4pt>[rrr]^{p_{12},p_{23},p_{31}}\ar[rrr]\ar@<-4pt>[rrr]&&&{\Phi_3}} }
\end{equation*}
un diagramme exact de champs sur~$E$. On obtient par image directe un
diagramme
$$
\xymatrix@C=.5cm{{f_* \Phi}\ar[r]&{f_* \Phi_1}\ar@<2pt>[r]\ar@<-2pt>[r]&{f_* \Phi_2}\ar@<4pt>[r]\ar[r]\ar@<-4pt>[r]&{f_* \Phi_3}}
$$
qui est \'evidemment exact.

Si $h \colon E'' \to E$ est un morphisme de sites, on a de m\^eme un
diagramme exact
$$
\xymatrix@C=.5cm{{h^* \Phi}\ar[r]^{p''}&{h^* \Phi_1}\ar@<2pt>[rr]^{p''_1,p''_2}\ar@<-2pt>[rr]&&{h^* \Phi_2}\ar@<4pt>[rrr]^{p''_{12},p''_{23},p''_{31}}\ar[rrr]\ar@<-4pt>[rrr]&&&{h^* \Phi_3}}
$$

V\'erifions d'abord la condition~a) de~\ref{XIII.1.10.2}. Soient
$F''$ un objet de~$E''$, $x''$ et $y''$ deux objets de~$(h^*
\Phi)_{F''}$, $x''_1$ et $y''_1$ leurs images respectives dans
$h^* \Phi_1$ et $x''_2$, $y''_2$ leurs images dans $h^*
\Phi_2$. Soit $u''_1 \colon x''_1 \to y''_1$ un morphisme tel que l'on
ait $p''_1(u''_1) = p''_2(u''_1)$ et prouvons que $u''_1$ provient
d'un unique morphisme $u'' \colon x'' \to y''$. La question \'etant
locale sur~$F''$, on peut supposer qu'on a un objet $F_1$ de~$E$, un
morphisme de~$F''$ dans l'image inverse $F''_1$ de~$F_1$ par $h$, des
objets $x$, $y$ de~$\Phi_{F_1}$ dont les images inverses sur~$F''$
sont $x''$ et $y''$. Soient $x_1$, $y_1$ (\resp $x_2$, $y_2$) les
images de~$x$, $y$ dans $\Phi_1$ (\resp~$\Phi_2$). On peut supposer
que $u''_1$ provient d'un morphisme $u_1 \colon x_1 \to y_1$, tel que
l'on ait $p_1(u_1) = p_2(u_1)$. Vu l'exactitude de
\eqref{eq:XIII.1.10.4.1}, on obtient un unique morphisme $u \colon x
\to y$ dont l'image inverse par $h$ est le morphisme $u''$
cherch\'e.

La
condition~b) de~\ref{XIII.1.10.2} se v\'erifie de fa\c con
analogue. Soient $x''_1$ un objet de~$(h^* \Phi_1)_{F''}$, $u''
\colon p''_1(x''_1) \to p''_2(x''_1)$ un morphisme satisfaisant \`a
la relation
$$
p''_{23}(u'')p''_{31}(u'')=p''_{12}(u'')^{-1} \text{,}
$$
et prouvons qu'il existe un objet $x''$ de~$(h^* \Phi)_{F''}$ et
un isomorphisme $i'' \colon p''(x'') \simeq x''_1$ rendant commutatif
un diagramme analogue \`a \eqref{eq:XIII.1.10.2.2}. Comme la
question est locale sur~$F''$, on peut supposer qu'on a un objet
$F_1$, un morphisme $F'' \to F''_1$ comme ci-dessus, et un objet $x_1$
de~$(\Phi_1)_{F_1}$ dont l'image inverse dans $(h^* \Phi_1)_{F''}$
est $x''_1$. De m\^eme on peut supposer que $u''$ provient d'un
morphisme $u \colon p_1(x_1) \to p_2(x_1)$ satisfaisant \`a
\eqref{eq:XIII.1.10.2.2}. L'existence d'un objet $x$ de~$\Phi_{F_1}$,
dont l'image inverse par $h$ soit un \'el\'ement $x''$
r\'epondant \`a la question, r\'esulte alors de l'exactitude de
\eqref{eq:XIII.1.10.4.1}.

\begin{exemples}
\label{XIII.1.11}
\begin{enumerate}
\item[1)] Soit $f \colon X_1 \to X$ un \emph{morphisme de descente}
pour la cat\'egorie des faisceaux \'etales sur des sch\'emas
variables (par exemple un morphisme universellement submersif (SGA~4
VIII~9.3)). Soient $X_2 = X_1 \times_X X_1$, $g \colon X_2 \to X$ la
projection canonique et $F$ un faisceau d'ensembles sur~$X$. Il
r\'esulte alors de \loccit que l'on a une suite exacte de
faisceaux d'ensembles
\begin{equation*}
\label{eq:XIII.1.11.1}
\tag{\thesubsection.1} { \xymatrix@C=.5cm{{F}\ar[r]&{f_* f^* F}\ar@<2pt>[r]\ar@<-2pt>[r]&{g_* g^* F \text{.}}}
}
\end{equation*}
Si $\Phi$ est le champ en cat\'egories discr\`etes associ\'e
\`a $F$ et $\Phi_3$ le champ final sur~$X$, \ie le champ dont
toutes les fibres sont r\'eduites \`a un seul \'el\'ement
ayant pour seul morphisme le morphisme identique, dire que la suite
\eqref{eq:XIII.1.11.1} est exacte revient \`a dire qu'il en est
ainsi du diagramme de champs
$$
\xymatrix@C=.5cm{{\Phi}\ar[r]&{f_* f^* \Phi}\ar@<2pt>[r]\ar@<-2pt>[r]&{g_* g^* \Phi}\ar[r]\ar@<4pt>[r]\ar@<-4pt>[r]&{\Phi_3}} \text{.}
$$

\item[2)]
Soit
$f \colon X_1 \to X$ un \emph{morphisme de descente effective}
pour la cat\'egorie des faisceaux \'etales sur des sch\'emas
variables (par exemple un morphisme propre surjectif, ou un morphisme
entier surjectif, ou un morphisme fid\`element plat localement de
pr\'esentation finie (SGA~4 VIII~9.4)). Soient $X_2 = X_1 \times_X
X_1$, $g \colon X_2 \to X$ la projection canonique, $X_3 = X_1
\times_X X_1 \times_X X_1$, $h \colon X_3 \to X$ le morphisme
canonique. Soient $\Phi$ un champ sur~$X$, $\Phi_1 = f_* f^* \Phi$,
$\Phi_2 = g_* g^* \Phi$, $\Phi_3 = h_* h^* \Phi$. On a alors un
diagramme exact
$$
\xymatrix@C=.5cm{{\Phi}\ar[r]&{\Phi_1}\ar@<2pt>[r]\ar@<-2pt>[r]&{\Phi_2}\ar[r]\ar@<4pt>[r]\ar@<-4pt>[r]&{\Phi_3 \text{,}}}
$$
o\`u les fl\`eches sont les morphismes canoniques associ\'es aux
projections.

Consid\'erons en effet $\Phi$ comme un champ sur la cat\'egorie
$\Sch_X$ des sch\'emas au-dessus de~$X$, munie de la topologie
\'etale. Alors, d'apr\`es \cite[VII 2.2.8]{XIII.2}, $\Phi$ est aussi un
champ pour la topologie la plus fine sur~$\Sch_X$ pour laquelle les
morphismes couvrants sont les morphismes de descente effective pour la
cat\'egorie des faisceaux \'etales. L'exactitude du diagramme
ci-dessus en r\'esulte aussit\^ot.
\end{enumerate}
\end{exemples}

\begin{proposition}
\label{XIII.1.12}
Soient $S$ un sch\'ema, $f\colon X \to Y$ un $S$-morphisme.
\begin{enumerate}
\item[1)] Soit
$$
\xymatrix@C=.5cm{{\Phi}\ar[r]&{\Phi_1}\ar@<2pt>[r]\ar@<-2pt>[r]&{\Phi_2}}
$$
un diagramme exact de champs sur~$X$. Supposons que $(\Phi_1, f)$ soit
cohomologique\-ment~propre relativement \`a~$S$ en dimension
$\leqslant 0$ et que $(\Phi_2, f)$ soit cohomologi\-quement~propre~relativement \`a~$S$ en dimension $\leqslant -1$. Alors $(\Phi,f)$
est cohomologiquement propre~relativement \`a~$S$ en dimension
$\leqslant 0$.

\item[2)] Soit
$$
\xymatrix@C=.5cm{{\Phi}\ar[r]&{\Phi_1}\ar@<2pt>[r]\ar@<-2pt>[r]&{\Phi_2}\ar@<4pt>[r]\ar[r]\ar@<-4pt>[r]&{\Phi_3}}
$$
un diagramme exact de champs sur~$X$. Supposons que $(\Phi,f)$ soit
cohomologi\-quement~propre relativement \`a~$S$ en dimension
$\leqslant 1$, que $(\Phi_2, f)$ soit cohomologiquement~propre
relativement \`a~$S$ en dimension $\leqslant 0$ et que
$(\Phi_3, f)$ soit cohomo\-logiquement~propre
relativement
\`a~$S$ en dimension $\leqslant -1$. Alors $(\Phi,f)$ est
cohomologique\-ment~propre~relativement \`a~$S$ en dimension
$\leqslant 1$.
\end{enumerate}
\end{proposition}

Pour tout $S$-sch\'ema $S'$, on consid\`ere le diagramme
commutatif suivant dont tous les carr\'es sont cart\'esiens:
$$
\xymatrix{{X'}\ar[r]^{f'}\ar[d]_h&{Y'}\ar[r]\ar[d]_g&{S'}\ar[d]\\{X}\ar[r]^f&{Y}\ar[r]&{S}}
$$

D\'emontrons 2), la d\'emonstration de 1) \'etant
analogue. Comme les foncteurs image directe et image inverse
transforment diagramme exact de champs en diagramme exact
\eqref{XIII.1.10.4}, on a le morphisme de diagrammes exacts de champs
suivant:
$$
\xymatrix{{g^* f_* \Phi}\ar[r]^\pi\ar[d]_{\varphi}&{g^*
f_* \Phi_1}\ar@<2pt>[r]^{\pi_1}\ar@<-2pt>[r]_{\pi_2}\ar[d]_{\varphi_1}&{g^* f_* \Phi_2}\ar@<4pt>[r]\ar[r]\ar@<-4pt>[r]\ar[d]_{\varphi_2}&{g^* f_* \Phi_3}\ar[d]_{\varphi_3}\\{f'_* h^* \Phi}\ar[r]&{f'_* h^*
\Phi_1}\ar@<2pt>[r]\ar@<-2pt>[r]&{f'_* h^* \Phi_2}\ar@<4pt>[r]\ar[r]\ar@<-4pt>[r]&{f'_* h^* \Phi_3}}
\text{.}
$$
Par hypoth\`ese $\varphi_1$ est une \'equivalence de
cat\'egories, $\varphi_2$ est pleinement fid\`ele et $\varphi_3$
fid\`ele. Il r\'esulte donc du diagramme pr\'ec\'edent que
$\varphi$ est une \'equivalence.

\begin{proposition}
\label{XIII.1.13}
Soit $f \colon X \to Y$ un $S$-morphisme.
\begin{enumerate}
\item[1)] Soit
$$
\xymatrix@C=.5cm{{F}\ar[r]&{G}\ar@<2pt>[r]\ar@<-2pt>[r]&{H}}
$$
un diagramme exact de faisceaux d'ensembles sur~$X$. Supposons que
$(G,f)$ soit cohomologiquement propre relativement \`a~$S$ en
dimension $\leqslant 0$ et que $(H,f)$ soit cohomologiquement propre
relativement \`a~$S$ en dimension $\leqslant -1$. Alors $(F,f)$ est
cohomologiquement propre relativement \`a~$S$ en dimension
$\leqslant 0$.
\item[2)]
Soit
$F \to G$ un monomorphisme de faisceaux en groupes sur~$X$. Si
$Y_1$ est un sch\'ema \'etale sur~$Y$, on pose $X_1 = Y_1 \times_Y
X$, et on note $f_1$ (\resp $F_1$, \resp $G_1$) l'image inverse de~$f$
(\resp $F$, \resp $G$) sur~$Y_1$ (\cf \ref{eq:XIII.1.0.2}). Supposons
que $(G,f)$ soit cohomologiquement propre relativement \`a~$S$ en
dimension $\leqslant 0$ (\resp en dimension $\leqslant 1$) et que,
pour tout sch\'ema $Y_1$ \'etale sur~$Y$ et pour tout torseur $Q$
sous $G_1$, $(Q/F_1, f_1)$ soit cohomologiquement propre relativement
\`a~$S$ en dimension $\leqslant -1$ (\resp en dimension $\leqslant
0$). Alors $(F,f)$ est cohomologiquement propre relativement \`a~$S$
en dimension $\leqslant 0$ (\resp en dimension $\leqslant 1$).
\item[3)] Soit $F \to G$ un monomorphisme de faisceaux en groupes sur~$X$. Supposons que $(F,f)$ soit cohomologiquement propre relativement
\`a~$S$ en dimension $\leqslant 1$ et que $(G,f)$ soit
cohomologiquement propre relativement \`a~$S$ en dimension
$\leqslant 0$. Alors, pour tout torseur $Q$ sous $G$, $(Q/F, f)$ est
cohomologiquement propre relativement \`a~$S$ en dimension
$\leqslant 0$.
\end{enumerate}
\end{proposition}

\subsubsection*{D\'emonstration}
1) Soit $\Phi$ (\resp $\Phi_1$, \resp $\Phi_2$) le champ en
cat\'egories discr\`etes associ\'e \`a $F$ (\resp $G$,
\resp $H$) et soit $\Phi_3$ le champ final sur~$X$. On a alors un
diagramme exact
$$
\xymatrix@C=.5cm{{\Phi}\ar[r]&{\Phi_1}\ar@<2pt>[r]\ar@<-2pt>[r]&{\Phi_2}\ar@<4pt>[r]\ar[r]\ar@<-4pt>[r]&{\Phi_3}}
\text{.}
$$
Par hypoth\`ese $(\Phi_1, f)$ est cohomologiquement propre
relativement \`a~$S$ en dimension $\leqslant 1$ et $(\Phi_2, f)$ est
cohomologiquement propre relativement \`a~$S$ en dimension
$\leqslant 0$ \eqref{XIII.1.2}. Comme $(\Phi_3, f)$ est \'evidemment
cohomologiquement propre relativement \`a~$S$ en dimension
$\leqslant -1$, il r\'esulte de~\ref{XIII.1.12} que $(\Phi, f)$ est
cohomologiquement propre relativement \`a~$S$ en dimension
$\leqslant 1$, \ie que $(F,f)$ est cohomologiquement propre
relativement \`a~$S$ en dimension $\leqslant 0$.

\smallskip
2) Montrons d'abord que, si $(G,f)$ est cohomologiquement
propre relativement \`a~$S$ en dimension $\leqslant 0$ et si les
$(Q/F_1, f_1)$ sont cohomologiquement
propres relativement \`a~$S$ en dimension $\leqslant -1$, alors
$(F,f)$ est cohomologiquement propre relativement \`a~$S$ en
dimension $\leqslant 0$. D'apr\`es~\ref{XIII.1.3.1} il suffit de
prouver que, pour tout sch\'ema $Y_1$ \'etale au-dessus de~$Y$ et
pour tout torseur $P$ sur~$X_1$ de groupe $F_1$, $(\leftexp{P\mkern-4mu}{F}_1,
f_1)$ est cohomologiquement propre relativement \`a~$S$ en dimension
$\leqslant 0$ quand on consid\`ere $\leftexp{P\mkern-4mu}{F}_1$ comme un
faisceau d'ensembles, et que le morphisme canonique
$$
d \colon g^*(\R^1 f_* F) \to \R^1 f'_* F'
$$
est injectif. La premi\`ere assertion r\'esulte aussit\^ot de 1)
car, si $Q$ d\'esigne le torseur d\'eduit de~$P$ par l'extension
$F_1 \to G_1$ du groupe structural, on a un isomorphisme
$\leftexp{Q}{G}_1 / \leftexp{P\mkern-4mu}{F}_1 \isomto Q/F_1$.

Montrons que $d$ est injectif. Il suffit de prouver que, si $Y_1$ est
un sch\'ema \'etale au-dessus de~$Y$, si $P$ et $\tilde{P}$ sont
deux torseurs sous $F_1$ dont les images inverses $P'$ et~$\tilde{P}'$
sur~$X'_1$ sont isomorphes, alors, quitte \`a faire une extension
\'etale surjective de~$Y_1$, $P$ et~$\tilde{P}$ deviennent
isomorphes. Choisissons un isomorphisme $p' \colon P' \isomto
\tilde{P}'$. Soient $Q$ (\resp~$\tilde{Q}$) le torseur d\'eduit de
$P$ (\resp $\tilde{P}$) par l'extension du groupe structural $F_1 \to
G_1$. Les images inverses $Q'$ (\resp $\tilde{Q}'$) de~$Q$
(\resp $\tilde{Q}$) sur~$X'_1$ se d\'eduisent de~$P'$
(\resp $\tilde{P}'$) par
l'extension
du groupe structural $F'_1 \to G'_1$; soit $q' \colon Q' \isomto
\tilde{Q}'$ l'isomorphisme que l'on obtient de m\^eme \`a partir
de~$p'$. Comme $(G,f)$ est cohomologiquement propre relativement \`a~$S$ en dimension $\leqslant 0$, on peut supposer, quitte \`a faire
une extension \'etale surjective de~$Y_1$, que $q'$ est l'image d'un
isomorphisme $q \colon Q \isomto \tilde{Q}$. Au torseur~$P$
(\resp~$\tilde{P}$) est associ\'ee une section $x$ de~$Q/F_1$
(\resp une section $\tilde{x}$ de~$\tilde{Q}/F_1$), et, pour que $P$
et $\tilde{P}$ soient isomorphes, il faut et il suffit que l'on ait un
isomorphisme $Q \isomto \tilde{Q}$ tel que l'isomorphisme
$$
e \colon \H^0(X_1, Q/F_1) \to \H^0(X_1,\tilde{Q}/F_1)
$$
qu'on en d\'eduit, transforme $x$ en $\tilde{x}$. On prend
l'isomorphisme $q$. Les sections $e(x)$ et $\tilde{x}$ de
$\H^0(X_1,\tilde{Q}/F_1)$ ont m\^eme image dans
$\H^0(X'_1,\tilde{Q}'/F'_1)$. Comme $(\tilde{Q}/F_1,f_1)$ est
cohomologiquement propre relativement \`a~$S$ en dimension
$\leqslant -1$,
quitte \`a faire une extension \'etale surjective de~$Y_1$, on a
bien $e(x) = \tilde{x}$, ce qui d\'emontre l'injectivit\'e de~$d$.

Pour achever la d\'emonstration, il reste \`a prouver que, si
$(G,f)$ est cohomologiquement propre relativement \`a~$S$ en
dimension $\leqslant 1$, et si, pour tout sch\'ema $Y_1$ \'etale
au-dessus de~$Y$ et tout torseur $Q$ sur~$X_1$ de groupe $F_1$,
$(Q/F_1, f_1)$ est cohomologiquement propre relativement \`a~$S$ en
dimension $\leqslant 0$, alors le morphisme $d$ est surjectif. Soient
$P'$ un torseur sur~$X'_1$ de groupe $F'_1$, $Q'$ le torseur sous
$G'_1$ obtenu \`a partir de~$P'$ par extension du groupe
structural. La donn\'ee de~$P'$ est \'equivalente \`a celle de
$Q'$ et d'une section $x'$ de~$\H^0(X'_1,Q'/F'_1)$. Il r\'esulte
alors de la surjectivit\'e du morphisme
$$
g^*(\R^1 f_* G) \to \R^1 f'_* G'
$$
que, quitte \`a faire une extension \'etale surjective de~$Y_1$,
il existe un torseur $Q$ sous~$G_1$, dont l'image inverse sur~$X'_1$
est isomorphe \`a $Q'$. Utilisant le fait que $(Q/F_1, f_1)$ est
cohomologiquement propre relativement \`a~$S$ en dimension
$\leqslant 0$, on peut de m\^eme supposer qu'il existe un
\'el\'ement $x$ de~$\H^0(X_1, Q/F_1)$ dont l'image dans
$\H^0(X'_1, Q'/F'_1)$ est $x'$. La donn\'ee de~$Q$ et de~$x$
d\'etermine un torseur $P$ sous $F_1$, dont l'image inverse sur~$X_1$ est isomorphe \`a $P'$, ce qui d\'emontre la
surjectivit\'e de~$d$.

\smallskip
3) Montrons que $(Q/F,f)$ est cohomologiquement propre
relativement \`a~$S$ en dimension $\leqslant -1$, \ie que, pour
tout $S$-sch\'ema $S'$, pour tout sch\'ema $Y_1$ \'etale
au-dessus de~$Y$, si $x$, $\tilde{x}$ sont deux \'el\'ements de
$\H^0(X_1,Q_1/F_1)$ dont les images $x'$, $\tilde{x}'$ dans
$\H^0(X'_1, Q'_1/F'_1)$ sont \'egales, alors, apr\`es extension
surjective de~$Y_1$, on a $x = \tilde{x}$. \`A~$x$ (\resp~$\tilde{x}$)
est associ\'e un torseur $P$ (\resp~$\tilde{P}$) sous $F_1$, tel que
$Q_1$ se d\'eduise de~$P$ (\resp~$\tilde{P}$) par l'extension $F_1
\to G_1$ du groupe structural. De la relation $x' = \tilde{x}'$
r\'esulte que l'on a un isomorphisme $u' \colon P' \isomto
\tilde{P}'$ tel que l'isomorphisme induit sur~$Q'_1$ par~$u'$ soit
l'identit\'e. Comme $(F,f)$ est cohomologiquement propre
relativement \`a~$S$ en dimension $\leqslant 0$, on en d\'eduit
que,
apr\`es extension \'etale surjective de~$Y_1$, on a un
isomorphisme $u \colon P \to \tilde{P}$ relevant $u'$; le fait que
$(G,f)$ soit cohomologiquement propre relativement \`a~$S$ en
dimension $\leqslant -1$ entra\^ine alors que l'on a $x =
\tilde{x}$.

Montrons que $(Q/F,f)$ est cohomologiquement propre relativement \`a~$S$ en dimension $\leqslant 0$. Soient $Y''$ un sch\'ema \'etale
sur~$Y'$ et $x''$ un \'el\'ement de~$\H^0(X'',Q''/F'')$. \`A~$x''$
est associ\'e un torseur $P''$ sur~$X''$ de groupe $F''$. Comme
$(F,f)$ est cohomologiquement propre relativement \`a~$S$ en
dimension $\leqslant 1$, on peut trouver des morphismes \'etales
surjectifs $Y''_1 \to Y''$ et $Y_1 \to Y$, tels que l'on ait un
morphisme $Y''_1 \to Y'_1$, et un torseur $P$ sur~$X_1$ de groupe
$F_1$ dont l'image inverse sur~$X''_1$ soit isomorphe \`a l'image
inverse de~$P''$. Il r\'esulte alors du fait que $(G,f)$ est
cohomologiquement propre relativement \`a~$S$ en dimension
$\leqslant 0$ que l'on peut m\^eme choisir $Y''_1$ et $Y_1$ tels que
le torseur d\'eduit de~$P$ par extension du groupe structural $F_1
\to G_1$ soit isomorphe \`a $Q_1$; il correspond \`a~$P$ un
\'el\'ement $x$ de~$\H^0(X_1,Q_1/F_1)$, dont l'image dans
$\H^0(X''_1,Q''_1/F''_1)$ est isomorphe \`a l'image inverse de
$x''$, ce qui ach\`eve la d\'emonstration.

\begin{proposition}
\label{XIII.1.14}
Soient $f \colon X \to S$ un $S$-sch\'ema, $F$ un faisceau
d'ensembles ou de groupes sur~$X$ (\resp un faisceau de
ind-$\LL$-groupes, o\`u $\LL$ est un ensemble de nombres
premiers). Supposons $F$ localement constant, $(F,f)$
cohomologiquement propre en dimension $\leqslant 0$ (\resp en
dimension $\leqslant 1$) et $f$ localement $0$-acyclique
(\resp localement $1$-asph\'erique pour $\LL$) \textup{(SGA~4 XV~1.11)}.
Alors, pour toute sp\'ecialisation $\overline{s}_1 \to
\overline{s}_2$ de points g\'eom\'etriques de~$S$, le morphisme de
sp\'ecialisation \textup{(SGA~4 VIII~7.1)}
$$
a_0 \colon (f_* F)_{\overline{s}_2} \to (f_*
F)_{\overline{s}_1}
$$
est un isomorphisme, et, si $F$ est un faisceau en groupes, le
morphisme
$$
a_1 \colon (\R^1 f_* F)_{\overline{s}_2} \to (\R^1 f_*
F)_{\overline{s}_1}
$$
est injectif (\resp les morphismes $a_0$ et $a_1$ sont des
isomorphismes).
\end{proposition}

La
d\'emonstration s'obtient en recopiant mot \`a mot celle de SGA~4
XVI~2.3, mais en y rempla\c cant l'expression \og propre\fg par
l'expression \og cohomologiquement propre\fg.

\enlargethispage{.5cm}

\begin{corollaire}
\label{XIII.1.15}
Soient $f \colon X \to S$ un morphisme, $\Phi$ un champ sur~$X$, $\LL$
un ensemble de nombres premiers. Supposons que, pour tout sch\'ema
$X_1$ \'etale sur~$X$ et pour tout couple d'objets $x$, $y$ de
$\Phi_{X_1}$, le faisceau $\SheafHom_{X_1}(x,y)$ soit localement
constant, que le faisceau $\SheafAut_{X_1}(x)$ soit un
ind-$\LL$-groupe localement constant, et que le faisceau des
sous-gerbes maximales $S\Phi$ de~$\Phi$ \textup{\cite[III 2.1.7]{XIII.1}} soit localement
constant. Supposons que $(\Phi,f)$ soit cohomologiquement propre en
dimension $\leqslant 1$ et que $f$ soit localement $1$-asph\'erique
pour $\LL$. Alors, pour toute sp\'ecialisation $\overline{s}_1 \to
\overline{s}_2$ de points g\'eom\'etriques de~$S$, le morphisme de
sp\'ecialisation
$$
a \colon (f_* \Phi)_{\overline{s}_2} \to (f_*
\Phi)_{\overline{s}_1}
$$
est une \'equivalence de cat\'egories.
\end{corollaire}

Soient $\overline{S}_1$ (\resp $\overline{S}_2$) le localis\'e
strict de~$S$ en $\overline{s}_1$ (\resp le localis\'e strict de~$S$
en $\overline{s}_2$), $\overline{X}_2$, $\overline{\Phi}_2$
(\resp $\overline{X}_1$, $\overline{\Phi}_1$) les images inverses de
$X_2$, $\Phi_2$ sur~$\overline{S}_2$ (\resp de~$X_1$, $\Phi_1$ sur~$\overline{S}_1$) et consid\'erons le carr\'e cart\'esien
$$
\xymatrix{{\overline{X}_1}\ar[r]^h\ar[d]_{\overline{f}_1}&{\overline{X}_2}\ar[d]^{\overline{f}_2}\\{\overline{S}_1}\ar[r]^g&{\overline{S}_2}}
\text{.}
$$
On doit montrer que le foncteur
$$
\varphi \colon \overline{\Phi}_2(\overline{X}_2) \to \overline{\Phi}_1
(\overline{X}_1)
$$
est une \'equivalence. Le foncteur $\varphi$ est pleinement
fid\`ele; soient

en
effet $x$, $y$ deux objets de~$(\overline{\Phi}_2)_{\overline{X}_2}$;
le morphisme canonique
$$
\Hom_{\overline{X}_2}(x,y) \to
\Hom_{\overline{X}_1}(\varphi(x),\varphi(y))
$$
s'identifie au morphisme canonique
$$
\H^0(\overline{X}_2, \SheafHom_{\overline{X}_2}(x,y)) \to
\H^0(\overline{X}_1,h^*(\SheafHom_{\overline{X}_2}(x,y)) \text{.}
$$
Ce morphisme est un isomorphisme d'apr\`es~\ref{XIII.1.14}.

Montrons que $\varphi$ est une \'equivalence. Soient $x_1$ un objet
de~$\overline{\Phi}_1(\overline{X}_1)$ et $G_1$ la sous-gerbe maximale
de~$\overline{\Phi}_1$ engendr\'ee par $x_1$. Le morphisme
$$
\H^0(\overline{X}_2, S\overline{\Phi}_2) \to \H^0(\overline{X}_1,
h^*(S \overline{\Phi}_2)) =
\H^0(\overline{X}_1,S\overline{\Phi}_1)
$$
est bijectif, et il existe donc une sous-gerbe maximale $G_2$ de
$\overline{\Phi}_2$ telle que $h^* G_2$ soit isomorphe \`a
$G_1$. Il suffit alors de prouver que le foncteur
$$
G_2 \to h_* h^* G_2
$$
est une \'equivalence. Mais, sous cette forme, la question est
locale pour la topologie \'etale sur~$\overline{X}_2$. On peut donc
supposer que $G_2$ est une gerbe de torseurs sous le groupe des
automorphismes d'un objet de~$G_2$, cas o\`u l'assertion r\'esulte
de~\ref{XIII.1.14}.

\begin{corollaire}
\label{XIII.1.16}
Les hypoth\`eses sont celles de~\ref{XIII.1.14}. Si l'on suppose de
plus que $F$ est un faisceau d'ensembles (\resp de ind-$\LL$-groupes)
et que $f_* F$ (\resp $\R^1 f_* F$) est constructible, alors
$f_* F$ (\resp $\R^1 f_* F$) est localement constant.
\end{corollaire}

Le corollaire r\'esulte de~\ref{XIII.1.14} gr\^ace \`a SGA~4
IX~2.11.

\begin{remarque}
\label{XIII.1.17}
Rappelons que la condition $f$ localement $0$-acyclique est satisfaite
si $f$ est plat \`a fibres s\'eparables, $X$ et $Y$ localement
noeth\'eriens (SGA~4 XV~4.1), et que la condition $f$ localement
$1$-asph\'erique
pour $\LL$ est satisfaite si $f$ est lisse, $\LL$ \'etant l'ensemble
des nombres premiers distincts des caract\'eristiques
r\'esiduelles de~$S$ (SGA~4 XV~2.1).
\end{remarque}

\section[Un cas particulier de propret\'e cohomologique]{Un cas particulier de propret\'e cohomologique: diviseurs
\`a croisements normaux relatifs}
\label{XIII.2}

\setcounter{subsection}{-1}

\subsection{}
\label{XIII.2.0}
Soient $R$ un anneau de valuation discr\`ete de corps des fractions
$K$ et $L$ une $K$-alg\`ebre \'etale; $L$ est alors produit direct
d'un nombre fini de corps $L_i$, o\`u $L_i$ est une extension
\'etale de~$K$. Si $L'_i$ d\'esigne l'extension galoisienne
engendr\'ee par $L_i$ dans une cl\^oture alg\'ebrique de~$L_i$,
on dit que $L$ est \emph{modérément ramifiée}
\index{mod\'er\'ement ramifi\'ee (alg\`ebre)|hyperpage}sur~$R$ si les $L'_i$ sont des extensions mod\'er\'ement
ramifi\'ees au sens de~X~\ref{X.3}, \ie si un groupe d'inertie~$I_i$ de~$L'_i|K$ est d'ordre premier \`a la caract\'eristique
r\'esiduelle $p$ de~$R$.

On sait que $I_i$ est en tous cas extension d'un groupe cyclique
d'ordre premier \`a $p$ par un $p$-groupe. (Cela r\'esulte de \cite[ch.\, IV prop.\, 7 cor.\, 4]{XIII.5} lorsque l'extension r\'esiduelle de~$R$ est
s\'eparable. La d\'emonstration donn\'ee dans
\loccit s'\'etend au cas g\'en\'eral de la fa\c con
suivante. Reprenons les hypoth\`eses et les notations de
\loccit mais sans supposer l'extension r\'esiduelle
s\'eparable. Soit $H_i$ le sous-groupe du groupe d'inertie $G_0$,
ensemble des \'el\'ements $s$ de~$G_0$ tels que l'on ait $s \pi /
\pi \in U^i$ pour toute uniformisante $\pi$ de~$A_L$. On v\'erifie
alors que $G_0/H_1$ est un groupe d'ordre premier \`a $p$ et que,
pour $i \geqslant 1$, les $H^i/H^{i+1}$ sont des $p$-groupes, d'o\`u
l'on d\'eduit le r\'esultat annonc\'e.)

\subsubsection{}
\label{XIII.2.0.1}
Dans le cas o\`u $R$ est strictement local, on a la
caract\'erisation simple suivante: la $K$-alg\`ebre $L$ est
mod\'er\'ement ramifi\'ee sur~$R$ si et seulement si les
$[L_i:K]$ sont premiers \`a $p$. De plus, si $L$ est
mod\'er\'ement ramifi\'ee sur~$R$, les $L_i$ sont des extensions
cycliques de~$K$. En effet, lorsque $R$ est strictement local, $I_i$
est \'egal au groupe de Galois
de~$L'_i$ sur~$K$. Comme on vient de le rappeler $I_i$ est extension
d'un groupe cyclique d'ordre premier \`a $p$ par un $p$-groupe. Si
l'on suppose $L'_i$ mod\'er\'ement ramifi\'ee sur~$R$, $I_i$ est
alors un groupe cyclique d'ordre premier \`a $p$. Il en r\'esulte
que $[L_i:K]$ est premier \`a $p$ et que l'on a $L_i =
L'_i$. Inversement, si $[L_i:K]$ est premier \`a $p$, $I_i$ ne peut
contenir de~$p$-sous-groupe distingu\'e non trivial; $I_i$ est donc
un groupe cyclique d'ordre premier \`a $p$, ce qui prouve que $L_i$
est mod\'er\'ement ramifi\'ee sur~$R$.

\subsubsection{}
\label{XIII.2.0.2}
Soient $R$ un anneau de valuation discr\`ete de corps des fractions
$K$, $L$ une $K$-alg\`ebre \'etale, et soient $\overline{R}$ un
localis\'e strict de~$R$, $\overline{K}$ son corps de fractions,
$\overline{L} = L \otimes_K \overline{K}$. Alors, pour que $L$ soit
mod\'er\'ement ramifi\'ee sur~$R$, il faut et il suffit que
$\overline{L}$ soit mod\'er\'ement ramifi\'ee sur~$\overline{R}$. On se ram\`ene en effet au cas o\`u $L$ est un
corps. Soient alors $\overline{L} = \prod_i \overline{L}_i$, o\`u
les $\overline{L}_i$ sont des corps extensions de~$\overline{K}$; si
$L'$ est l'extension galoisienne engendr\'ee par $L$, et, si
$\overline{L'} = L' \otimes_K \overline{K}$,
on a de m\^eme une d\'ecomposition de~$\overline{L'}$ en produit
de corps,
$\overline{L'} = \prod_j \overline{L'_j}$
et chaque $\overline{L}_i$ est sous-extension d'au moins l'un des
$\overline{L'_j}$. Comme $L'$ est une extension galoisienne de~$K$,
les $\overline{L'_j}$ sont des extensions galoisiennes de
$\overline{K}$. Supposons $L$ mod\'er\'ement ramifi\'ee sur~$R$;
comme le groupe de Galois de~$\overline{L'_j}|K$ est isomorphe au
groupe d'inertie de~$L'|K$, les $L'_j$ sont aussi mod\'er\'ement
ramifi\'ees sur~$\overline{R}$, et il en est donc de m\^eme des
$\overline{L}_i$. Inversement, supposons $\overline{L}$
mod\'er\'ement ramifi\'ee sur~$\overline{R}$. Pour chaque $j$,
soit $v_j$ la valuation discr\`ete de~$\overline{L'_j}$ qui prolonge
la valuation de~$\overline{K}$ et notons encore $v_j$ la valuation
induite sur~$L'$. Quand $j$ varie, $v_j$ parcourt l'ensemble des
valuations de~$L$ qui prolongent la valuation de~$K$. Soient $G =
\Gal(L'|K)$, $H=\Gal(L'|L)$, $I_j$ le groupe d'inertie de~$L'|K$ en
$v_j$, $J_j$ le groupe d'inertie de~$L'|L$ en $v_j$. Le groupe $I_i$
est extension d'un groupe cyclique d'ordre premier \`a $p$ par un
$p$-groupe $P_j$. Comme les $\overline{L}_i$ sont mod\'er\'ement
ramifi\'ees sur~$\overline{R}$, $I_j/J_j$ est d'ordre premier \`a
$p$, donc on a $P_j \subset J_j$. Par suite le groupe $H$ contient
tous les $P_j$ donc aussi le groupe $P$ engendr\'e par
les $P_j$ pour $j$ variable. Mais le groupe~$P$ est invariant dans $G$
car un automorphisme int\'erieur de~$G$ transforme les~$I_j$ entre
eux donc aussi les $P_j$ entre eux. Il en r\'esulte que $P$ est un
sous-groupe de~$H$ distingu\'e dans~$G$, donc, puisque $L'$ est
l'extension galoisienne engendr\'ee par $L$, que l'on a $P=1$, ce
qui prouve que $L$ est mod\'er\'ement ramifi\'ee sur~$R$.

Soient plus g\'en\'eralement $R \to R'$ un morphisme d'anneaux de
valuation discr\`ete tel que l'image d'une uniformisante $\pi$ de
$R$ soit une uniformisante $\pi'$ de~$R'$ et que l'extension
r\'esiduelle $\kres(R')$ soit une extension s\'eparable de
$\kres(R)$. Soient $K$ le corps des fractions de~$R$, $K'$ le corps des
fractions de~$R'$, $L$ une $K$-alg\`ebre \'etale, $L' = L
\otimes_K K'$. Alors, pour que $L$ soit mod\'er\'ement
ramifi\'ee sur~$R$, il faut et il suffit que $L'$ soit
mod\'er\'ement ramifi\'ee sur~$R'$. On peut en effet supposer
$R$ et $R'$ strictement locaux. D'apr\`es~\ref{XIII.2.0.1} il suffit
de prouver que, lorsque $L$ est un corps, il en est de m\^eme de~$L'$. Soient $\tilde{R}$ le normalis\'e de~$R$ dans $L$,
$\tilde{\pi}$ une uniformisante de~$\tilde{R}$, $\tilde{R}' =
\tilde{R} \otimes_R R'$. L'extension $\kres(\tilde{R})|\kres(R)$ \'etant
radicielle et l'extension $\kres(R')|\kres(R)$ \'etale,
$\kres(\tilde{R})\otimes_{\kres(R)}\kres(R')$ est un corps [EGA IV 4.3.2 et
4.3.5]. Ceci prouve que $R'$ est un anneau local, et, comme $\pi$ a
pour image $\pi'$ dans $R'$, on a $\kres(\tilde{R}')= R'/(\tilde{\pi})$;
par suite $R'$ est un anneau de valuation discr\`ete \cite[ch.\, I \S 2
prop.\, 2]{XIII.5} donc $L'$ est un corps.

\subsubsection{}
\label{XIII.2.0.3}
Par r\'eduction au cas strictement local, on voit qu'une
sous-alg\`ebre d'une alg\`ebre mod\'er\'ement ramifi\'ee est
mod\'er\'ement ramifi\'ee, que le produit tensoriel de deux
alg\`ebres mod\'er\'ement ramifi\'ees est mod\'er\'ement
ramifi\'e, qu'une alg\`ebre mod\'er\'ement ramifi\'ee le
reste apr\`es extension de l'anneau de valuation discr\`ete,
qu'une alg\`ebre qui devient mod\'er\'ement ramifi\'ee
apr\`es une extension mod\'er\'ement ramifi\'ee est
mod\'er\'ement ramifi\'ee.

\subsection{}
\label{XIII.2.1}
Soient
$X$ un $S$-sch\'ema, $D$ un diviseur $\geq 0$ sur~$X$.
Rappelons (SGA~5 II~4.2) qu'on dit que $D$ est strictement \`a
\emph{croisements normaux relativement à}~$S$
\index{croisements normaux (diviseur \`a)|hyperpage}\index{diviseur \`a croisements normaux|hyperpage}\index{strictement \`a croisements normaux (diviseur)|hyperpage}s'il existe une famille finie $(f_i)_{i\in I}$ d'\'el\'ements de
$\Gamma(X,\cal{O}_X)$, telle que l'on ait $D=\sum_{i\in
I}\divisor(f_i)$ et que la condition suivante soit r\'ealis\'ee:

\setcounter{subsubsection}{-1}
\subsubsection{}
\label{XIII.2.1.0}
Pour tout point $x$ de~$\Supp D$, $X$ est lisse sur~$S$ en $x$, et, si
l'on note $I(x)$ l'ensemble des $i\in I$ tels que $f_i(x)=0$, le
sous-sch\'ema $V((f_i)_{i\in I(x)})$ est lisse sur~$S$ de
codimension
$\card I(x)$
dans~$X$.

Le diviseur $D$ est dit \`a \emph{croisements normaux relativement
à} $S$ si, localement sur~$X$ pour la topologie \'etale, il est
strictement \`a croisements normaux.

Soit $D$ un diviseur \`a croisements normaux relativement
\`a~$S$. On pose $Y=\Supp D$, $U=X-Y$, et on note $i\colon U\to X$
l'immersion canonique. Pour tout point g\'eom\'etrique
$\overline{s}$ de~$S$ et pour tout point maximal $y$ de la fibre
g\'eom\'etrique $Y_{\overline{s}}$, l'anneau
$R=\cal{O}_{X_{\overline{s}},y}$ est un anneau de valuation
discr\`ete.

Dans la suite de ce num\'ero, nous utiliserons la d\'efinition
technique suivante:

\begin{subdefinition}
\label{XIII.2.1.1}
Soit $F$ un faisceau d'ensembles sur~$U$. On dit que $F$ est
\emph{modérément ramifié sur}~$X$
\index{mod\'er\'ement ramifi\'e (faisceau)|hyperpage}(\emph{le long de}~$D$) \emph{relativement à}~$S$ si, pour tout
point g\'eom\'etrique $\overline{s}$ de~$S$, la condition suivante
est satisfaite:

Pour tout point maximal $y$ de~$Y_{\overline{s}}$, la restriction de
$F$ au corps des fractions $K$ de~$\cal{O}_{X_{\overline{s}}}$ est
repr\'esentable par le spectre d'une $K$-alg\`ebre \'etale $L$,
mod\'er\'ement ramifi\'ee sur~$\cal{O}_{X_{\overline{s},y}}$.
\end{subdefinition}

Le plus souvent, quand il ne pourra en r\'esulter de confusion, nous
omettrons la mention de~$D$ dans la terminologie.

\begin{subdefinition}
\label{XIII.2.1.2}
Si
$F$ est un faisceau en groupes sur~$U$, mod\'er\'ement
ramifi\'e sur~$X$ relativement \`a~$S$, on d\'esigne par
$$
\H^1_{\tame}(U,F)
$$
\label{indnot:mc}le sous-ensemble de~$\H^1(U,F)$ form\'e des classes de torseurs sous
$F$ qui sont mod\'er\'ement ramifi\'es sur~$X$ relativement
\`a~$S$.
\end{subdefinition}

Soit
$$
\xymatrix{{U}\ar[r]^i\ar[d]_g&{X}\ar[dl]^f\\{T}}
$$
un diagramme commutatif de~$S$-sch\'emas, avec $i$ comme
dans~\ref{XIII.2.1}; on d\'esigne par
$$
\R^1_{\tame} g_* F
$$
\label{indnot:md}le faisceau sur~$T$ associ\'e au pr\'efaisceau
$T'\mto\H^1_{\tame}(U',F)$, o\`u $T'$ parcourt les sch\'emas
\'etales au-dessus de~$T$ et o\`u $U'=U\times_T T'$; $\R^1_{\tame} g_*F$
est un sous-faisceau de~$\R^1g_*F$.

Notons que, si $g$ est coh\'erent, si $\overline{t}$ est un point
g\'eom\'etrique de~$T$, $\overline{T}$ le localis\'e strict de
$T$ en $\overline{t}$, $\overline{U}=U\times_T\overline{T}$, on a un
isomorphisme
\begin{equation*}
\label{eq:XIII.2.1.2.1}
\tag{\thesubsubsection.1} (\R^1_{\tame} g_*F)_{\overline{t}}\simeq
\H^1_{\tame}(\overline{U},\overline{F})
\end{equation*}

\subsubsection{}
\label{XIII.2.1.3}
Soit $C_t((U,X)/S)$ ou simplement $C_t$
\label{indnot:me}la cat\'egorie des rev\^etements \'etales de~$U$ qui sont
mod\'er\'ement ramifi\'es sur~$X$ relativement
\`a~$S$. Supposons $U$ connexe et soit $a$ un point
g\'eom\'etrique de~$U$. Soit $\Gamma_t$ le foncteur qui, \`a un
rev\^etement \'etale~$U'$ de~$U$ mod\'er\'ement ramifi\'e
sur~$X$ relativement \`a~$S$ fait correspondre l'ensemble des points
g\'eom\'etriques de~$U'$ au-dessus de~$a$. De~\ref{XIII.2.0}
r\'esulte que le couple $(C_t,\Gamma_t)$ satisfait aux axiomes
($G_1$) \`a ($G_6$) de V~\ref{V.4}. Par suite $\Gamma_t$ est
repr\'esentable par un pro-objet qu'on appelle le
\emph{revêtement universel modérément ramifié
de~$(U,X)$ relativement à~$S$ ponctué en~$a$}.
\index{mod\'er\'ement ramifi\'e (rev\^etement universel)|hyperpage}\index{revetement universel moderement ramifie@rev\^etement universel mod\'er\'ement ramifi\'e|hyperpage}Le groupe oppos\'e au groupe des $U$-automorphismes du
rev\^etement universel mod\'er\'ement ramifi\'e est appel\'e
le \emph{groupe fondamental modérément ramifié}
\index{groupe fondamental mod\'er\'ement ramifi\'e|hyperpage}\index{mod\'er\'ement ramifi\'e (groupe fondamental)|hyperpage}et not\'e
$$
\pi_1^{\tame}((U,X)/S,a)\text{ ou simplement }\pi_1^{\tame}(U,a)\text{ ou
même }\pi_1^{\tame}(U).
$$
\label{indnot:mf}C'est \'evidemment un quotient du groupe fondamental $\pi_1(U,a)$
(V~\ref{V.6.9}).

\subsubsection{}
\label{XIII.2.1.4}
Soient $F$ un faisceau en groupes sur~$U$, $P$ un torseur \`a droite
de groupe $F$, $Q$ un torseur \`a gauche de groupe $F$, et supposons
$P$ et $Q$ mod\'er\'ement ramifi\'es sur~$X$ relativement \`a~$S$. Alors il en est de m\^eme de~$P\stackrel{F}{\wedge}Q$. On se
ram\`ene en effet \`a montrer que, si $R$ est un anneau de
valuation discr\`ete de corps des fractions $K$ et si~$F$ est un
sch\'ema en groupes \'etale fini sur~$K$, et, si $P$ et $Q$ sont
deux torseurs sous~$F$ mod\'er\'ement ramifi\'es sur~$R$, alors
il en est de m\^eme de~$P\stackrel{F}{\wedge} Q$. Or
$T=P\stackrel{F}{\wedge} Q$ est un quotient de~$P\times_K Q$. Si $L$,
$M$, $N$ d\'esignent les $K$-alg\`ebres repr\'esentant
respectivement $T$, $P$, $Q$, alors $L$ est une sous-alg\`ebre de
$M\otimes_K N$, et il r\'esulte de~\ref{XIII.2.0.3} que $L$ est
mod\'er\'ement ramifi\'ee sur~$R$.

On d\'eduit de ce qui pr\'ec\`ede que, si $F$ est un faisceau en
groupes sur~$U$ et s'il existe un torseur de groupe $F$,
mod\'er\'ement ramifi\'e sur~$X$ relativement \`a~$S$, alors
$F$ est mod\'er\'ement ramifi\'e sur~$X$ relativement \`a~$S$. En effet le torseur $P^{\circ}$ oppos\'e de~$P$ est
mod\'er\'ement ramifi\'e sur~$X$ relativement \`a~$S$,
puisqu'il est isomorphe \`a $P$ en tant que faisceau d'ensembles. Si
$\leftexp{P\mkern-4mu}{F}$ est le groupe tordu de~$F$ par $P$, on a un
isomorphisme
$$
F\simeq P^{\circ}\stackrel{\leftexp{P\mkern-4mu}{F}}{\wedge}P
$$
et par suite $F$ est mod\'er\'ement ramifi\'e sur~$X$
relativement \`a~$S$.

On voit comme pr\'ec\'edemment que, si $F\to F'$ est un morphisme
de faisceaux en groupes sur~$U$, mod\'er\'ement ramifi\'es sur~$X$ relativement \`a~$S$, et si $P$ est un torseur sous $F$
mod\'er\'ement ramifi\'e sur~$X$ relativement
\`a~$S$, alors le torseur $P'$ d\'eduit de~$P$ par l'extension du
groupe structural $F\to F'$ est mod\'er\'ement ramifi\'e sur~$X$
relativement \`a~$S$.

En particulier le morphisme canonique
$$
\H^1(U,F)\to \H^1(U,F')
$$
donne par restriction \`a $\H^1_{\tame}(U,F)$ un morphisme canonique
$$
\H^1_{\tame}(U,F)\to\H^1_{\tame}(U,F')
$$

\subsubsection{}
\label{XIII.2.1.5}
Soient $S'\to S$ un morphisme et notons $U'$ (\resp $X'$, etc.)
l'image inverse de~$U$ (\resp $X$, etc.) sur~$S'$. Si $F$ est un
faisceau d'ensembles sur~$U$ mod\'er\'ement ramifi\'e sur~$X$
relativement \`a~$S$, il r\'esulte de la
d\'efinition~\ref{XIII.2.1.1} et de~\ref{XIII.2.0.3} que $F'$ est
mod\'er\'ement ramifi\'e sur~$X'$ relativement \`a~$S'$.

Si maintenant $F$ est un faisceau en groupes sur~$U$, l'image inverse
sur~$S'$ d'un torseur sous $F$ mod\'er\'ement ramifi\'e sur~$X$
relativement \`a~$S$ est un torseur sous $F'$ mod\'er\'ement
ramifi\'e sur~$X'$ relativement \`a~$S'$. En particulier on a un
foncteur canonique
\begin{equation*}
\label{eq:XIII.2.1.5.1}
\tag{\thesubsubsection.1} C_t((U,X)/S)\to C_t((U',X')/S').
\end{equation*}
Supposons $U$ et $U'$ connexes et soient $a$ un point
g\'eom\'etrique de~$U$, $a'$ un point g\'eom\'etrique de~$U'$
au-dessus de~$a$; on d\'eduit de ce qui pr\'ec\`ede un morphisme
canonique
\begin{equation*}
\label{eq:XIII.2.1.5.2}
\tag{\thesubsubsection.2} \pi_1^{\tame}(U',a')\to\pi_1^{\tame}(U,a).
\end{equation*}

Si $S'\to S$ est un morphisme et $h\colon T'\to T$ la projection
canonique, le morphisme
$$
h^*(\R^1g_*F)\to\R^1g_*'F'
$$
donne par restriction un morphisme canonique
\begin{equation*}
\label{eq:XIII.2.1.5.3}
\tag{\thesubsubsection.3} h^*(\R^1_{\tame}g_*F)\to\R^1_{\tame}g_*'F'
\end{equation*}

\subsubsection{}
\label{XIII.2.1.6}
Soit
$F$ un faisceau en groupes sur~$U$, mod\'er\'ement
ramifi\'e sur~$X$ relativement \`a~$S$. Les notations \'etant
celles de~\ref{XIII.2.1.2}, on a des suites exactes canoniques:
\begin{equation*}
\label{eq:XIII.2.1.6.1}
\tag{\thesubsubsection.1}
\begin{array}{ccccccc}
1 & \longrightarrow & \H^1(X,i_*F)& \longrightarrow & \H^1_{\tame}(U,F) &
\longrightarrow & \H^0(X,\R^1_{\tame} i_*F)\\[10pt] 1 &
\longrightarrow & \R^1f_*(i_*F) & \longrightarrow & \R^1_{\tame}g_*F &
\longrightarrow & f_*(\R^1_{\tame}i_*F).\end{array}
\end{equation*}

La premi\`ere s'obtient \`a partir de la suite exacte
(SGA~4 III~3.2):
$$
1\to \H^1(X,i_*F)\to \H^1(U,F)\to
\H^0(X,\R^1i_*F).
$$
Il suffit en effet de montrer que l'image de~$\H^1(X,i_*F)$ dans
$\H^1(U,F)$ est en fait contenue dans $\H^1_{\tame}(U,F)$ et que l'image de
$\H^1_{\tame}(U,F)$ dans $\H^0(X,\R^1i_*F)$ est en fait contenue dans
$\H^0(X,\R^1_{\tame}i_*F)$. Or l'image inverse sur~$U$ d'un torseur sous
$i_*F$ est un torseur sous $i^*i_*F$ qui est \'evidemment
mod\'er\'ement ramifi\'e sur~$X$ relativement \`a~$S$, donc il
en est de m\^eme apr\`es l'extension du groupe structural
$i^*i_*F\to F$, ce qui prouve l'existence de la fl\`eche
$\H^1(X,i_*F)\to\H^1_{\tame}(U,F)$. Le fait que l'image de~$\H^1_{\tame}(U,F)$
dans $\H^0(X,\R^1i_*F)$ soit contenue dans $\H^0(X,\R^1_{\tame}i_*F)$
r\'esulte aussit\^ot de la d\'efinition de~$\R^1_{\tame}i_*F$. Ceci
prouve l'existence de la premi\`ere suite exacte, et la deuxi\`eme
s'en d\'eduit par localisation.

\subsection{}
\label{XIII.2.2}
On conserve les notations de~\ref{XIII.2.1}. Nous allons d\'efinir
une notion d'objet mod\'er\'ement ramifi\'e d'un champ $\Phi$
sur~$U$ lorsque celui-ci est donn\'e, localement \emph{sur} $X$ et
sur~$S$ pour la topologie \'etale, comme \emph{image inverse d'un
champ} $\Psi$ \emph{sur}~$S$.

Soit d'abord $G$ une gerbe sur~$U$ et supposons donn\'e un morphisme
\'etale surjectif $S_1\to S$, un morphisme \'etale surjectif
$X_2\to X\times_S S_1$, une gerbe triviale $H$ sur~$S_1$ et un
isomorphisme
$$
G|{U_2}\to H|{U_2},
$$
o\`u
$U_1=U\times_X X_1$, $U_2=U\times_X X_2$. Quand on choisit une
trivialisation de~$H|{X_2}$, l'isomorphisme ci-dessus identifie
$G|{U_2}$ au champ des torseurs sous un faisceau en groupes $F$. On
dit qu'un \'el\'ement $x$ de~$G_U$ est mod\'er\'ement
ramifi\'e sur~$X$ relativement
\`a~$S$
si la restriction de~$x$
\`a $U_2$ est un torseur mod\'er\'ement ramifi\'e sur~$X$
relativement \`a~$S$. D'apr\`es~\ref{XIII.2.1.4} cette notion ne
d\'epend pas de la fa\c con dont on a trivialis\'e~$H|X_2$.

Soit maintenant $\Phi$ un champ sur~$U$ et supposons donn\'es un
morphisme \'etale surjectif $S_1\to S$, un morphisme \'etale
surjectif $X_2\to X\times_S S_1$, un champ $\Psi$ sur~$S_1$ et un
isomorphisme
$$
i\colon \Phi|U_2\to \Psi|U_2.
$$
Soit $x$ un \'el\'ement de~$\Phi_U$, $G_x$ la sous-gerbe maximale
de~$\Phi$ engendr\'ee par $x$
\cite[III~2.1.7]{XIII.1}, $S\Phi$ le faisceau des
sous-gerbes maximales de~$\Phi$. L'isomorphisme $i$ induit un
isomorphisme
$$
S\Phi|U_2\to S\Psi|U_2.
$$
Il r\'esulte de~\ref{XIII.5.7} que, quitte \`a remplacer $S_1$ par
une extension \'etale surjective, on a une unique sous-gerbe
maximale $H$ de~$\Psi$, que l'on peut supposer triviale, telle que $i$
d\'efinisse un isomorphisme
$$
G_x|U_2\to H|U_2.
$$
On dit que l'\'el\'ement $x$ est mod\'er\'ement ramifi\'e
sur~$X$ relativement \`a~$S$ s'il l'est en tant qu'\'el\'ement
de~$G_x$ munie de l'isomorphisme ci-dessus.

\subsubsection{}
\label{XIII.2.2.1}
Soit $\Phi$ un champ sur~$U$ donn\'e, localement sur~$X$ et $S$,
comme image inverse d'un champ sur~$S$, et soit
$$
\xymatrix{{U}\ar[r]^i\ar[d]_g&{X}\ar[dl]^f\\{T}}
$$
un diagramme comme dans~\ref{XIII.2.1.2}. Pour tout sch\'ema $T'$
\'etale sur~$T$, si $U'=U\times_T T'$, on consid\`ere le
sous-ensemble $(g_*^{\tame}\Phi)_{T'}$
\label{indnot:mg}de~$(g_*\Phi)_{T'}=\Phi_{U'}$
form\'e
des \'el\'ements de~$\Phi_{U'}$ qui sont mod\'er\'ement
ramifi\'es
\index{element moderement@\'el\'ement mod\'er\'ement ramifi\'e d'un champ|hyperpage}sur~$X$ relativement \`a~$S$. On appelle image directe
mod\'er\'ement ramifi\'ee
\index{image directe mod\'er\'ement ramifi\'ee d'un champ|hyperpage}\index{mod\'er\'ement ramifi\'ee (image directe d'un champ)|hyperpage}de~$\Phi$ par~$g$, et on note
$$
g_*^{\tame}\Phi
$$
la sous-cat\'egorie pleine de~$g_*\Phi$ dont les objets au-dessus
d'un sch\'ema $T'$ \'etale sur~$T$ sont les \'el\'ements de
$(g_*^{\tame}\Phi)_{T'}$. Il est clair que $g_*^{\tame}\Phi$ est un sous-champ de
$g_*\Phi$.

\subsubsection{}
\label{XIII.2.2.2}
Par r\'eduction au cas d'un champ de torseurs, on voit que, si
$\Phi$ est un champ sur~$U$ qui est localement pour la topologie
\'etale de~$S$ et $X$ image inverse d'un champ sur~$S$, le morphisme
canonique
$$
h^*(g_*\Phi)\to g_*'\Phi'
$$
donne par restriction un morphisme canonique
$$
h^*(g_*^{\tame}\Phi)\to g_*^{\prime\tame}\Phi'
$$

\begin{remarques}
\label{XIII.2.3}

a) Si $F$ est un faisceau d'ensembles localement constant
constructible sur~$U$, pour que $F$ soit mod\'er\'ement
ramifi\'e sur~$X$ relativement \`a~$S$, il suffit que la condition
de~\ref{XIII.2.1.1} soit satisfaite pour les points
g\'eom\'etriques de~$S$ au-dessus des points maximaux de~$S$. Pour
le voir on peut supposer le diviseur $D$ strictement \`a croisements
normaux. Le faisceau $F$ est repr\'esentable par un rev\^etement
\'etale $V$ de~$U$. Si $\overline{s}$ est un point
g\'eom\'etrique de~$S$, $y$ un point maximal de
$Y_{\overline{s}}$, on note $\overline{S}$ le localis\'e strict de
$S$ en $\overline{s}$, $\overline{X}$ le localis\'e strict de~$X$ en
$\overline{y}$, $\overline{U}=U\times_X \overline{X}$,
$\overline{V}=V\times_X\overline{X}$. Si la condition
de~\ref{XIII.2.1.1} est satisfaite aux points g\'eom\'etriques
au-dessus des points maximaux de~$S$, il r\'esulte de~\ref{XIII.5.5}
plus bas que $\overline{V}$ est un rev\^etement de~$\overline{U}$
mod\'er\'ement ramifi\'e sur~$\overline{X}$ relativement \`a
$\overline{S}$; par suite $V$ est un rev\^etement \'etale de~$U$
mod\'er\'ement ramifi\'e sur~$X$ relativement \`a~$S$.

b)
Soit $F$ un faisceau en groupes sur~$U$ mod\'er\'ement
ramifi\'e sur~$X$ relativement \`a~$S$. Si $\overline{s}$ est un
point g\'eom\'etrique de~$S$, $y$ un point maximal de
$Y_{\overline{s}}$, on note $K$ le corps des fractions de
$\cal{O}_{X_{\overline{s}},y}$. Supposons que, pour tout point
$\overline{s}$ et pour tout point $y$, la $K$-alg\`ebre $L$ dont le
spectre repr\'esente $F|K$ soit de rang premier \`a la
caract\'eristique r\'esiduelle $p$ de
$\cal{O}_{X_{\overline{s}}}$. On dira parfois, par abus de langage,
que $F$ est premier aux caract\'eristiques r\'esiduelles de~$S$.
\index{premier aux caract\'eristiques r\'esiduelles (faisceau)|hyperpage}Lorsqu'il en est ainsi, tout torseur $P$ sous $F$ est
mod\'er\'ement ramifi\'e sur~$X$ relativement \`a~$S$.

Soient en effet $\overline{R}$ le localis\'e strict de
$\cal{O}_{X_{\overline{s}},y}$ en $\overline{y}$, $\overline{K}$ son
corps des fractions, $\overline{F}$ l'image inverse de~$F$ sur~$\overline{K}$. Montrons que l'on peut supposer $F$ constant. Comme
$F$ est mod\'er\'ement ramifi\'e sur~$X$ relativement \`a~$S$,
$\overline{F}$ est repr\'esentable par le spectre d'une
$\overline{K}$-alg\`ebre $L=\prod L_i$, o\`u les $L_i$ sont des
extensions de~$\overline{K}$ de degr\'e premier \`a $p$. On peut
donc trouver une extension $K'$ de~$\overline{K}$ de degr\'e premier
\`a $p$ telle que $\overline{F}|K'$ soit un faisceau
constant. D'apr\`es~\ref{XIII.2.0.3}, pour prouver que
$P|\overline{K}$ est mod\'er\'ement ramifi\'e sur~$\overline{R}$, il suffit de voir que $P|K'$ est mod\'er\'ement
ramifi\'e sur la cl\^oture int\'egrale de~$\overline{R}$ dans
$K'$, d'o\`u la r\'eduction au cas o\`u $\overline{F}$ est
constant. Supposons d\'esormais $\overline{F}$ constant. La
$\overline{K}$-alg\`ebre $H$ qui repr\'esente $P|\overline{K}$ est
alors produit d'extensions $H_i$ de~$\overline{K}$ isomorphes entre
elles. Comme le rang de~$H$ est premier \`a $p$, il en est de
m\^eme de~$[H_1:\overline{K}]$, ce qui prouve que $H$ est
mod\'er\'ement ramifi\'e sur~$X$ relativement \`a~$S$.

c) Soient $X$ un sch\'ema r\'egulier, $D$ un diviseur \`a
croisements normaux de~$X$ (SGA~5 I~3.1.5), $U=X-\Supp D$, $F$ un
faisceau d'ensembles sur~$U$. Si $y$ est un point maximal de~$\Supp
D$, on d\'esigne par $K$ le corps des fractions de
$\cal{O}_{X,y}$. On dit que $F$ est \emph{modérément
ramifié relativement à}~$D$,
\index{mod\'er\'ement ramifi\'e (faisceau)|hyperpage}si, pour tout point maximal $y$ de~$\Supp D$, $F|K$ est
repr\'esentable par une $K$-alg\`ebre mod\'er\'ement
ramifi\'ee sur~$\cal{O}_{X,y}$.
\end{remarques}

\begin{theoreme}
\label{XIII.2.4}
Soient
$f\colon X\to S$ un $S$-sch\'ema, $D$ un diviseur sur~$X$
\`a croisements normaux relativement \`a~$S$ \eqref{XIII.2.1},
$Y=\Supp D$, $U=X-Y$, $i\colon U\to X$ l'immersion canonique. Soit $F$
un faisceau d'ensembles (\resp de groupes) sur~$U$, satisfaisant \`a
l'une des conditions suivantes:
\begin{enumerate}
\item[a)] $F$ est localement pour la topologie \'etale \emph{sur}
$X$ et sur~$S$ l'image inverse d'un faisceau d'ensembles (\resp d'un
faisceau en groupes constructible) sur~$S$.
\item[b)] $F$ est localement constant constructible sur~$U$ et
mod\'er\'ement ramifi\'e sur~$X$ relativement \`a~$S$.
\end{enumerate}
Alors on a les conclusions suivantes:
\begin{enumerate}
\item[1)]$(F,i)$ est cohomologiquement propre relativement \`a~$S$
en dimension $\leq 0$ (\resp pour tout morphisme $h\colon S'\to S$, si
$i'\colon U'\to X'$ est l'image inverse de~$i$ sur~$S'$, si $F'=F|U'$
et si $k=h_{(X)}$, le morphisme canonique
$$
\Psi\colon k^*(\R^1_{\tame}i_*F)\to\R^1_{\tame}i_*'F'
$$
est un isomorphisme).

Si $F$ est un faisceau en groupes premier aux caract\'eristiques
r\'esiduelles de~$S$ (\ref{XIII.2.3}.b)) (mod\'er\'ement
ramifi\'e sur~$X$ relativement \`a~$S$), alors $(F,i)$ est
cohomologiquement propre relativement \`a~$S$ en dimension $\leq 1$.
\item[2)]Si $F$ est un faisceau d'ensembles (\resp de groupes)
constructible, $i_*F$ (\resp $\R^1_{\tame}i_*F$) est constructible.
\end{enumerate}
\end{theoreme}

\subsubsection*{D\'emonstration}
Pour tout $S$-sch\'ema $S'$, on consid\`ere le diagramme suivant
dont tous les carr\'es sont cart\'esiens:
$$
\xymatrix{{U}\ar[dd]_-g\ar[dr]^-i& & &{U'}\ar[lll]\ar[dd]_{g'}\ar[dr]^-{i'}& \\ &{X}\ar[dl]_{f\mkern-8mu}& & &{X'}\ar[lll]_{k}\ar[dl]_{f'\mkern-10mu}\\{S}& & &{S'}\ar[lll]_-h& & }
$$
Comme
la question est locale sur~$X$ pour la topologie \'etale, on peut
supposer que $D$ est un diviseur strictement \`a croisements normaux
relativement \`a~$S$ \eqref{XIII.2.1}; de plus, quitte \`a
restreindre $X$ \`a un voisinage de~$Y$, on peut supposer $X$ lisse
sur~$S$.

\subsubsection*{D\'emonstration de~\ref{XIII.2.4}~\textup{1)}}
\subsubsection{}
\label{XIII.2.4.1}
\emph{Cas d'un faisceau d'ensembles satisfaisant à} a). On peut
supposer que l'on~a $F=g^*G$, o\`u $G$ est un faisceau sur~$S$. Il
r\'esulte alors de (SGA~4 XVI~3.2) que le morphisme canonique
\begin{equation*}
\label{eq:XIII.2.4.1.1}
\tag{\thesubsubsection.1} f^*G\to i_*F
\end{equation*}
est un isomorphisme. Pour tout $S$-sch\'ema $h\colon S'\to S$, on a
de m\^eme un isomorphisme $f'^*G'\to i_*'F'$; par suite le morphisme
canonique
$$
\varphi\colon k^*(i_*F)\to i_*'F'
$$
s'identifie \`a l'isomorphisme naturel
$$
k^*f^*G\simeq f'^*G'
$$

\subsubsection{}
\label{XIII.2.4.2}
\emph{Cas d'un faisceau d'ensembles satisfaisant à} b). On doit
montrer que $\varphi$ est un isomorphisme et il suffit pour cela de
voir qu'il en est ainsi en chaque point g\'eom\'etrique
$\overline{x}'$ de~$X'$. Soit $\overline{S}$ (\resp $\overline{X}$,
\resp $\overline{S'}$, \resp $\overline{X'}$) le localis\'e strict
de~$S$ (\resp~$X$, \resp~$S'$, \resp $X'$) en $\overline{x}'$ et posons
$\overline{U}=U_{(\overline{X})}$,
$\overline{U'}=U_{(\overline{X'})}$, etc. Le morphisme
$\varphi_{\overline{x}}$ s'identifie au morphisme canonique
$$
\overline{\varphi}\colon \H^0(\overline{U},\overline{F})\to
\H^0(\overline{U'}, \overline{F'})
$$
On peut trouver un rev\^etement principal $V$ de~$\overline{U}$, du
type figurant dans~\ref{XIII.5.4}, tel que l'image inverse de
$\overline{F}$ sur~$V$ soit un faisceau constant de valeur $C$. Si
$\Pi$ est le groupe de Galois de~$V$ sur~$\overline{U}$, $\Pi$
op\`ere sur~$\overline{F}|V$, et l'on a
\begin{equation*}
\label{eq:XIII.2.4.2.1}
\tag{\thesubsubsection.1}
\H^0(\overline{U},\overline{F}))\simeq \H^0(V,C_V)^{\Pi}
\end{equation*}
\label{indnot:mh}o\`u le deuxi\`eme membre d\'esigne l'ensemble des
\'el\'ements de~$\H^0(V,C_V)$ invariants
sous~$\Pi$. Comme $V'=V\times_{\overline{U}}\overline{U'}$ est une
rev\^etement principal de~$\overline{U'}$ de groupe de Galois
$\Pi'\simeq \Pi$, on voit que le morphisme $\overline{\varphi}$
s'obtient, en prenant les invariants sous $\Pi$, \`a partir du
morphisme canonique
$$
\H^0(V,C_V)\to \H^0(V',C_{V'}).
$$
Comme $V$ et $V'$ sont connexes \eqref{XIII.5.4}, ce morphisme, donc
aussi $\overline{\varphi}$, est un isomorphisme.

Notons que si de plus $F$ est un faisceau en groupes et si $P$ est un
torseur sur~$\overline{U}$ de groupe $\overline{F}$,
mod\'er\'ement ramifi\'e relativement \`a $D$, il r\'esulte
de la d\'emonstration pr\'ec\'edente et de~\ref{XIII.2.2} que
$(\leftexp{P}{\overline{F}},\overline{i})$ est cohomologiquement
propre relativement \`a~$S$ en dimension $\leq 0$.

\subsubsection{}
\label{XIII.2.4.3}
\emph{Cas d'un faisceau en groupes.} Pour montrer que $\Psi$ est un
isomorphisme, il suffit de prouver que, pour tout point
g\'eom\'etrique $\overline{y}'$ de~$Y'$, le morphisme
$$
\Psi_{\overline{y}'}\colon
(k^*(\R^1_{\tame}i_*F))_{\overline{y}'}\to (\R^1_{\tame}
i_*'F')_{\overline{y}'}
$$
est un isomorphisme. Or, d'apr\`es~\ref{XIII.2.1.2},
$\Psi_{\overline{y}'}$ s'identifie au morphisme canonique
$$
\overline{\Psi}\colon
\H^1_{\tame}(\overline{U},\overline{F})\to
\H^1_{\tame}(\overline{U'},\overline{F'}).
$$
Soient $\widetilde{U}$ le rev\^etement universel mod\'er\'ement
ramifi\'e de~$\overline{U}$ \eqref{XIII.2.1.3} et $\widetilde{F}$
l'image inverse de~$\overline{F}$ sur~$\widetilde{U}$. Il r\'esulte
de~\ref{XIII.5.7} dans le cas a) et de~\ref{XIII.5.5} dans le cas b),
que $\H^1_{\tame}(\overline{U},\overline{F})$ s'identifie au
sous-ensemble $\H^1(\overline{U},\overline{F})$ form\'e des classes
de~$\overline{F}$-torseurs dont l'image inverse sur~$\widetilde{U}$
est triviale. D'autre part un raisonnement
classique (\cf IX~\ref{IX.5}, p.\,\pageref{page-300})
montre que
l'ensemble des \'el\'ements de~$\H^1(\overline{U},\overline{F})$
dont l'image inverse sur~$\widetilde{U}$ est triviale s'identifie
\`a
$\H^1(\pi_1^{\tame}(\overline{U}),\H^0(\widetilde{U},\widetilde{F})).$
On obtient ainsi un isomorphisme canonique
\begin{equation*}
\label{eq:XIII.2.4.3.1} \tag{\thesubsubsection.1} {}
\H^1_{\tame}(\overline{U},\overline{F})\isomto
\H^1(\pi_1^{\tame}(\overline{U}),\H^0(\widetilde{U},\widetilde{F})).
\end{equation*}

Par suite le morphisme $\overline{\Psi}$ s'identifie au morphisme
canonique
$$
\H^1(\pi_1^{\tame}(\overline{U}),\H^0(\widetilde{U},\widetilde{F}))
\to
\H^1(\pi_1^{\tame}(\overline{U'}),\H^0(\widetilde{U}',\widetilde{F}')).
$$
Montrons
que ce morphisme est un isomorphisme. Le morphisme
$\pi_1^{\tame}(\overline{U'})\to\pi_1^{\tame}(\overline{U})$
est un isomorphisme d'apr\`es~\ref{XIII.5.6}, et il en est de
m\^eme du morphisme $\H^0(\widetilde{U},\widetilde{F})\to
\H^0(\widetilde{U}',\widetilde{F}')$. En effet, cela est \'evident
dans le cas b) car $\widetilde{F}$ est constant et $\widetilde{U}$ et
$\widetilde{U}'$ sont connexes. Dans le cas a), soit $\overline{G}$ un
faisceau en groupes constructible sur~$\overline{S}$ tel que l'on ait
$\overline{F}=\overline{g}^*\overline{G};$ comme les morphismes
$\widetilde{U}\to\overline{S}$ et
$\widetilde{U}'\to\overline{S'}$ sont $0$-acycliques
\eqref{XIII.5.7}, on a
$$
\H^0(\widetilde{U},\widetilde{F})\simeq
\H^0(\overline{S},\overline{G}) \simeq
\H^0(\overline{S'},\overline{G'})\simeq
\H^0(\widetilde{U}',\widetilde{F}'),
$$
ce qui entra\^ine que $\Psi$ est un isomorphisme. La derni\`ere
assertion de \ref{XIII.2.4}~1) r\'esulte de ce qui pr\'ec\`ede,
compte tenu de \ref{XIII.2.3}~b).

\subsubsection*{D\'emonstration de~\ref{XIII.2.4}~\textup{2)}}

Le cas d'un faisceau d'ensembles constructible satisfaisant \`a a)
r\'esulte aussit\^ot de \eqref{eq:XIII.2.4.1.1}. Soit $F=g^*G$ un
faisceau en groupes satisfaisant \`a a), o\`u $G$ est un faisceau
constructible; on peut supposer $S$ affine; soient $(S_j)_{j\in J}$
une famille finie de sous-sch\'emas ferm\'es r\'eduits de~$S$
dont la r\'eunion recouvre $S$, tels que l'image inverse de~$G$ sur~$S_j$ soit un faisceau localement constant. Compte tenu de
\ref{XIII.2.4}~1), il suffit, pour \'etablir que $\R^1_{\tame}i_*F$
est constructible, de voir qu'il en est ainsi apr\`es le changement
de base $S_j\to S$, pour chaque $j\in J$. On est donc
ramen\'e au cas~b) o\`u $F$ est localement constant.

On suppose d\'esormais que $F$ est un faisceau d'ensembles ou de
groupes satisfaisant \`a b). Comme la question est locale pour la
topologie \'etale sur~$X$, on peut supposer $X$ de pr\'esentation
finie sur~$S$, et, par passage \`a la limite, on peut supposer $X$
et $S$ noeth\'eriens.

Soit $D=\sum_{1\leq i\leq r}\divisor f_i$, o\`u, pour chaque point
$x$ de~$\Supp D$, si $I(x)$ est l'ensemble des $i$ tels que
$f_i(x)=0$, le sous-sch\'ema $V((f_i)_{i\in I(x)})$ est lisse sur~$S$ de codimension $\card I(x)$ dans~$X$. Soit $\mathcal{P}$
l'ensemble des parties de~$[ 1,r ]$ et, pour chaque $I\in\mathcal{P}$,
posons
$$
X_I=\textstyle\Big(\bigcap_{i\in I}V(f_i)\Big)\cap\Big(\bigcap_{i\not\in I}X_{f_i}\Big).
$$
Soit
$z$ un point de~$X_I$. Quitte \`a se restreindre d'abord \`a un
voisinage \'etale de~$z$, on peut trouver un ouvert $W$ de~$X$
contenant $z$ et un rev\^etement principal $V$ de~$U\cap W$,
mod\'er\'ement ramifi\'e sur~$W$ relativement \`a~$S$, du type
consid\'er\'e dans~\ref{XIII.5.6.1}, tel que l'image
r\'eciproque de~$F$ sur~$V$ soit un faisceau constant de
valeur~$C$. Soit $\pi$ le groupe de Galois du rev\^etement~$V$. Pour
tout point g\'eom\'etrique $\overline{x}$ de~$X_I$, on a alors
d'apr\`es \eqref{eq:XIII.2.4.2.1}
$$
\H^0(\overline{U},\overline{F})\simeq \H^0(V,C_V)^{\pi}.
$$
Il en r\'esulte que $i_*F|X_I\cap W$ est localement constant (SGA~4
IX~2.13) et par suite $i_*F$ est constructible.

Montrons enfin que si $F$ est un faisceau en groupes localement
constant, $\R^1_{\tame}i_*F|_{X_I}$ est constructible. Si
$\overline{x}$ est un point g\'eom\'etrique de~$X$, on a obtenu
dans \eqref{eq:XIII.2.4.3.1} l'expression
$$
\H^1_{\tame}(\overline{U},\overline{F})\isomto
\H^1(\pi^{\tame}_1(\overline{U}),\H^0(\widetilde{U},\widetilde{F})).
$$
Si $p$ est la caract\'eristique r\'esiduelle de~$\overline{X}$, on
a, d'apr\`es~\ref{XIII.5.6},
$\pi_1^{\tame}(\overline{U})=\prod_{l\neq p }\ZZ_l(1)^{\card I}$. D\'esignons par $\LL$ l'ensemble des nombres premiers qui
divisent l'ordre du groupe fini $\H^0(\widetilde{U},\widetilde{F})$ et
soit $K=\prod_{l\in\LL -\{p\}\cap\LL }\ZZ_l(1)^{\card I}$. Il
r\'esulte de \cite[I \S5 ex.2]{XIII.4} que l'on a
$$
\H^1_{\tame}(\overline{U},\overline{F})\simeq
\H^1(K,\H^0(\widetilde{U},\widetilde{F})).
$$
Comme $K$ est topologiquement de type fini et
$\H^0(\widetilde{U},\widetilde{F})$ fini, on en d\'eduit tout
d'abord que les fibres du faisceau $\R^1_{\tame}i_*F|_{X_I}$ sont
finies. D'autre part, l'ensemble $\LL$ ne d\'epend pas du point
$\overline{x}$. Pour tout $q\in\LL$, soit $X_{I,q}$ le ferm\'e de
$X_I$ d'\'equation $q=0$ et soit $X_{I'}$ l'ouvert de~$X_I$
compl\'ementaire de la r\'eunion des $X_{I,q}$. Alors
$\R^1_{\tame}i_*F|_{X_{I,q}}$ et $\R^1_{\tame}i_*F|_{X_I'}$ sont
localement constants; en effet une fl\`eche de sp\'ecialisation de
points g\'eom\'etriques de~$X_{I,q}$ (\resp de~$X_{I'}$) induit un
isomorphisme sur les groupes $K$ \eqref{XIII.5.6.1} donc aussi sur les
ensembles $\H^1_{\tame}(\overline{U},\overline{F})$, et l'on peut appliquer
SGA~4 IX~2.13.

\begin{corollaire}
\label{XIII.2.5}
Soient
$f\colon X\to S$ un morphisme, $D$ un diviseur sur~$X$
\`a croisements normaux relativement \`a~$S$ \eqref{XIII.2.1},
$Y=\Supp D$, $U=X- Y$, $i\colon U\to X$ l'immersion
canonique. Soit $\Phi$ un champ sur~$U$ et supposons donn\'es des
morphismes \'etales surjectifs $S_1\to S$ et $X_2\to
X\times_{S}S_1$, un champ $\Psi$ sur~$S_1$ et un isomorphisme
$\Phi|U_2\simeq\Psi|U_2$ (\cf \ref{XIII.2.2}).

Alors, pour tout morphisme $h\colon S'\to S$, si $k=h_{(X)}$,
le foncteur canonique
$$
\varphi\colon k^*i_*\Phi\to i_*'\Phi'
$$
est pleinement fid\`ele. Si $\Psi$ est constructible, le foncteur
canonique
$$
\psi\colon k^*i_*^{\tame}\Phi\to i_*^{\prime\tame}\Phi'
$$
est une \'equivalence de cat\'egories.

De plus, si le champ $\Psi$ est constructible (\resp si $\Psi$ est
$1$-constructible \eqref{XIII.0} et $S$ localement noeth\'erien),
$i_*\Phi$ est constructible (\resp $i_*^t\Phi$ est
$1$-constructible).
\end{corollaire}

Montrons que $\varphi$ est pleinement fid\`ele. Il suffit de voir
que, pour tout point g\'eom\'etrique $\overline{x}'$ de~$X'$, il
en est ainsi du foncteur
$$
\overline{\varphi}\colon
\Phi(\overline{U})\to\Phi'(\overline{U'})
$$
(on a repris les notations de~\ref{XIII.2.4.2}). Soient $a$, $b$ deux
\'el\'ements de~$\Phi(\overline{U})$, $a',\ b'$ leurs images par
$\overline{\varphi}$. Comme le morphisme
$\overline{U}\to\overline{S}$ est localement $0$-acyclique
\eqref{XIII.5.7}, on a un isomorphisme
$$
\H^0(\overline{U},\overline{S\Phi})
\overset{\sim}{\to}\H^0(S,\overline{S\Psi}).
$$
Par suite $a$ et $b$ proviennent par image inverse d'\'el\'ements
de~$\Psi$, et il en est donc de m\^eme de
$F=\SheafHom_{\overline{U}}(a,b)$. Comme le faisceau
$F'=\SheafHom_{\overline{U'}}(a',b')$ est l'image inverse sur~$\overline{U'}$ de~$F$, il r\'esulte de~\ref{XIII.2.4.1} que le
morphisme canonique
$$
\H^0(\overline{U},F)\to \H^0(\overline{U'},F')
$$
est un isomorphisme, ce qui prouve que $\overline{\varphi}$ est
pleinement fid\`ele.

Montrons
que, pour tout point g\'eom\'etrique $\overline{x}'$ de~$X'$, le
foncteur
$$
\overline{\psi}\colon i^{\tame}_*\Phi(\overline{X'})\to
i_*^{\prime\tame}\Phi'(\overline{X'})
$$
est une \'equivalence. D'apr\`es ce qui pr\'ec\`ede,
$\overline{\psi}$ est pleinement fid\`ele. Montrons que
$\overline{\psi}$ est essentiellement surjectif. Soit $a'$ un
\'el\'ement de~$\Phi'(\overline{U'})$ mod\'er\'ement
ramifi\'e sur~$X'$ relativement \`a~$S'$ \eqref{XIII.2.2} et
montrons qu'il est l'image d'un \'el\'ement mod\'er\'ement
ramifi\'e de~$\Phi(\overline{U})$. Il r\'esulte de~\ref{XIII.2.4}
1) que le morphisme canonique
$$
\H^0(\overline{U},\overline{S}\overline{\Phi})\to
\H^0(\overline{U'},\overline{S}\overline{\Phi}')
$$
est un isomorphisme. Soit $G'$ la sous-gerbe maximale de~$\Phi'$
engendr\'ee par $a'$; il existe alors une sous-gerbe maximale $G$ de
$\overline{\Phi}$, image inverse d'une gerbe sur~$\overline{S}$, telle
que l'on ait
$$
\overline{m}^*G\simeq G',
$$
o\`u $m$ est le morphisme $\overline{U'}\to
\overline{U}$. Le foncteur canonique
\begin{equation*}
\label{eq:XIII.2.5.*}
\tag{$*$} {\overline{k}^*\overline{i}_*^{\tame}G\to
\overline{i}_*^{\prime\tame}G'}
\end{equation*}
est une \'equivalence, car $G$ s'identifie \`a une gerbe de
torseurs sous un faisceau en groupes constructible provenant de
$\overline{S}$, et l'on peut appliquer \ref{XIII.2.4}~1). Il
r\'esulte alors de \eqref{eq:XIII.2.5.*} qu'il existe un
\'el\'ement $a$ de~$G(\overline{U})$, mod\'er\'ement
ramifi\'e sur~$X$ relativement \`a~$S$, dont l'image inverse sur~$\overline{U'}$ est $a'$, ce qui prouve que $\psi$ est une
\'equivalence.

Si $\Psi$ est constructible, il en est de m\^eme de~$i_*\Phi$; en
effet un objet $x$ de~$i_*\Phi$ est, localement pour la topologie
\'etale de~$S$, image inverse d'un objet $y$ de~$\Psi$; il
r\'esulte donc de~\ref{eq:XIII.2.4.1.1} que $\SheafAut (x)$ est
l'image inverse de~$\SheafAut (y)$, donc est constructible. Enfin, si
$\Psi$ est $1$-constructible, il en est de m\^eme de
$i^{\tame}_*\Phi$ d'apr\`es~\ref{XIII.6.3} ci-dessous.

\begin{corollaire}
\label{XIII.2.6}
Les
notations sont celles de~\ref{XIII.2.4}. Supposons que $S$ soit de
caract\'eristique nulle en tout point $s$ tel que l'on ait $Y_s\neq
\emptyset.$ Alors, si $\cal{ F}$ est un faisceau en groupes localement
constant constructible sur~$U$ (\resp un champ constructible sur~$U$ qui
est localement sur~$X$ et $S$ image inverse d'un champ constructible
sur~$S$), $(\cal{ F},i)$ est cohomologiquement propre relativement
\`a~$S$ en dimension $\leq 1.$
\end{corollaire}

Comme tout faisceau d'ensembles constructible sur~$U$ est
mod\'er\'ement ramifi\'e sur~$X$ relativement \`a~$S,$ le
corollaire r\'esulte de~\ref{XIII.2.4} (\resp \ref{XIII.2.5}).

\begin{corollaire}
\label{XIII.2.7}
Les notations sont celles de~\ref{XIII.2.4}, mais on se donne de plus
un $S$-sch\'ema $T$, un morphisme propre $p\colon X\to T,$
et on suppose $X$ et $T$ de pr\'esentation finie sur~$S$; soit
$q=pi$. Soit $\cal{ F}$ un faisceau d'ensembles constructible sur~$U$
satisfaisant \`a l'une des conditions \textup{a)} ou \textup{b)} de~\ref{XIII.2.4}
(\resp un faisceau en groupes satisfaisant \`a l'une des condition
\textup{a)}, \textup{b)} de~\ref{XIII.2.4}, \resp un champ sur~$U$ qui est localement sur~$X$ et $S$ image inverse d'un champ constructible $G$ sur~$S$). Alors
on a les conclusions suivantes:
\begin{enumerate}
\item[1)] $(\cal{ F},q)$ est cohomologiquement propre relativement \`a~$S$ en dimension $\leq 0$ (\resp pour tout morphisme $h\colon
S^{'}\to S$, si $m=h_{(T)}$, le morphisme canonique:
$$
\Theta\colon m^*(\R^1_{\tame}q_*\cal{ F})\to
\R^1_{\tame}q_*^{'}\cal{ F}^{'}
$$
est un isomorphisme, \resp pour tout morphisme $S^{'}\to S$,
le morphisme canonique
$$
\xi \colon m^*(q_*^{\tame}\cal{ F})\to
{q'}_*^{{\tame}}\cal{F}^{'}
$$
est une \'equivalence).
\item[2)] le faisceau $q_*\cal{ F}$ (\resp le faisceau
$\R_{\tame}^1q_*\cal{ F}$, \resp le champ $q_*^{\tame}\cal{ F}$)
est constructible. Dans le dernier cas, si l'on suppose $S$ localement
noeth\'erien et $G$ $1$-constructible, il en est de m\^eme de
$q_*^{\tame}\cal{ F}.$
\end{enumerate}
\end{corollaire}

La
premi\`ere partie r\'esulte aussit\^ot de~\ref{XIII.2.4},
\ref{XIII.2.5} et de la d\'emonstration
de~\ref{XIII.1.8}. D\'emontrons~2). Si $\cal{ F}$ est un faisceau
d'ensembles constructible sur~$U$ satisfaisant \`a \ref{XIII.2.4}~a)
ou \ref{XIII.2.4}~b), il r\'esulte de \ref{XIII.2.4}~2) que
$i_*\cal{ F}$ est constructible; il en est donc de m\^eme de
$q_*\cal{ F}=p_*(i_*\cal{ F})$ (SGA~4 XIV~1.1).

Soit $\cal{F}$ un faisceau en groupes constructible sur~$U$
satisfaisant \`a \ref{XIII.2.4}~a) ou \ref{XIII.2.4}~b) et prouvons
que $\R_{\tame}^1q_*\cal{ F}$ est constructible. Par passage \`a
la limite (EGA~IV 8.10.5 et 17.7.8) et en utilisant 1), on peut
supposer $S$ noeth\'erien. Soit alors $\Phi$ le champ sur~$X$ dont
la fibre en tout sch\'ema $X^{'}$ \'etale sur~$X$ est form\'ee
des torseurs sur~$U^{'}=U\times_{X}X^{'}$, de groupe $\cal{F}|U$, qui
sont mod\'er\'ement ramifi\'es sur~$X$ relativement \`a~$S$; on
a donc
$$
S(i_*^{\tame}\Phi)\simeq \R_{\tame}^1i_*\cal{ F}
$$
et ce faisceau est constructible d'apr\`es \ref{XIII.2.4}~2).
Il r\'esulte donc de~\ref{XIII.6.3} ci-dessous que $S(p_*\Phi)$
est constructible, \ie que $\R^1_{\tame}q_*\cal{ F}$ est
constructible.

Enfin, si $\cal{ F}$ est un champ sur~$U$ qui est localement sur~$X$ et
$S$ image inverse d'un champ constructible sur~$S,$
$i_*^{\tame}\cal{ F}$ est constructible, et il en est donc de
m\^eme de~$q_*^{\tame}\cal{ F}=p_*(i_*^{\tame}\cal{ F})$. Si
de plus $S$ est localement noeth\'erien et $SG$ constructible,
$S(i_*^{\tame}\cal{ F})$ est constructible d'apr\`es 6.3; il en
est donc de m\^eme de~$S(q_*i_*^{\tame}\cal{ F})$, \ie de
$S(q_*^{\tame}\cal{ F})$~\ref{XIII.6.2}.

\begin{corollaire}
\label{XIII.2.8}
Soit
$$
\xymatrix{{U}\ar[r]^{i}\ar[d]_{g}&{X}\ar[ld]^{f}\\{S}}
$$
un diagramme commutatif de sch\'emas, dans lequel $U$ est l'ouvert
compl\'ementaire
dans $X$ d'un diviseur \`a croisements normaux relativement \`a~$S$, $f$ un morphisme \emph {propre} de pr\'esentation finie. Soit
$\LL$ un ensemble de nombres premiers. On suppose~$g$ localement
$0$-acyclique (\resp localement $1$-asph\'erique pour $\LL$). Alors, si
$\cal{ F}$ est un faisceau d'ensembles sur~$U$ (\resp un faisceau de
$\LL$-groupes) sur~$U,$ localement constant constructible,
mod\'er\'ement ramifi\'e sur~$X$ relativement \`a~$S$,
$f_*\cal{ F}$ (\resp $\R^1_{\tame}f_*\cal{ F}$) est localement
constant constructible et $(\cal{ F},f)$ est cohomologiquement propre
relativement \`a~$S$ en dimension $\leq 0$ (\resp la formation de
$\R^1_{\tame}f_*\cal{ F}$ commute \`a tout changement de base
$S^{'}\to S$). Dans le cas non resp\'e, si $\cal{ F}$ est un
faisceau en groupes, pour toute sp\'ecialisation $\overline{s}_1
\to \overline{s}_2$ de points g\'eom\'etriques de~$S$, le
morphisme de sp\'ecialisation
$$
(\R^1_{\tame}f_*\cal{ F})_{\overline{s}_2}\to
(\R^1_{\tame}f_*\cal{ F})_{\overline{s}_1}
$$
est injectif.
\end{corollaire}
Le corollaire r\'esulte aussit\^ot de~\ref{XIII.2.6} et
de~\ref{XIII.1.14} (\resp de l'analogue de~\ref{XIII.1.14} pour le
$\R^1_{\tame}f_*\cal{ F}$, lequel se d\'emontre comme \loccit).

\begin{corollaire} \label{XIII.2.9}
Soit
$$
\xymatrix{{U}\ar[r]^{i}\ar[d]_{g}&{X}\ar[ld]^{f}\\{S}}
$$
un diagramme commutatif de sch\'emas, dans lequel $U$ est l'ouvert
compl\'ementaire dans~$X$ d'un diviseur \`a croisements normaux
relativement \`a~$S$, $f$~un morphisme propre lisse de
pr\'esentation finie. Soit~$\LL $ l'ensemble des nombres premiers
distincts des caract\'eristiques r\'esiduelles de~$S$. Soit~$\cal{F}$ un faisceau de~$\LL $-groupes localement constant constructible
sur~$U$, mod\'er\'ement ramifi\'e sur~$X$ relativement
\`a~$S$. Alors $\R^1 f_*\cal{ F}$ est localement constant
constructible et $(\cal{ F},f)$ est cohomologiquement propre
relativement
\`a~$S$ en dimension $\leq 1$.
\end{corollaire}
Le corollaire r\'esulte de~\ref{XIII.2.8} et du fait que l'on a
$\R_{\tame}^1f_*\cal{ F}=\R^1f_*\cal{ F}$ (\ref{XIII.2.3}~b)).

\subsection{}
\label{XIII.2.10}
Si $U$ est un sch\'ema connexe, $a$ un point g\'eom\'etrique
de~$U$, $\LL$~un ensemble de nombre premiers, on note
\begin{equation*}
\label{eq:XIII.2.10.0}
\tag{2.10.0} \pi_1^{\LL}(U,a)
\end{equation*}
\label{indnot:mi}la limite projective des quotients finis de~$\pi_1(U,a)$ dont les
ordres ont tous leurs facteurs premiers dans~$\LL$.

Nous allons d\'efinir des morphismes de sp\'ecialisation pour le
groupe fondamental, g\'e\-n\'e\-ra\-li\-sant X.\ref{X.2}.

Soit $g\colon U\to S$ un morphisme coh\'erent \`a fibres
g\'eom\'etriquement connexes (\resp un morphisme de la forme $g=
fi$, o\`u $f\colon X \to S$ est un morphisme propre de
pr\'esentation finie et o\`u $i\colon U\to X$ est une
immersion ouverte telle que $U$ soit le compl\'ementaire dans $X$
d'un diviseur \`a croisements normaux relativement \`a~$S$
(\cf \ref{XIII.2.8})). Soit ${\LL}$ un ensemble de nombres premiers et
supposons, dans le cas non resp\'e, que, pour tout ${\LL}$-groupe
constant fini $C$, $(C_U,g)$ soit cohomologiquement propre
relativement \`a~$S$ en dimension $\leq 1$. Soient $\overline{s}_1
\to \overline{s}_2$ un morphisme de sp\'ecialisation de
points g\'eom\'etriques de~$S$, $\overline{S}$ le localis\'e
strict de~$S$ en $\overline{s}_2$,
$\overline{U}=U\times_S\overline{S}$. On a un diagramme commutatif
$$
\xymatrix{{U_{\bar{s}_1}}\ar[r]^{h_1}\ar[d]&{\bar{U}}\ar[d]^{\bar{g}}&{U_{\bar{s}_1}}\ar[l]_{h_2}\ar[d]\\{\bar{s}_1}\ar[r]&{\bar{S}}&{\,\bar{s}_2\,.}\ar[l]}
$$
Si $a_1$ est un point g\'eom\'etrique de~$U_{\bar{s}_1},\ a_2$ un
point g\'eom\'etrique de~$U_{\bar{s}_2}$ les morphismes $h_1$ et
$h_2$ d\'efinissent des morphismes canoniques
\begin{align*}
\pi_1\colon \pi_1^{\LL}(U_{\bar{s}_1},a_1)&\to
\pi_1^{\LL}(\bar{U},a_1) & \pi_2\colon
\pi_1^{\LL}(U_{\bar{s}_2},a_2)&\to \pi_1^{\LL}(\bar{U},a_2)\\
(\text{\resp~} \pi_1\colon \pi_1^\tame (U_{\bar{s}_1},a_1)&\to
\pi_1^\tame (\bar{U},a_1) & \pi_2\colon
\pi_1^\tame(U_{\bar{s}_2},a_2)&\to \pi_1^\tame(\bar{U},a_2))
\end{align*}
(V~\ref{V.7} et~\ref{eq:XIII.2.1.5.2}). Les hypoth\`eses de
propret\'e cohomologique (\resp \ref{XIII.2.8}) prouvent que $\pi_2$
est un isomorphisme. Si l'on choisit une classe de chemins de~$a_1$
\`a $a_2$, on obtient un isomorphisme
$$
\pi_{12}\colon \pi_1^{\LL}(\bar{U},a_1)\isomto
\pi_1^{\LL}(\bar{U},a_2) \quad (\text{\resp~} \pi_{12}\colon
\pi_1^\tame (\bar{U},a_1)\isomto \pi_1^\tame(\bar{U},a_2))
$$
d'o\`u un morphisme $\pi=\pi_2^{-1}\pi_{12}\pi_1$
$$
\pi\colon \pi_1^{\LL}(U_{\bar{s}_1},a_1)\to
\pi_1^{\LL}(U_{\bar{s}_2},a_2) \quad (\text{\resp~}\pi\colon
\pi_1^\tame (U_{\bar{s}_1},a_1)\to \pi_1^\tame
(U_{\bar{s}_2},a_2)).
$$
Changer la classe de chemins de~$a_1$ \`a $a_2$ revient \`a
modifier $\pi$ par un automorphisme int\'erieur de
$\pi_1^{\LL}(X_{\bar{s}_2},a_2)$ (\resp de
$\pi_1^\tame(X_{\bar{s}_2},a_2)$). On appelle \emph{morphisme de
spécialisation pour le groupe fondamental}
\index{homomorphisme de sp\'ecialisation pour le groupe fondamental|hyperpage}\index{sp\'ecialisation pour le groupe fondamental (homomorphisme de)|hyperpage}associ\'e au morphisme $\bar{s}_1\to \bar{s}_2$ et on note
simplement
$$
\pi\colon \pi_1^{\LL}(X_{\bar{s}_1})\to
\pi_1^{\LL}(X_{\bar{s}_2}) \quad \text{(\resp~} \pi\colon
\pi_1^{\tame}(X_{\bar{s}_1})\to \pi_1^{\tame}(X_{\bar{s}_2}))
$$
l'un des morphismes d\'efini ci-dessus.

\begin{lemme}
\label{XIII.2.11}
Soient $f\colon X\to S$ un morphisme propre de
pr\'esentation finie, $D$ un diviseur sur~$X$ \`a croisements
normaux relativement \`a~$S$, $Y=\Supp D$, $U=X-Y$, $i\colon
U\to X$ le morphisme canonique, $\bar{s}_1\to
\bar{s}_2$ un morphisme de sp\'ecialisation de points
g\'eom\'etriques de~$S$, $y_1$ un point g\'eom\'etrique de
$Y_{\bar{s}_1}$, $y_2$ un point g\'eom\'etrique de
$Y_{\bar{s}_2}$, tel que la projection $z_1$ de~$y_1$ sur~$X$ soit une
g\'en\'erisation de la projection $z_2$ de~$y_2$. Soit $I_{y_1}$
un sous-groupe d'inertie de~$\pi_1^{\tame}(\bar{U}_{\bar{s}_1})$ en
$y_1.$ Alors l'image de~$I_{y_{1}}$ par le morphisme de
sp\'ecialisation
$$
\pi\colon \pi_1^{\tame}(U_{\bar{s}_1})\to
\pi_1^{\tame}(U_{\bar{s}_2})
$$
est
un sous-groupe d'inertie de~$\pi_1^{\tame}(U_{\overline{s}_2})$
en~$y_2$.
\end{lemme}

Soient en effet $\overline{X}$ (\resp $\widetilde{X}$) le localis\'e
strict de~$X$ en $y_2$ (\resp en $y_1$), $\overline{U}=U \times_X
\overline{X}$ (\resp $\widetilde{U}=U\times_X \widetilde{X}$). On a un
morphisme canonique $\widetilde{U}\to \overline{U}$, et il
r\'esulte de~\ref{XIII.1.10} que l'on a un diagramme commutatif
$$
\xymatrix{{\pi_1^{\tame} \left(\widetilde{U}_{\overline{s}_1} \right)}\ar[r]^{\pi '}\ar[d]&{\pi_1^{\tame} \left(
\overline{U}_{\overline{s}_2} \right)}\ar[d]\\{\pi_1^{\tame} \left(
U_{\overline{s}_1} \right)}\ar[r]^{\pi }&{\pi_1^{\tame} \left(
U_{\overline{s}_2} \right)}}
$$
o\`u $\pi'$ est compos\'e du morphisme canonique
$\pi_1^{\tame}(\widetilde{U}_ {\overline{s}_1})\to
\pi_1^{\tame}(\overline{U}_{\overline{s}_1})$ et du morphisme de
sp\'ecialisation. Comme
$\pi_1^{\tame}(\widetilde{U}_{\overline{s}_1})$
(\resp $\pi_1^{\tame}(\overline{U}_{\overline{s}_1})$)
est un groupe d'inertie de
$\pi_1^{\tame}(U_{\overline{s}_1})$ en~$y_1$ et (\resp de~$\pi_1^{\tame}(U_{\overline{s}_2})$ en~$y_2$), il suffit de
prouver que $\pi'$ est surjectif. Mais cela r\'esulte de
l'expression obtenue dans~\ref{XIII.5.6}.

\begin{corollaire}
\label{XIII.2.12}
Soit $X$ une courbe propre et lisse connexe de genre $g$ sur un corps
s\'e\-pa\-ra\-blement clos $k$ de caract\'eristique $p\geq
0$. Soit $U$ l'ouvert obtenu en enlevant \`a $X$ $n$ points
ferm\'es distincts $a_1,\dots,a_n$. Alors le groupe fondamental
mod\'er\'ement ramifi\'e $\pi_1^{\tame}(U)$ \eqref{XIII.2.1.3}
peut \^etre engendr\'e par $2g+n$ \'el\'ements $x_i,
y_i,\sigma_j$, avec $1\leq i\leq g$, $1\leq j\leq n$, tel que
$\sigma_j$ soit un g\'en\'erateur d'un groupe d'inertie
correspondant \`a $a_j$, et que l'on ait la relation
\begin{equation*}
\label{eq:XIII.2.12.*}
\tag{$*$} \prod_{1\leq i\leq
g}(x_iy_ix^{-1}_iy^{-1}_i)\cdot\prod_{1\leq j\leq n}\sigma_j=1.
\end{equation*}
Pour tout groupe fini $G$ \emph{d'ordre premier à} $p$,
engendr\'e par des \'el\'ements $\overline{x}_i,
\overline{y}_i,\overline{\sigma}_j$ satisfaisant \`a la relation
\eqref{eq:XIII.2.12.*}, il existe un rev\^etement \'etale de~$U$,
de groupe $G$, correspondant \`a un homomorphisme
$\pi_1^{\tame}(U)\to G$ qui
envoie $x_i, y_i,\sigma_j$ sur~$\overline{x}_i,
\overline{y}_i,\overline{\sigma}_j$ respectivement. En d'autres
termes, si $p'$ d\'esigne l'ensemble des nombres premiers distincts
de~$p$, $\pi_1^{p'}(U)$ est le pro-$p'$-groupe engendr\'e par les
g\'en\'erateurs $x_i, y_i,\sigma_j$ li\'es par la seule relation~\eqref{eq:XIII.2.12.*}.
\end{corollaire}

\subsubsection*{D\'emonstration}
On peut supposer $k$ alg\'ebriquement clos. Supposons d'abord $k$ de
carac\-t\'eristique z\'ero. Il existe alors une sous-extension
alg\'ebriquement close $k'$ de~$k$, de degr\'e de transcendance
fini sur~$\QQ$, telle que $X$ provienne par extension des scalaires
d'une courbe propre et lisse $X'$ d\'efinie sur~$k'$, et l'on peut
supposer que les points $a_1,\dots,a_n$ proviennent de points rationnels
$a'_1,\dots,a'_n$ de~$X'$. Comme $k'$ est de degr\'e de transcendance
fini sur~$\QQ$, on peut trouver un plongement de~$k'$ dans le corps
des nombres complexes $\CC$; soit $\widetilde{U}=U'\times_{k'}\CC.$
Soit $k''$ une extension alg\'ebriquement close de~$k'$ telle que
l'on ait des $k'$-morphismes de~$k$ et de~$\CC$ dans $k''$. Si
$g'\colon U'\to k'$ est le morphisme structural, et si
$\cal{F}$ est un faisceau en groupes constant fini sur~$U''$, il
r\'esulte de~\ref{XIII.2.9} que les morphismes de sp\'ecialisation
$$
(\R^1g'_*\cal{F}')_{\CC}\to(\R^1g'_*\cal{F}')_{k''}
\from(\R^1g'_*\cal{F}')_k
$$
sont des isomorphismes. En termes de groupes fondamentaux, cela montre
que l'on a un isomorphisme, d\'efini \`a automorphisme
int\'erieur pr\`es
$$
\pi_1(U)\to\pi_1(\widetilde{U}),
$$
et il est clair que cet isomorphisme transforme un groupe d'inertie
relatif \`a un point de~$X'-U'$ en un groupe d'inertie relatif au
m\^eme point. On peut donc supposer que l'on a $k=\CC$. Dans ce
dernier cas il r\'esulte du th\'eor\`eme d'existence de Riemann
(XII~\ref{XII.5.2}) que le groupe fondamental $\pi_1(U)$ n'est autre
que le compl\'et\'e pour la topologie des sous-groupes d'indice
fini du groupe fondamental de l'espace analytique associ\'e \`a~$U$. Or ce dernier peut se calculer par voie transcendante \cite[ch.\, 7
\S 47]{XIII.3};
il peut \^etre engendr\'e par $2g+n$ \'el\'ements $x_i,
y_i,\sigma_j$ tels que $\sigma_j$ soit l'image d'un g\'en\'erateur
du groupe fondamental local $\pi_1(D_j)$ d'un petit disque centr\'e
en $a_j$, \ie un g\'en\'erateur d'un groupe d'inertie
correspondant au point $a_j$, ces \'el\'ements satisfaisant \`a
la seule relation \eqref{eq:XIII.2.12.*}.

Si maintenant $k$ est de
caract\'eristique $p>0$, on peut trouver un anneau de valuation
discr\`ete complet $A$, de corps r\'esiduel $k$, de corps des
fractions $K$ de caract\'eristique z\'ero, et un sch\'ema
connexe $X_1$, propre et lisse sur~$S=\Spec A$, tel que l'on ait
$X_1\times_S \Spec k\simeq X$ (III~\ref{III.7.4}). Les points $a_j$ se
rel\`event alors en des sections $s_j$ de~$X_1$ au-dessus de~$S$;
soit $Y_{1j}$ le sous-sch\'ema ferm\'e r\'eduit d'espace
sous-jacent $s_j(S)$, $Y_1$ la r\'eunion des $Y_{1j}$,
$U_1=X_1-Y_1$, $g_1\colon U_1\to S$ le morphisme
structural. Soient $\overline{K}$ une extension alg\'ebriquement
close de~$K$, $\overline{U}= U_1\times_S\overline{K}$. Si $C$ est un
groupe constant fini, il r\'esulte de~\ref{XIII.2.8} que le
morphisme de sp\'ecialisation
$$
(\R^1g_{1*}C_{U_1})_k\to(\R^1g_{1*}C_{U_1})_{\overline{K}}
$$
est injectif et m\^eme bijectif si $C$ est d'ordre premier \`a
$p$. Or cela signifie, en termes de groupes fondamentaux, que le
morphisme de sp\'ecialisation \eqref{XIII.1.10}
$$
\pi \colon \pi_1(\overline{U})\to\pi^{\tame}_1(U)
$$
est surjectif, et que le morphisme de sp\'ecialisation
$$
\pi^{p'}_1(\overline{U})\to\pi^{p'}_1(U)
$$
est bijectif. Enfin, si $x_i, y_i,\sigma_j$ sont des
g\'en\'erateurs de~$\pi_1(\overline{U})$ tels que $\sigma_j$ soit
un g\'en\'erateur d'un groupe d'inertie correspondant au point
$b_j=Y_{1j}(\overline{K})$ de~$\overline{X}$, alors,
d'apr\`es~\ref{XIII.1.11}, $\pi(\sigma_j)$ est un g\'en\'erateur
d'un groupe d'inertie correspondant \`a $a_j$, ce qui ach\`eve la
d\'emonstration.


\begin{remarqueMR}
\label{XIII.2.13}
Soient $k$ un corps alg\'ebriquement clos de caract\'eristique $p>0$, $X$~une courbe alg\'ebrique propre lisse sur~$k$, connexe, de genre $g$, $U$~un ouvert affine de~$X$, compl\'ementaire de~$r\geq1$ points rationnels
de~$X$. On dispose du groupe fondamental $\pi_1(U)$, de son quotient
$\pi_1^{\tame}(U)$ qui classifie les rev\^etements finis \'etales de~$U$,
mod\'er\'ement ramifi\'es aux points de~$X-U$, et du quotient
$\pi_1^{p'}(U)$ de~$\pi_1^{\tame}(U)$ qui classifie les rev\^etements \'etales
galoisiens de~$U$, de groupe de Galois d'ordre premier \`a~$p$.

Dans (Coverings of algebraic curves, Amer.\ J.~Math.\ \textbf{79}
(1957), p.\,825--856) S.\, Abhyankar formulait un certain nombre de conjectures sur la structure des groupes finis, groupes de Galois de rev\^etements finis
\'etales de~$U$.

Les conjectures concernant les groupes finis qui sont groupes de
Galois de rev\^etements \'etales connexes de~$U$ d'ordre premier \`a~$p$,
sont d\'emontr\'ees et pr\'ecis\'ees dans le corollaire \ref{XIII.2.12}. Avant d'aborder les quotients finis de~$\pi_1(U)$, commen\c{c}ons par donner quelques
indications sur la \og taille\fg de ce groupe. Soit $X = \Spec(A)$. On
sait que $\Hom_{\mathrm{cont}}(\pi_1(U), \ZZ/p\ZZ)$ est d\'ecrit par
la th\'eorie d'Artin-Schreier (cor.\, XI~\ref{XI.6.9}). Ce groupe est
isomorphe \`a $A/\wp(A)$ o\`u $\wp$ est l'application de~$A$ dans $A$, qui
envoie $a$ sur~$a^p-a$. Supposons d'abord que $U$ soit la droite
affine $\AA^1$, d'anneau $k[T]$. Soit~$E$ l'ensemble des
\'el\'ements de~$k[T]$ de la forme $\sum_i a_iT^i$, avec $a_i = 0$
lorsque $p$ divise~$i$. Il est imm\'ediat que l'application compos\'ee
$E\to k[T]\to k[T]/\wp k[T]$ est bijective. D\`es lors les coefficients
$a_i$, $(i,p)=1$, se comportent comme des coordonn\'ees d'un espace qui
param\`etre les rev\^etements de la droite affine cycliques de degr\'e
$p$. En particulier, on en d\'eduit que $\pi_1(\AA^1)$ n'est pas
topologiquement de type fini et que, si l'on effectue une extension
$k\to k'$ de corps alg\'ebriquement clos, le morphisme naturel $\pi_1(\AA^1\times_kk')\to\pi_1(\AA^1)$, qui est surjectif, n'est pas bijectif,
contrairement au cas propre. Dans le cas g\'en\'eral, la courbe $U$ se
r\'ealise comme sch\'ema fini sur la droite affine et on conclut que les
m\^emes ph\'enom\`enes se produisent pour $\pi_1(U)$ (et d'ailleurs, plus
g\'en\'eralement, pour $\pi_1(V)$ pour tout $k$-sch\'ema affine connexe $V$,
de type fini, de dimension $\geq1$).

Ceci \'etant, si $G$ est un groupe fini, notons $G^{(p')}$ le plus grand
groupe quotient de~$G$, d'ordre premier \`a~$p$. Pour qu'un groupe fini
$G$ soit un quotient topologique de~$\pi_1(U)$, il faut que $G^{(p')}$
soit un quotient topologique de~$\pi_1^{p'}(U)$, condition \`a laquelle
on sait en principe r\'epondre gr\^ace au corollaire \ref{XIII.2.12}. Dans l'article
pr\'ecit\'e, Abhyankar conjecture que cette condition n\'ecessaire est
\'egalement suffisante. Cette conjecture a \'et\'e d\'emontr\'ee par M.\, Raynaud
dans le cas de la droite affine et par D.\, Harbater dans le cas g\'en\'eral
(Invent. Math.\ \textbf{116} (1994), p.\,425--462 et \textbf{117}, p.\,1--25). Par exemple,
dans le cas de la droite affine $\pi_1^{p'}(\AA^1)=1$ et on
conclut qu'un groupe fini $G$ est groupe de Galois d'un rev\^etement
\'etale connexe de~$\AA^1$ si et seulement si $G^{(p')} = 1$,
c'est-\`a-dire si et seulement si $G$ est engendr\'e par ses
$p$-sous-groupes de Sylow. Ainsi tout groupe fini simple d'ordre
multiple de~$p$ convient.
\end{remarqueMR}


\section{Propret\'e cohomologique et locale acyclicit\'e g\'en\'erique}
\label{XIII.3}

\begin{theoreme}
\label{XIII.3.1}
Soient $S$ un sch\'ema irr\'eductible de point
g\'en\'erique~$s$, $X$ et $Y$ deux $S$-sch\'emas de
pr\'esentation finie, $f\colon X \to Y$ un $S$-morphisme. Pour tout
$S$-sch\'ema~$S'$, on note $Y'$, $X'$, etc. l'image inverse de~$Y$,
$X$, etc. par le morphisme $S' \to S$. On a les propri\'et\'es
suivantes:

\begin{enumerate}
\item[1)]
{\upshape a)} On peut trouver un ouvert non vide~$S'$ de~$S$ tel que, pour
tout faisceau d'ensembles constant fini $F'$ sur~$X'$, $f'_*F'$ soit
constructible, et que $(F',f')$ soit cohomologiquement propre
relativement \`a~$S'$ en dimension $\le 0$.

{\upshape b)} Soit $F$ un faisceau d'ensembles constructible sur~$X$. Alors
on peut trouver un ouvert non vide~$S'$ de~$S$ (d\'ependant de~$F$)
tel que $f'_*F'$ soit constructible et que $(F',f')$ soit
cohomologiquement propre relativement \`a~$S'$ en dimension $\le 0$.

\item[2)] Supposons que les sch\'emas de type fini de dimension $\le
\dim X_s$ sur une cl\^oture alg\'ebrique $\overline{k}$ de~$\kres(s)$
soient fortement d\'esingularisables \textup{(SGA~5 I~3.1.5)}. Alors on a de
plus les propri\'et\'es suivantes:

{\upshape a)} On peut trouver un ouvert non vide~$S'$ de~$S$ tel que, pour
tout faisceau en groupes constant fini $F'$ sur~$X'$, d'ordre premier
aux caract\'eristiques r\'esiduelles de~$S$, si $\Phi'$ est le
champ des torseurs sous $F'$, $f'_*\Phi'$ soit $1$-constructible, et
que $(F',f')$ soit cohomologiquement propre relativement \`a~$S'$ en
dimension $\le 1$.

{\upshape b)} Soient $\LL$ l'ensemble des nombres premiers distincts des
caract\'eristiques r\'esiduelles de~$S$ et $\Phi$ un
ind-$\LL$-champ $1$-constructible sur~$X$ \eqref{XIII.0}, tel que,
pour tout sch\'ema~$X_1$ \'etale sur~$X$ et pour tout couple
d'objets $x, x_1$ de~$\Phi_{X_1}$, le faisceau $\SheafHom
_{X_1}(x,x_1)$ soit constructible. On suppose de plus $S$ localement
noeth\'erien. Alors on peut trouver un ouvert non vide~$S'$ de~$S$
tel que $f'_*\Phi'$ soit $1$-constructible, que, pour tout couple
d'objets $y, y_1$ d'une fibre $(f'_*\Phi')_{Y_1}$,
$\SheafHom_{Y_1}(y,y_1)$ soit constructible,
et que $(\Phi',f')$ soit cohomologiquement propre relativement \`a~$S'$ en dimension $\le 1$.
\end{enumerate}
\end{theoreme}

\subsubsection*{D\'emonstration}
On peut suppose $S$ affine; d'apr\`es SGA~4 VIII~1.1, on peut
supposer $S$ int\`egre; enfin, par passage \`a la limite, on peut
supposer que $S$ est le spectre d'une alg\`ebre de type fini sur~$\ZZ$; en particulier $S$ est alors noeth\'erien. Comme la question
est locale sur~$Y$, on peut supposer $Y$ affine. De plus il suffit,
pour d\'emontrer le th\'eor\`eme, de le faire apr\`es
extension finie $S'\to S$, o\`u $S'$ est un sch\'ema int\`egre
et o\`u $S'\to S$ est compos\'e de morphismes \'etales et de
morphismes finis radiciels surjectifs.

\subsection*{1) Cas des faisceaux d'ensembles constants}

1) 1. R\'eduction au cas o\`u $X$ est normal sur~$S$. Soit
$X_{1\overline{s}}$ le normalis\'e de~$(X_{\overline{s}})_{\red}$;
quitte \`a restreindre $S$ \`a un ouvert non vide et \`a faire
une extension radicielle de~$S$, on peut supposer que
$X_{1\overline{s}}$ provient d'un sch\'ema $X_1$ normal sur~$S$, et
que le morphisme $X_{1\overline{s}}\to X_{\overline{s}}$ provient d'un
morphisme fini surjectif $p\colon X_1\to X$ (EGA~IV 8.8.2 et
9.6.1). Supposons le th\'eor\`eme d\'emontr\'e pour
$fp$. Quitte \`a restreindre $S$ \`a un ouvert, on peut supposer
que, pour tout faisceau d'ensembles constant $F$ sur~$X$, $(p^*F,fp)$
est cohomologiquement propre relativement \`a~$S$ en dimension $\le
0$ et que $f_*p_*(p^*F)$ est
constructible. D'apr\`es~\ref{XIII.1.9}, $(p_*p^*F,f)$ est alors
cohomologiquement propre relativement \`a~$S$ en dimension $\le
0$. Le morphisme
$$
F \to p_* p^* F=G
$$
est un monomorphisme. Il en r\'esulte d\'ej\`a, $f_*F$ \'etant
un sous-faisceau de~$f_*G$, que $f_*F$ est constructible (SGA~4
IX~2.9~(ii)) et que $(F,f)$ est cohomologiquement propre relativement
\`a~$S$ en dimension $\le -1$.

Soient $X_2=X_1 \times_{X} X_1$, $q \colon X_2 \to X$ le morphisme
canonique. D'apr\`es \ref{XIII.1.11}~1),
on a une suite exacte
$$
\xymatrix@C=.5cm{{F}\ar[r]&{G}\ar@<2pt>[r]\ar@<-2pt>[r]&{q_*q^*F}}.
$$

D'apr\`es ce que l'on vient de d\'emontrer, appliqu\'e \`a
$fq$ au lieu de~$f$, on peut supposer, quitte \`a restreindre $S$
\`a un ouvert non vide, que, pour tout faisceau d'ensembles constant
$F$ sur~$X$, $(q^*F,fq)$ est cohomologiquement propre relativement
\`a~$S$ en dimension $\le -1$, donc que $(q_*q^*F,f)$ est
cohomologiquement propre relativement \`a~$S$ en dimension $\le
-1$. Il r\'esulte alors de \ref{XIII.1.13}~1) que $(F,f)$ est
cohomologiquement propre relativement \`a~$S$ en dimension $\le 0$.

1) 2. R\'eduction au cas o\`u $X$ est normal affine sur~$S$.

Soit $U_s$ un ouvert affine de~$X_s$ dense dans $X_s$. Quitte \`a
restreindre $S$ \`a un ouvert non vide, on peut supposer que $U_s
\to X_s$ se rel\`eve en une immersion ouverte $i\colon U\to X$,
sch\'ematiquement dominante relativement \`a~$S$ (EGA~IV
8.9.1). Comme le morphisme $X \to S$ est normal, on a d'apr\`es
SGA~2 XIV~1.18:
$$
\prof \et_{S- U} (X) \ge 2 \quoi ;
$$
par suite, pour tout faisceau constant $F$ sur~$X$, le morphisme
canonique
$$
F \to i_*i^*F
$$
est un isomorphisme. Il en r\'esulte que, si l'on suppose le
th\'eor\`eme d\'emontr\'e pour $fi$ et $i$, alors, apr\`es
restriction de~$S$ \`a un ouvert non vide, $(i^*F,i)$ et $(i^*F,fi)$
sont cohomologiquement propres relativement \`a~$S$ en dimension
$\le 0$. Il en est donc de m\^eme de~$(F,f)$ (\ref{XIII.1.6}~2)). Comme de plus $f_* F=(fi)_*(i^*F)$ est constructible, ceci
ach\`eve la r\'eduction.

1) 3. Fin de la d\'emonstration.

On peut supposer $S$ normal (EGA~IV 7.8.3). On peut trouver une
compactification de~$X_s$:
$$
\xymatrix{{X_s}\ar[r]^{j_s}\ar[d]_{f_s}&{P_s}\ar[ld]^{g_s}\\{Y_s}}
$$
o\`u $j_s$ est une immersion ouverte dominante et $g_s$ un morphisme
propre; quitte \`a faire une extension radicielle de~$\kres(s)$ et \`a
remplacer $P_s$ par son normalis\'e, ce qui ne change pas $X_s$, on
peut supposer $P_s$ g\'eom\'etriquement normal. Quitte \`a
restreindre $S$ \`a un ouvert non vide et \`a faire une extension
radicielle surjective, on peut supposer que le diagramme ci-dessus
provient d'un diagramme
$$
\xymatrix{{X}\ar[r]^{j}\ar[d]_{f}&{P}\ar[ld]^{g}\\{Y}}
$$
o\`u $P$ est un sch\'ema normal sur~$S$, $j$ une immersion ouverte
sch\'ematiquement dominante relativement \`a~$S$ et $g$ un
morphisme propre (EGA~IV 6.9.1, 9.9.4 et 9.6.1). Pour tout faisceau
d'ensembles constant fini $F$ sur~$X$ de valeur $C$, $j_*F$ est le
faisceau constant de valeur $C$ (SGA~4 2.14.1), et il en est de
m\^eme apr\`es tout changement de base $S' \to S$; il en
r\'esulte que $(F,j)$ est cohomologiquement propre relativement
\`a~$S$ en dimension $\le 0$. Il en est donc de m\^eme de~$(F,f)$,
puisque $g$ est propre \eqref{XIII.1.8}. Comme $g$ est propre
$f_*F=g_*C_P$ est constructible, ce qui ach\`eve la
d\'emonstration de 1) a).

\subsection*{2) Cas d'un faisceau d'ensembles constructible}

Soit $F$ un faisceau d'ensembles constructible sur~$X$. D'apr\`es
SGA~4 IX 2.14~(ii), on peut trouver une famille finie de morphismes
$p_i \colon Z_i \to X$, et, sur chaque $Z_i$, un faisceau d'ensembles
constant fini $C_i$, de sorte que l'on ait un monomorphisme
$$
j \colon F \to \prod_{i} {p_i}_* C_i =G \quoi.
$$
D'apr\`es~1) a),
on peut supposer, quitte \`a restreindre $S$ \`a un ouvert non
vide, que les $(C_i,fp_i)$ sont cohomologiquement propres relativement
\`a~$S$ en dimension $\le 0$, et que les $f_*{p_i}_*C_i$ sont
constructibles. On en conclut d\'ej\`a que $(G,f)$ est
cohomologiquement propre relativement \`a~$S$ en dimension $\le 0$
\eqref{XIII.1.9}, donc que $(F,f)$ est cohomologiquement propre
relativement \`a~$S$ en dimension $\le -1$, et que $f_*F$ est
constructible. Soit~$K$ la somme amalgam\'ee $K= G \coprod_{F} G$;
comme $F$ et $G$ sont constructibles, il en est de m\^eme de~$K$. On
conclut donc de ce qui pr\'ec\`ede que, quitte \`a restreindre
$S$ \`a un ouvert non vide, on peut supposer que $(K,f)$ est
cohomologiquement propre relativement \`a~$S$ en dimension $\le
-1$. Il r\'esulte alors de \ref{XIII.1.13}~1) que $(F,f)$ est
cohomologiquement propre relativement \`a~$S$ en dimension $\le 0$.

\subsection*{3) Cas des faisceaux en groupes constants}

Si $F$ est un faisceau en groupes constant sur~$X$, on note $\Phi$ le
champ des torseurs sous $F$.

3) 1. Montrons d'abord que, quitte \`a restreindre $S$ \`a un
ouvert non vide, pour tout faisceau en groupes constant $F$ sur~$X$,
d'ordre premier aux caract\'eristiques r\'esiduelles de~$S$,
$(F,f)$ est cohomologiquement propre relativement \`a~$S$ en
dimension $\le 0$, et que $f_*\Phi$ est constructible.

On se ram\`ene pour cela au cas o\`u $X$ est lisse sur~$S$. Quitte
\`a faire une extension finie de~$\kres(s)$, ce qui est loisible car on
peut la consid\'erer comme compos\'ee d'une extension \'etale et
d'une extension radicielle, on peut trouver un morphisme propre
surjectif $p_s \colon X_{1s}\to X_s$, o\`u $X_{1s}$ est un
sch\'ema lisse sur~$S$ de m\^eme dimension que $X_s$, et, quitte
\`a restreindre $S$ \`a un ouvert non vide, on peut supposer que
$p_s$ provient d'un morphisme propre surjectif $p \colon X_1 \to X$,
o\`u $X_1$ est un sch\'ema lisse sur~$S$
(EGA~IV 9.6.1 et 12.1.6). Soient $X_2=X_1 \times_X X_1$, $q \colon
X_2 \to X$ le morphisme canonique. On a un diagramme exact de champs
sur~$X$
$$
\xymatrix@C=.5cm{{\Phi}\ar[r]&{p_*p^* \Phi}\ar@<2pt>[r]\ar@<-2pt>[r]&{q_*q^*\Phi}}
$$
(\ref{XIII.1.11}~2)). Le th\'eor\`eme \'etant suppos\'e
d\'emontr\'e dans le cas lisse, on voit tout d'abord que l'on peut
supposer $f_*p_*p^*\Phi$ constructible; il en est donc de m\^eme de
$f_*\Phi$ (\ref{XIII.3.1.1} ci-dessous). De plus, d'apr\`es \ref{XIII.1.6}~2),
on peut supposer $(p_*p^*\Phi,f)$ cohomologiquement propre
relativement \`a~$S$ en dimension $\le 0$; il en r\'esulte que
$(\Phi,f)$ est cohomologiquement propre relativement \`a~$S$ en
dimension $\le -1$. On peut donc supposer que $(q_*q^*\Phi,f)$ est
cohomologiquement propre relativement \`a~$S$ en dimension $\le -1$,
et cela entra\^ine que $(\Phi,f)$ est cohomologiquement propre
relativement \`a~$S$ en dimension $\le 0$ (\ref{XIII.1.12}~1)).

On se ram\`ene ensuite comme dans 1)2 au cas o\`u $X$ est lisse
et affine sur~$S$. Soit alors
$$
\xymatrix{{X}\ar[r]^{i}\ar[d]_{f}&{P}\ar[ld]^{q}\\{Y}}
$$
une compactification de~$X$, o\`u $i$ est une immersion ouverte
dominante et $q$ un morphisme propre. Comme on a $\dim P_s=\dim X_s$,
on peut appliquer l'hypoth\`ese de r\'esolution des
singularit\'es \`a $P_{\overline s}$. Quitte \`a faire une
extension \'etale et une extension radicielle de~$S$, on peut
trouver un morphisme propre $r \colon Z \to P$, o\`u $Z$ est lisse
sur~$S$, $r^{-1}(X)\simeq X$, et o\`u $r^{-1}(X)$ est le
compl\'ementaire dans $Z$ d'un diviseur \`a croisements normaux
relativement \`a~$S$. Tout torseur sous $F$ est alors
mod\'er\'ement ramifi\'e sur~$Z$ relativement \`a~$S$
(\ref{XIII.2.3}~b)). Il r\'esulte donc de~\ref{XIII.2.7} que $(F,f)$
est cohomologiquement propre relativement \`a~$S$ en dimension $\le
0$, ce qui d\'emontre notre assertion.

\pagebreak[4] 3) 2. R\'eduction au cas o\`u $X$ est lisse sur~$S$.\par\nopagebreak
Quitte \`a faire une extension finie de~$\kres(s)$, on peut trouver un
morphisme propre surjectif $p_s\colon X_{1s}\to X$, o\`u~$X_{1s}$
est un sch\'ema lisse sur~$s$, et on peut supposer que $p_s$
provient d'un morphisme propre surjectif $p\colon X_1 \to X$, o\`u
$X_1$ est lisse sur~$S$. Supposons le th\'eor\`eme
d\'emontr\'e pour $fp$ et montrons-le pour $f$. Soit $F$ un
faisceau en groupes constant fini sur~$X$, d'ordre premier aux
caract\'eristiques r\'esiduelles de~$S$ et $\Phi$ le champ des
torseurs sous $F$. Soient $X_2=X_1 \times_X X_1$, $X_3=X_1 \times_X
X_1 \times_X X_1$, $q\colon X_2 \to X$, $r\colon X_3 \to X$ les
morphismes canoniques. D'apr\`es \ref{XIII.1.11}~2), on a un
diagramme exact de champs
$$
\xymatrix@C=.5cm{{\Phi}\ar[r]&{p_*p^*\Phi}\ar@<2pt>[r]\ar@<-2pt>[r]&{q_*q^*\Phi}\ar@<4pt>[r]\ar[r]\ar@<-4pt>[r]&{r_*r^*\Phi}}\quoi.
$$
Pour prouver que $(\Phi,f)$ est cohomologiquement propre relativement
\`a~$S$ en dimension $\le 1$, il suffit de montrer qu'il en est de
m\^eme de~$(p_*p^*\Phi,f)$, que $(q_*q^*\Phi, f )$ est
cohomologiquement propre relativement \`a~$S$ en dimension $\le 0$
et que $(r_*r^*\Phi,f)$ est cohomologiquement propre relativement
\`a~$S$ en dimension $\le -1$ (\ref{XIII.1.12}~2)). D'apr\`es~3) 1
ci-dessus on peut supposer que, pour tout faisceau en groupes constant
fini~$F$, $(q^*\Phi,fq)$, $(q^*\Phi,q)$, $(r^*\Phi,fr)$, $(r^*\Phi,r)$
sont cohomologiquement propres relativement \`a~$S$ en dimension
$\le 0$. Il r\'esulte alors de \ref{XIII.1.6}~2) que
$(q_*q^*\Phi,f)$ et $(r_*r^*\Phi,f)$ sont cohomologiquement propres
relativement \`a~$S$ en dimension $\le 0$. Le th\'eor\`eme
\'etant suppos\'e d\'emontr\'e dans le cas lisse, $(p^*\Phi,
fp)$ et $(p^*\Phi,p)$ donc aussi $(p_*p^*\Phi,f)$ (\ref{XIII.1.6}~2))
sont cohomologiquement propres relativement \`a~$S$ en dimension
$\le 1$. Ceci montre bien que $(\Phi,f)$ est cohomologiquement propre
relativement \`a~$S$ en dimension $\le 1$.

De plus $f_*p_*p^*\Phi$ est $1$-constructible par hypoth\`ese;
d'apr\`es~\ref{XIII.3.1} on peut supposer que $f_*q_*q^*\Phi$ est
constructible; il r\'esulte donc de~\ref{XIII.3.1.1} ci-dessous que
$f_*\Phi$ est $1$-constructible.

3) 3. R\'eduction au cas o\`u $X$ est lisse affine sur~$S$.

D'apr\`es 3)2, on peut supposer $X$ lisse sur~$S$. Soit $(U_i)_{i\in
I}$ un recouvrement fini de~$X$ par des ouverts affines, et soit $X_1$
la somme directe des $U_i$, $p\colon X_1 \to X$ le morphisme
canonique. Comme $p$ est un morphisme de descente effective pour la
cat\'egorie des faisceaux \'etales de type fini sur des
sch\'emas variables, on voit comme dans 3)2 que, si l'on suppose le
th\'eor\`eme d\'emontr\'e pour les $U_i$, \ie pour $X_1$, il
est aussi vrai pour~$X$.

3) 4. Cas o\`u $X$ est lisse affine sur~$S$.

On voit comme dans 3)1 que, quitte \`a restreindre $S$ \`a un
ouvert non vide et \`a faire une extension \'etale et une
extension radicielle surjectives de~$S$, on peut trouver un diagramme
commutatif
$$
\xymatrix{{X}\ar[r]^{i}\ar[d]_{f}&{P}\ar[ld]^{q}\\{Y}}
$$
o\`u $P$ est un sch\'ema lisse sur~$S$, $X$ le compl\'ementaire
dans $P$ d'un diviseur \`a croisements normaux relativement \`a~$S$ et $g$ un morphisme propre. Si $F$ est un faisceau constant
d'ordre premier aux caract\'eristiques r\'esiduelles de~$S$, tout
torseur sous $F$ est mod\'er\'ement ramifi\'e sur~$P$
relativement \`a~$S$ (\ref{XIII.2.3}~b)). Le fait que $(F,f)$ soit
cohomologiquement propre relativement \`a~$S$ en dimension $\le 1$
et que $f_*\Phi$ soit constructible r\'esulte alors
de~\ref{XIII.2.7}.

\subsection*{4) D\'emonstration de 2) b)}

4) 1. Cas o\`u $\Phi$ est une gerbe.

On peut trouver un morphisme \'etale surjectif de type fini $p
\colon X_1 \to X$, tel que $p^*\Phi$ soit un gerbe triviale. Par
descente, comme dans 3)2, on voit qu'il suffit de prouver le
th\'eor\`eme pour $X_1$, $X_1\times_X X_1$ et
$X_1\times_X X_1\times_X X_1$. On est donc ramen\'e au cas o\`u
$\Phi$ est le gerbe des torseurs sous un faisceau en groupes
constructible $F$, dont les fibres sont d'ordre premier aux
caract\'eristiques r\'esiduelles de~$S$.

D'apr\`es SGA~4 IX~2.14, on peut trouver une famille finie de
morphismes finis $p_i \colon Z_i \to X$, et, pour chaque $i$, un
faisceau en groupes constant fini $C_i$, d'ordre premier aux
caract\'eristiques r\'esiduelles de~$S$, de sorte que l'on ait un
monomorphisme
$$
j\colon F \to \prod_{i} p_{i*}C_i=G \quoi.
$$
Soit $\Phi_i$ le champ de torseurs sous $C_i$ et $\Psi$ le champ des
torseurs sous $G$. Il r\'esulte de 2) a) que, quitte \`a
restreindre $S$ \`a un ouvert non vide, on peut supposer que les
$(C_i,fp_i)$ sont cohomologiquement propres relativement \`a~$S$ en
dimension $\le 1$, et que les champs $f_*p_{i*}\Phi_i$ sont
$1$-constructibles. Il r\'esulte alors de~\ref{XIII.1.9} que les
$(p_{i*}C_i,f)$ sont cohomologiquement propres relativement \`a~$S$
en dimension $\le 1$; il en est donc de m\^eme de~$(G,f)$. De plus,
comme $p_{i*}\Phi_i$ est \'equivalent au champ des torseurs sous le
groupe $p_{i*}C_i$ (SGA~4 VIII~5.8), on voit que $f_*\Psi$ est
$1$-constructible.

Comme $\R^1f_*G$ est constructible, on peut trouver un faisceau
repr\'esentable par un $Y$-sch\'ema \'etale de type fini $T$ et
un \'epimorphisme
$$
a \colon T \to \R^1f_*G
$$
(SGA~4 IX~2.7); de plus on peut supposer que l'image de la section
identique de~$T(T)$ est d\'efinie par un torseur $Q$ sur~$X\times_Y
T=X_T$, de groupe $G|X_T$. Soient $f_T \colon X_T \to T$ le morphisme
canonique, $F_T=F|X_T$, etc. D'apr\`es 1) b) on peut supposer,
quitte \`a restreindre $S$ \`a un ouvert non vide, que
$(Q/F_T,f_T)$ est cohomologiquement propre relativement \`a~$S$ en
dimension $\le 0$ et que $f_{T*}(Q/F_T)$ est constructible. Il
r\'esulte alors de~\ref{XIII.3.1.2} que $f_*\Psi$ est constructible.

Montrons que $(F,f)$ est cohomologiquement propre relativement \`a
$S$ en dimension $\le 1$. D'apr\`es \ref{XIII.1.13}~2) il suffit de
prouver que, pour tout sch\'ema $Y_1$ \'etale sur~$Y$ et pour tout
torseur $Q_1$ sur~$X_1=X\times_Y Y_1$, si $f_1 \colon X_1 \to Y_1$ est
le morphisme canonique, alors $(Q_1/F_1,f_1)$ est cohomologiquement
propre relativement \`a~$S$ en dimension $\le 0$. Or, par
d\'efinition de~$T$, $Q_1$ est, localement pour la topologie
\'etale de~$Y_1$, image inverse de~$Q$, ce qui d\'emontre notre
r\'eduction.

4) 2. Cas g\'en\'eral.\par\nopagebreak
On voit en utilisant le lemme~\ref{XIII.6.1.1}, 4)1 et 1) a) que,
quitte \`a restreindre $S$ \`a un ouvert non vide, on peut
supposer $S(f_*\Phi)$ constructible et $(S\Phi,f)$ cohomologiquement
propre relativement \`a~$S$ en dimension $\le 0$. On peut alors
trouver un faisceau repr\'esentable par un $Y$-sch\'ema \'etale
de type fini $T$ et un \'epimorphisme
$$
a \colon T \to S(f_*\Phi)\quoi.
$$
On reprend les notations de 4)1, et on pose de plus $Z=T\times_Y T$,
$X_Z=X\times_Y Z$ et on note $f_Z \colon X_Z \to Z$ le morphisme
canonique. On peut supposer $T$ choisi de sorte que l'image $q$ de la
section identique de~$T(T)$ par $a$ soit d\'efinie par un objet~$p$
de~$(f_*\Phi)_T=\Phi_{X_T}$. Soient $p_1$ et $p_2$ (\resp~$q_1$ et
$q_2$) les images inverses de~$p$ (\resp~$q$) par les deux projections
de~$Z$ dans $T$. On peut supposer, quitte \`a restreindre $S$ \`a
un ouvert non vide, que $(\SheafAut_{X_T}(p),f_T)$ est
cohomologiquement propre relativement \`a~$S$ en dimension $\le 1$,
que $(\SheafHom_{X_Z}(p_1,p_2),f_Z)$ est cohomologiquement propre
relativement \`a~$S$ en dimension $\le 0$, et que
$f_{T*}(\SheafAut_{X_T}(p))$, $f_{Z*}(\SheafHom_{X_Z}(p_1,p_2))$ sont
constructibles ( a)1) et 4)1)).

On en d\'eduit d'abord que, pour tout sch\'ema $Y_1$ \'etale sur~$Y$ et pour tout couple d'objets $y$, $y_1$ de~$(f_*\Phi)_{Y_1}$, le
faisceau $\SheafAut_{Y_1}(y)$ (\resp $\SheafHom_{Y_1}(y,y_1)$) est
constructible, un tel faisceau \'etant, localement pour la topologie
\'etale de~$Y_1$, image inverse de~$f_{T*}(\SheafAut_{X_T}(p))$
(\resp de~$f_{Z*}(\SheafHom_{X_Z}(p_1,p_2))$).

Il reste \`a prouver que $(\Phi,f)$ est cohomologiquement propre
relativement
\`a~$S$ en dimension $\le 1$. Il suffit pour cela de montrer que,
pour tout $S'$-sch\'ema $S$ et pour tout point g\'eom\'etrique
$\overline{y}'$ de~$Y$, si l'on note $\overline{Y}$ le localis\'e
strict de~$Y$ en $y'$, $\overline{X}=X \times_Y \overline{Y}$, etc.,
alors le foncteur canonique
$$
\overline{\varphi} \colon \Phi(\overline{X}) \to
\Phi'(\overline{X'})
$$
est une \'equivalence de cat\'egories.

Montrons que $\overline{\varphi}$ est pleinement fid\`ele. Soient
$x$, $y \in \Phi(\overline{X})$, $x'$, $y'$ leurs images dans
$\Phi'(\overline{X'})$ et montrons que le morphisme canonique
\begin{equation*}
\label{eq:XIII.4.*} \tag{$*$} {\Hom_{\overline{X}}(x,y)
\to \Hom_{\overline{X'}}(x',y') }
\end{equation*}
est bijectif. Par d\'efinition de~$T$ il existe deux morphismes de
$\overline{Y}$ dans $T$ tels que $x$ et~$y$ soient les images inverses
de~$p$ par ces deux morphismes. Cela revient \`a dire qu'il y a un
morphisme $\overline {Y} \to Z$, d'o\`u un morphisme $h\colon
\overline{X} \to X_Z$ tel que l'on ait
$$
h^*(p_1)=x \qquad h^*(p_2)=y.
$$
Par suite on a un isomorphisme canonique
$$
\SheafHom_{\overline{X}}(x,y)=h^*(\SheafHom_{X_Z}(p_1,p_2) )\quoi.
$$
mais, compte tenu du fait que $(\SheafHom_{X_Z}(p_1,p_2),f_Z)$ est
cohomologiquement propre relativement \`a~$S$ en dimension $\le 0$,
on voit qu'il en est de m\^eme de
$(\SheafHom_{\overline{X}}(x,y),\overline{f})$, ce qui prouve que le
morphisme \eqref{eq:XIII.4.*} est bijectif.

Montrons que $\overline{\varphi}$ est essentiellement surjectif. Soit
$x'\in \Phi'(\overline{X'})$. Comme $(S\Phi,f)$ est cohomologiquement
propre relativement \`a~$S$ en dimension $\le 0$, le morphisme
canonique
$$
\H^0(\overline{X},\overline{S\Phi}) \to
\H^0(\overline{X'},\overline{S\Phi}')
$$
est bijectif. Soit $G'$ la sous-gerbe maximale de~$\Phi'$
engendr\'ee par $x'$; il existe alors une sous-gerbe maximale $G$ de
$\overline{\Phi}$ dont l'image inverse sur~$\overline{X'}$ est
$G'$. D'apr\`es 4)1 $(G,\overline{f})$ est cohomologiquement propre
relativement \`a~$S$ en dimension $\le 1$; par suite le foncteur
canonique
$$
G(\overline{X}) \to G'(\overline{X'})
$$
est une \'equivalence de cat\'egories, ce qui prouve l'existence
d'un \'el\'ement $x$ de~$\Phi(\overline{X})$ dont l'image dans
$\Phi'(X')$ soit isomorphe \`a $x'$ et ach\`eve la
d\'emonstration du th\'eor\`eme.

\begin{sublemme}
\label{XIII.3.1.1}
Soient $S$ un sch\'ema localement noeth\'erien et
$$
\xymatrix@C=.5cm{{\Phi}\ar[r]^p&{\Phi_1}\ar@<2pt>[rr]^{p_1,p_2}\ar@<-2pt>[rr]&&{\Phi_2}}
$$
un diagramme exact de champs sur~$S$ \eqref{XIII.1.10.1}. Si $\Phi_1$
est constructible, il en est de m\^eme de~$\Phi$. Si, pour tout
sch\'ema $S'$ \'etale sur~$S$ et pour tout couple
d'objets
$x_1,~y_1$ de~$(\Phi_1)_{S'}$ le faisceau $\SheafHom_{S'}(x_1,y_1)$
est constructible, alors, pour tout couple d'objets $x,~y$ de
$\Phi_{S'}$, il en est de m\^eme de~$\SheafHom_{S'}(x,y)$. Supposons
que $\Phi_1$ soit $1$-constructible (\ref{XIII.0}) et que $\Phi_2$ soit
constructible, alors $\Phi$ est $1$-constructible.
\end{sublemme}

Pour tout sch\'ema $S'$ \'etale sur~$S$ et pour tout objet $x$ de
$\Phi_{S'}$, on a un monomorphisme
$$
\SheafAut_{S'}(x)\to\SheafAut_{S'}(p(x)).
$$
Il r\'esulte donc de SGA~4 IX~2.9 que, si $\Phi_1$ est
constructible, il en est de m\^eme de~$\Phi$, et la deuxi\`eme
assertion du lemme se d\'emontre de la m\^eme fa\c con.

Supposons maintenant $\Phi_1$ $1$-constructible et $\Phi_2$
constructible. Le morphisme $p$ induit sur les faisceaux de
sous-gerbes maximales un morphisme
$$
\varphi\colon S\Phi\to S\Phi_1.
$$
Soit $G$ l'mage de~$S\Phi$ par $\varphi$; d'apr\`es SGA~4 IX~2.9,
$G$ est un faisceau constructible. On peut donc trouver un faisceau
repr\'esentable par un $S$-sch\'ema \'etale de type fini $T$ et
un \'epimorphisme
$$
a\colon T\to G
$$
(SGA~4 IX~2.7).
De plus on peut choisir $T$ de sorte que l'image $y$ de section
identique de~$T(T)$ par $a$ soit d\'efinie par un objet $x_1$ de
$(\Phi_1)_T$ de la forme $x_1=p(x)$, o\`u $x\in\Phi_T$.

Il suffit de montrer que, pour tout point $s$ de~$S$, il existe un
ouvert non vide~$U$ de~$\overline{\{ s\} }$ tel que $S\Phi|U$ soit
localement constant constructible. Soient $s\in S$, $\overline{s}$ un
point g\'eom\'etrique au-dessus de~$s$ et $\overline{y}_1,\dots,\overline{y}_n$ les \'el\'ements de~$G_{\overline{s}}$. Par
d\'efinition de~$T$ il existe des morphismes
$h_i\colon\overline{s}\to T$ tels que l'on ait
$h_i^*(y)=\overline{y}_i$. Soit $S'$ le produit fibr\'e sur~$S$ de
$n$ sch\'emas isomorphes \`a $T$, $h\colon\overline{s}\to S'$ le
produit fibr\'e des $h_i$, $y_i$ (\resp $x_i$) l'image inverse de
$y$ (\resp $x$) par la $i$-\`eme projection de~$S'$ sur~$T$. Soit
$F_i$ le sous-faisceau de~$S\Phi |S'$ image r\'eciproque de~$y_i$
et montrons que les $F_i$ sont constructibles. Le faisceau $F_i$ est
un quotient du faisceau $F'_i$ tel que, pour tout sch\'ema $S''$
\'etale au-dessus de~$S'$, on ait
\begin{multline*}
F'_i(S'') = \{\text{classes mod. isomorphisme d'objet $z$ de
$\Phi_{S''}$ muni}\\
\text{ d'un isomorphisme $i\colon p(z)\simeq p(x_i|S'')$}\}.
\end{multline*}
Il suffit de montrer que les $F_i'$ sont constructibles. Or, si l'on
pose $z_i=p_1p(x_i)$, $z_i'=p_2p(x_i)$, on a un monomorphisme
$$
\Psi_i\colon F'_i\to\SheafIsom_{S'}(z_i,z_i'),
$$
obtenu en associant, \`a tout sch\'ema $S''$ \'etale sur~$S'$ et
\`a tout objet $z$ de~$\Phi_{S''}$ tel que l'on ait un isomorphisme
$i\colon p(z)\simeq p(x_i|S'')$, l'isomorphisme de~$z_i$ dans $z_i'$
d\'efini par la condition que le diagramme
$$
\xymatrix@C=1.2cm{{p_1p(z)}\ar[r]^-{p_1(i)}\ar[d]_-j&{z_i}\ar[d]\\{p_2p(z)}\ar[r]^-{p_2(i)}&{z'_i}}
$$
soit commutatif ($j$ est le morphisme canonique associ\'e au
diagramme exact).

Le morphisme $\Psi_i$ est injectif car dire que deux objets $z,~z'$,
tels que l'on ait des isomorphismes $i\colon p(z)\isomto
p(x_i|S''),~i'\colon p(z')\isomto p(x_i|S'')$, d\'efinissent le
m\^eme \'el\'ement de~$\SheafIsom_{S''}(z_i,z'_i)$ revient \`a
dire que l'on a $p_1(i'^{-1}i)=p_2(i'^{-1}i)$, \ie que $i'^{-1}i$
provient d'un isomorphisme $z\to z'$. Le faisceau $F'_i$ \'etant un
sous-faisceau de~$\SheafIsom_{S'}(z_i,z'_i)$ est constructible.

Comme on peut trouver un ouvert non vide~$U$ de~$\overline{\{ s\}}$
tel que $y_1|U,\dots,y_n|U$ engendrent $G$, il r\'esulte du
lemme~\ref{XIII.6.1.2} ci-dessous que $S\Phi |U$ est constructible,
donc, quitte \`a restreindre $U$, $S\Phi |U$ est localement constant
constructible.

\begin{sublemme}
\label{XIII.3.1.2}
Soient $S$ un sch\'ema localement noeth\'erien, $f\colon X\to S$
un morphisme, $F\to G$ un monomorphisme de faisceaux en groupes sur~$X$, $\Psi$ (\resp $\Psi_1$) le champ des torseurs sous $F$
(\resp sous $G$), $\Phi=f_*\Psi$, $\Phi_1=f_*\Psi_1$. On suppose
donn\'e un faisceau sur~$S$ repr\'esentable par un $S$-sch\'ema
\'etale de type fini $T$ et un morphisme surjectif
$$
a\colon T\to S\Phi_1\simeq\R^1f_*G,
$$
de sorte qu'il existe un torseur $Q$ sur~$X_T=X\times_ST$ de groupe
$G|X_T$ qui d\'efinisse dans $\R^1f_*G(T)$ l'image par $a$ de la
section identique de~$T(T)$. Soit $f_T\colon X_T\to T$ le morphisme
canonique et posons $F_T=F|X_T$. On suppose que $\Phi_1$ est
$1$-constructible et que $f_{T*}(Q/F_T)$ est constructible. Alors
$\Phi$ est $1$-constructible.
\end{sublemme}

Il suffit en effet de recopier la d\'emonstration
de~\ref{XIII.3.1.1}, le fait que l'on ait des morphismes $\Psi_i$
\'etant remplac\'e par le fait que l'on a des isomorphismes
$$
F'_i\isomto f_*(Q/F_T)|S'.
$$

\begin{subremarque}
\label{XIII.3.1.3}
Supposons
que $k=\kres(s)$ soit de caract\'eristique nulle; alors les
sch\'emas de type fini sur~$\overline{k}$ de dimension $\leq\dim X$
sont fortement d\'esingularisables et la d\'emonstration
de~\ref{XIII.3.1} permet de prouver les r\'esultats suivants:
\begin{enumerate}
\item[a)] Il existe un ouvert non vide~$S_1$ de~$S$ tel que, pour tout
sch\'ema $S'$ au-dessus de~$S_1$ dont les points maximaux sont de
caract\'eristique nulle et pour tout faisceau d'ensembles localement
constant constructible $F$ sur~$X'=X\times_SS'$, $(F,f')$ soit
cohomologiquement propre relativement \`a~$S'$ en dimension $\leq
0$.

\item[b)] Si toutes les caract\'eristiques r\'esiduelles de~$S$ sont
nulles, il existe un ouvert non vide~$S_1$ de~$S$ tel que, pour tout
sch\'ema $S'$ au-dessus de~$S_1$ et pour tout faisceau en groupes
localement constant constructible $F$ sur~$X'$, $(F,f')$ soit
cohomologiquement propre relativement \`a~$S'$ en dimension $\leq1$.
\end{enumerate}

Il suffit en effet de recopier la d\'emonstration de
\ref{XIII.3.1}.2) a). La proposition~\ref{XIII.2.7} utilis\'ee dans
3) 4 s'applique au cas d'un faisceau localement constant $F$, car,
reprenant les notations de 3)4, tout torseur sous $F$ est
mod\'er\'ement ramifi\'e sur~$P$ relativement \`a~$S$ puisque
toutes les caract\'eristiques r\'esiduelles de~$S$ sont nulles.
\end{subremarque}

\begin{corollaire}
\label{XIII.3.2}
Soient $k$ un corps de caract\'eristique $p\geq 0$, $p'$ l'ensemble
des nombres premiers distincts de~$p$, $f\colon X\to k$ un morphisme
coh\'erent.
\begin{enumerate}
\item[1)] Pour tout faisceau d'ensembles $F$, $(F,f)$ est
cohomologiquement propre en dimension $\leq 0$.
\item[2)] Supposons satisfaite l'une des deux conditions suivantes:
\begin{enumerate}
\item[a)] $f$ est de type fini et les sch\'emas de type fini de
dimension $\leq\dim X$ sur une cl\^oture alg\'ebrique de~$k$ sont
fortement d\'e\-sin\-gu\-la\-ri\-sables.
\item[b)] Les sch\'emas de type fini sur une cl\^oture
alg\'ebrique de~$k$ sont
fortement d\'esingularisables.
\end{enumerate}
Alors, pour tout faisceau de ind-$p'$-groupe $F$, $(F,f)$ est
cohomologiquement propre en dimension $\leq 1$.
\end{enumerate}
\end{corollaire}

Soit $F$ un faisceau d'ensembles (\resp de
ind-$p'$-groupes). D'apr\`es SGA~4 IX~2.7.2 on peut \'ecrire
$F$ comme limite inductive filtrante
$$
F=\varinjlim F_i,
$$
o\`u les $F_i$ sont des faisceaux d'ensembles (\resp de
ind-$p'$-groupes) constructibles. Comme $f$ est coh\'erent, $f_*$
(\resp $\R^1f_*$) commute aux limites inductives (SGA~4 VII~3.3). Si
l'on sait que les $(F_i,f)$ sont cohomologiquement propres en
dimension $\leq 0$ (\resp que tout faisceau d'ensemble est
cohomologiquement propre en dimension $\leq 0$ et que les $(F_i,f)$
sont cohomologiquement propres en dimension $\leq 1$), il en sera de
m\^eme de~$(F,f)$. On peut donc supposer $F$ constructible.

Si l'on suppose $f$ de type fini (\resp satisfaisant \`a a)), la
proposition r\'esulte de \ref{XIII.3.1}~1) b) (\resp de
\ref{XIII.3.1}~2) b)). Prouvons maintenant~\ref{XIII.3.2} quand on ne
suppose plus $f$ de type fini. Pour tout sch\'ema $S'$ au-dessus de
$k$ et pour tout point g\'eom\'etrique $\overline{s}$ de~$S'$, on
note $\overline{k}$ (\resp $\overline{S'}$) le localis\'e strict de
$k$ en $\overline{s}$ (\resp de~$S'$ en $\overline{s}$),
$\overline{X}$ l'image inverse de~$X$ sur~$\overline{k}$, et on
consid\`ere le carr\'e cart\'esien
$$
\xymatrix@C=1.2cm{{\overline{X}}\ar[d]_-{\overline{f}}&{\overline{X'}}\ar[l]_g\ar[d]^-{\overline{f}'}\\{\overline{k}}&{\overline{S'}}\ar[l]}
$$
Il suffit de prouver que, pour tout $S'$, pour tout $\overline{s}$, le
morphisme canonique
$$
\H^0(\overline{X},\overline{F})\to
\H^0(\overline{X'},\overline{F'})\qquad
\text{(\resp $\H^1(\overline{X},\overline{F})\to
\H^1(\overline{X'},\overline{F'})$)}
$$
est un isomorphisme. Il suffit de montrer que l'on a les relations
\begin{equation*}
\label{eq:XIII.3.*}
\tag{$*$} \overline{F}\simeq g_*g^*\overline{F}\qquad
\text{(\resp $\R^1g_*(g^*\overline{F})=0$).}
\end{equation*}
Or, sous cette forme, la question est locale sur~$X$ pour la topologie
\'etale. On peut donc supposer $X$ affine; par passage \`a la
limite on peut supposer $X$ de type fini sur~$k$. On sait alors que
$(F,f)$ est cohomologiquement propre en dimension $\leq 0$
(\resp $\leq 1$) et qu'il en est de m\^eme quand on remplace $X$ par
un sch\'ema \'etale de type fini sur~$X$, ce qui prouve
\eqref{eq:XIII.3.*}.

\begin{theoreme}
\label{XIII.3.3}
Soient $S$ un sch\'ema irr\'eductible de point g\'en\'erique
$s$, $f\colon X\to S$ un morphisme de pr\'esentation
finie. Supposons que les sch\'emas de type fini de dimension
$\leqslant\dim X_s$ sur une cl\^oture alg\'ebrique $\bar{k}$ de
$k$ soient d\'esingularisables \textup{(EGA~IV 7.9.1)}. Alors, si $\LL$
d\'esigne l'ensemble des nombres premiers distincts des
caract\'eristiques r\'esiduelles de~$S$, on peut trouver un ouvert
non vide~$S_1$ de~$S$ tel que le morphisme $f|S_1$ soit
universellement localement $1$-asph\'erique pour~$\LL$.
\end{theoreme}

On peut supposer $S$ int\`egre et $X$ r\'eduit (SGA~4
VIII~1.1). Par passage \`a la limite on peut supposer $S$
noeth\'erien. De plus, pour d\'emontrer le th\'eor\`eme, il
suffit de le faire apr\`es extension finie $S_1\to S$, o\`u $S_1$
est un sch\'ema int\`egre et o\`u $S_1\to S$ est compos\'e
d'extensions \'etales et d'extensions radicielles surjectives.

Montrons d'abord que, quitte \`a restreindre $S$ \`a un ouvert non
vide, $f$ est universellement localement $0$-acyclique. Quitte \`a
faire une extension radicielle de~$\kres(s)$, on peut supposer que le
morphisme $(X_s)_\red\to s$ est s\'eparable; on peut donc supposer,
quitte \`a restreindre $S$ \`a un ouvert non vide et \`a faire
une extension radicielle surjective de~$S$, que le morphisme $f$ est
plat, \`a fibres g\'eom\'etriques s\'eparables (EGA~IV
12.1.1), ce qui entra\^ine que $f$ est universellement
$0$-acyclique (SGA~4 XV~4.1).

Montrons que, quitte \`a restreindre $S$ \`a un ouvert non vide,
$f$ est universellement localement $1$-asph\'erique pour
$\LL$. Quitte \`a faire une
extension finie de~$\kres(s)$, ce qui est loisible car on peut la
consid\'erer comme compos\'ee d'une extension \'etale et d'une
extension radicielle, on peut trouver un morphisme propre surjectif
$p_s\colon Y_s\to X_s$, o\`u $Y_s$ est un sch\'ema lisse sur~$S$,
de m\^eme dimension que $X_s$, et on peut supposer que $p_s$
provient d'un morphisme propre surjectif $p\colon Y\to X$, o\`u $Y$
est un sch\'ema lisse sur~$S$ (EGA~IV 9.6.1 et 12.1.6). Il suffit de
montrer que, quitte \`a restreindre $S$ \`a un ouvert non vide,
pour tout diagramme \`a carr\'es cart\'esiens
$$
\xymatrix{{S''}\ar[d]_i&{X ''}\ar[d]^j\ar[l]_{f''}\\{S'}\ar[d]&{X'}\ar[d]\ar[l]_{f'}\\{S}&{X,}\ar[l]_f}
$$
o\`u $i$ est \'etale de pr\'esentation finie, et pour tout
faisceau de ind-$\LL$-groupes $F$ sur~$S''$, si $\Phi$ est le champ
des torseurs sous $F$, alors le morphisme canonique
$$
{f'}^*i_*\Phi \to j_*{f''}^*\Phi
$$
est une \'equivalence. Soient $Z=Y\times_{X}Y,
T=Y\times_{X}Y\times_{X}Y$. On a de fa\c con naturelle un diagramme
commutatif
$$
\xymatrix{{S''}\ar[d]_i&{X''}\ar[l]_{f''}\ar[d]_j&{Y''}\ar[l]_{p''}\ar[d]&{Z''}\ar[d]\ar@<-.4ex>[l]\ar@<.4ex>[l]&{T''}\ar[d]\ar[l]\ar@<-.8ex>[l]\ar@<.8ex>[l]\\{S'}\ar[d]&{X'}\ar[l]_{f'}\ar[d]&{Y'}\ar[d]\ar[l]_{p'}&{Z'}\ar[d]\ar@<-.4ex>[l]\ar@<.4ex>[l]&{T'}\ar[d]\ar[l]\ar@<-.8ex>[l]\ar@<.8ex>[l]\\{S}&{X}\ar[l]_f&{Y}\ar[l]_p&{Z'}\ar@<-.4ex>[l]\ar@<.4ex>[l]&{T'}\ar[l]\ar@<-.8ex>[l]\ar@<.8ex>[l]}
$$
Soient
$q\colon Z\to X$ et $r\colon T\to X$ les morphismes canoniques,
les notations $q', r', q'', r''$ ayant un sens
\'evident. D'apr\`es \ref{XIII.1.11}~2), on a le diagramme
essentiellement commutatif suivant, dont les lignes sont exactes:
$$
\xymatrix{{f'^* i_*\Phi}\ar[r]\ar[d]_a&{p'_*p'^*(f'^*i_*\Phi)}\ar[d]_b\ar@<-.4ex>[r]\ar@<.4ex>[r]&{q'_*q'^*(f'^*i_*\Phi)}\ar[d]_c\ar[r]\ar@<-.8ex>[r]\ar@<.8ex>[r]&{r'_*r'^*(f'^*i_*\Phi)}\ar[d]_d\\{j_*f''^*\Phi}\ar[r]&{j_*p''_*p''^*(f''^*\Phi)}\ar@<-.4ex>[r]\ar@<.4ex>[r]&{j_*q''_*q''^*(f''^*\Phi)}\ar[r]\ar@<-.8ex>[r]\ar@<.8ex>[r]&{j_*r''_*r''^*(f''^*\Phi)}}
$$
Comme $Y$ est lisse sur~$S$, le morphisme $fp$ est universellement
localement $1$-asph\'erique pour $\LL$ (SGA~XV~2.1),
et il r\'esulte de \cite[VII 2.1.7]{XIII.2} que $b$ est une \'equivalence de
cat\'egories. D'autre part, quitte \`a restreindre $S$ \`a un
ouvert non vide, on peut supposer que les morphismes $Z\to S$ et $T\to
S$ sont universellement localement $0$-acycliques. Il en r\'esulte
que les foncteurs $c$ et $d$ sont pleinement fid\`eles, et le
diagramme ci-dessus montre alors que $a$ est une \'equivalence, ce
qui ach\`eve la d\'emonstration.
\begin{corollaire}
\label{XIII.3.4}
Soient $k$ un corps de caract\'eristique $p\geqslant0$, $p'$
l'ensemble des nombres premiers distincts de~$p$, $f\colon X\to k$ un
morphisme coh\'erent. Supposons satisfaite l'une des deux
conditions suivantes:
\begin{enumerate}
\item[a)] $f$ est de type fini et les sch\'emas de type fini de
dimension $\leqslant\dim X$ sur une cl\^oture alg\'ebrique de~$k$
sont d\'esingularisables.
\item[b)] Les sch\'emas de type fini sur une cl\^oture
alg\'ebrique de~$k$ sont d\'esingularisables.
\end{enumerate}
Alors $f$ est universellement localement $1$-asph\'erique pour $p'$.
\end{corollaire}

Le cas a) r\'esulte de~\ref{XIII.3.3}. Dans le cas b), la question
\'etant locale sur~$X$, on peut supposer $X$ affine; par passage
\`a la limite (SGA~4 XV~1.3), on se ram\`ene au cas o\`u $X$ est
de type fini sur~$k$.
\begin{corollaire}
\label{XIII.3.5}
Soient
$S$ un sch\'ema irr\'eductible de point g\'en\'erique
$s$, $f\colon X\to S$ un morphisme de pr\'esentation
finie. Supposons que les sch\'emas de type fini de dimension
$\leqslant\dim X_s$ sur une cl\^oture alg\'ebrique $\overline{k}$
de~$\kres(s)$ soient fortement d\'esingularisables \textup{(SGA~5 I~3.1.5)}. Si
$\LL$ d\'esigne l'ensemble des nombres premiers distincts des
caract\'eristiques r\'esiduelles de~$S$, on peut trouver un ouvert
non vide~$S_1$ de~$S$ tel que, pour toute sp\'ecialisation
$\bar{s}_1\to\bar{s}_2$ de points g\'eom\'etriques de~$S_1$, le
morphisme de sp\'ecialisation \eqref{XIII.2.10}
$$
\pi_1^{\LL}(X_{\bar{s}_1})\to\pi_1^{\LL}(X_{\bar{s}_2})
$$
soit bijectif.
\end{corollaire}

D'apr\`es~\ref{XIII.3.1} et~\ref{XIII.3.3}, on peut, quitte \`a
restreindre $S$ \`a un ouvert non vide, supposer que $f$ est
localement $1$-asph\'erique pour $\LL$, et que, pour tout faisceau
de~$\LL$-groupes constant fini $F$ sur~$X$, $(F,f)$ est
cohomologiquement propre en dimension $\leqslant 1$. Il r\'esulte
alors de~\ref{XIII.1.14} que, pour toute sp\'ecialisation
$\bar{s}_1\to\bar{s}_2$ de points g\'eom\'etriques de~$S_1$, le
morphisme de sp\'ecialisation
$$
\left(\R^1 f_* F\right)_{\bar{s}_2} \to \left(\R^1 f_*
F\right)_{\bar{s}_1}
$$
est bijectif. Le corollaire n'est autre que la traduction de ce qui
pr\'ec\`ede en termes de groupes fondamentaux.

\section{Suites exactes d'homotopie}\label{XIII.4}\setcounter{subsection}{-1}
\subsection{}
\label{XIII.4.0}
Soient $X$ et $S$ deux sch\'emas connexes, $f\colon X\to S$ un
morphisme, $a$ un point g\'eom\'etrique de~$X$, $\LL$ un ensemble
de nombres premiers. Soit $K$ le noyau de l'homomorphisme canonique
$\pi_1(X,a)\to\pi_1(S,a)$ et $N$ le plus petit pro-sous-groupe
distingu\'e de~$K$ tel que $K/N$ soit un pro-$\LL$-groupe
$K^{\LL}$. Alors $N$ est distingu\'e dans $\pi_1(X,a)$ et on note
$$
\pi_1'(X,a)
$$
\label{indnot:mj}le
quotient de~$\pi_1(X,a)$ par $N$. Si $a$ est un point
g\'eom\'etrique d'une fibre g\'eom\'etrique~$X_{\bar{s}}$, les
morphismes canoniques
$$
\pi_1(X_{\bar{s}},a)\to \pi_1(X,a) \to
\pi_1(S,a)
$$
permettent de d\'efinir des morphismes canoniques
$$
\pi_1^{\LL}(X_{\bar{s}},a)\lto{u} \pi_1'(X,a) \lto{v} \pi_1(S,a)
$$
On a $vu=0$.

\begin{proposition}
\label{XIII.4.1}
Soient $S$ un sch\'ema connexe, $f\colon X\to S$ un morphisme
localement $0$-acyclique \textup{(SGA~4 XV~1.11)}; supposons de plus $f$
$0$-acyclique (ce qui, lorsque $f$ est coh\'erent, revient \`a
dire que les fibres g\'eom\'etriques de~$f$ sont connexes \textup{(SGA~4
XV~1.16)}). Soit $\LL$ un ensemble de nombres premiers. Si $S'$ est un
sch\'ema \'etale sur~$S$, on note $X'$, $f'$ les images inverses
de~$X$, $f$ sur~$S'$. Supposons que, pour tout rev\^etement
\'etale $S'$ de~$S$ et pour tout rev\^etement \'etale $E$ de
$X'$, quotient d'un rev\^etement galoisien de groupe un
$\LL$-groupe, $(E,f')$ soit cohomologiquement propre relativement
\`a~$S'$ en dimension $\leqslant 0$ et $f'_* E$
constructible. Alors, si $\bar{s}$ est un point g\'eom\'etrique de
$S$ et $a$ un point g\'eom\'etrique de la fibre $X_{\bar{s}}$, la
suite d'homomorphismes de groupes
\begin{equation*}
\label{eq:XIII.4.1.1}
\tag{\thesubsection.1}
\pi_1^{\LL}(X_{\bar{s}},a)\lto{u} \pi_1'(X,a)\lto{v}\pi_1(S,a)\to1
\end{equation*}
est exacte
\end{proposition}

Cet \'enonc\'e g\'en\'eralise X.\ref{X.1.4}, dont on va copier
la d\'emonstration.

Montrons d'abord que $v$ est surjectif. Il suffit de montrer que, pour
tout rev\^etement \'etale connexe $S'$ de~$S$, $X'$ est aussi
connexe (V~\ref{V.6.9}). Soit $C$ un ensemble ayant au moins deux
\'el\'ements. Il r\'esulte du fait que $f$ est $0$-acyclique que
le morphisme canonique
$$
\H^0(S',C_{S'})\to \H^0(X',C_{X'})
$$
est bijectif, donc que $X'$ est connexe, d'o\`u la surjectivit\'e
de~$v$.

Par
d\'efinition de~$K^{\LL}$ \eqref{XIII.4.0}, on a la suite exacte
$$
1\to K^{\LL} \to \pi_1'(X,a)\to \pi_1(S,a)\to 1
$$
Soient $\widetilde{S}$ le rev\^etement universel de~$S$ et
$\widetilde{X}=\widetilde{S}\times_SX$; le groupe $K^{\LL}$ classe les
rev\^etements galoisiens $P$ de groupe un $\LL$-groupe, tels qu'il
existe un rev\^etement \'etale $S'$ de~$S$ et un rev\^etement
galoisien $Q$ de~$X'=X\times_SS'$ tels que l'on ait un isomorphisme
$P\simeq Q\times_{X'}\widetilde{X}$. Pour que la suite
\eqref{eq:XIII.4.1.1} soit exacte, il faut et il suffit que le
morphisme canonique
$$
\pi_1^{\LL}(X_{\bar{s}},a)\to K^{\LL}
$$
soit surjectif. D'apr\`es l'interpr\'etation de~$K^{\LL}$ cela
revient \`a dire que, pour tout rev\^etement \'etale $S'$ de~$S$
et pour tout rev\^etement galoisien $Q$ de~$X'$ de groupe un
$\LL$-groupe, tel que $P=Q\times_{X'}\widetilde{X}$ soit connexe,
alors $Q|X_{\bar{s}}$ est connexe. Montrons que cette derni\`ere
condition est satisfaite. Soient en effet $S'$ un rev\^etement
\'etale de~$S$, $Q$ un rev\^etement galoisien de~$X'$ de groupe un
$\LL$-groupe $F$, tel que $Q|X_{\bar{s}}$ soit disconnexe et montrons
que, quitte \`a remplacer $S'$ par un rev\^etement \'etale, $Q$
devient disconnexe. Il existe un sous-groupe $G$ de~$F$ distinct de
$F$ et un torseur $R$ sous $G|X_{\bar{s}}$ tel que $Q|X_{\bar{s}}$
s'obtienne par extension du groupe structural $G\to F$ \`a partir de
$R$. Le rev\^etement \'etale $E=Q/G$ de~$X'$ est tel que
$E|X_{\bar{s}}$ ait une section. D'apr\`es~\ref{XIII.1.16} $f_*E$
est localement constant constructible et, quitte \`a remplacer $S'$
par un rev\^etement \'etale, on peut m\^eme supposer que $f_*E$
est constant. Comme $(E,f')$ est cohomologiquement propre relativement
\`a~$S'$ en dimension $\leqslant0$ et comme $\H^0(X_{\bar{s}},
E|X_{\bar{s}})$ est non vide, on voit que $E$ a une section. Mais ceci
prouve que $Q$ est disconnexe, ce qui ach\`eve la d\'emonstration.

On d\'eduit de~\ref{XIII.4.1} le lemme suivant, qui sera utilis\'e
dans~\ref{XIII.4.6}.

\begin{lemme}
\label{XIII.4.2}
Soient $S$ un sch\'ema connexe, $f\colon X\to S$ un morphisme
localement $0$-acyclique et $0$-acyclique, $\LL$ un ensemble de
nombres premiers. Supposons
que, pour tout faisceau de~$\LL$-groupes constant fini $F$ sur~$X$,
$(F,f)$ soit cohomologiquement propre relativement \`a~$S$ en
dimension $\leqslant1$, et que, pour tout rev\^etement \'etale
$S'$ de~$S$ et pour tout rev\^etement \'etale $E$ de~$X'$,
quotient d'un rev\^etement galoisien de groupe un $\LL$-groupe,
$f'_*E$ soit constructible. Alors, si $\bar{s}$ est un point
g\'eom\'etrique de~$S$ et $a$ un point g\'eom\'etrique de la
fibre $X_{\bar{s}}$, la suite d'homomorphismes de groupes
$$
\pi_1^{\LL}(X_{\bar{s}},a)\to\pi_1'(X,a)\to\pi_1(S,a)\to1
$$
est exacte.
\end{lemme}

les hypoth\`eses de~\ref{XIII.4.1} sont satisfaites. Il r\'esulte
en effet de~\ref{XIII.1.13} 3) que, pour tout sch\'ema $S'$
\'etale sur~$S$ et pour tout rev\^etement \'etale $E$ de~$X'$,
quotient d'un rev\^etement galoisien de groupe un $\LL$-groupe,
$(E,f')$ est cohomologiquement propre relativement \`a~$S'$ en
dimension $\leqslant 0$.
\begin{proposition}
\label{XIII.4.3}
Soient $S$ un sch\'ema connexe, $\LL$ un ensemble de nombres
premiers, $f\colon X\to S$ un morphisme $0$-acyclique, localement
$1$-asph\'erique pour $\LL$ \textup{(SGA~4 XV~1.11)}, $g\colon S\to X$ une
section de~$f$. Soient $\bar{s}$ un point g\'eom\'etrique de~$S$,
$a$ un point g\'eom\'etrique de la fibre $X_{\bar{s}}$. On suppose
que, pour tout faisceau de
$\LL$-groupes
constant~$F$, $(F,f)$ est cohomologiquement propre en dimension
$\leqslant1$, que l'image directe par~$f$ du champ des torseurs sous
$F$ est un champ $1$-constructible \eqref{XIII.0} et que, pour tout
rev\^etement \'etale $E$ de~$X'$, quotient d'un rev\^etement
galoisien de groupe un $\LL$-groupe, $f'_*E$ est constructible. Alors
la suite d'homomorphismes de groupes
\begin{equation*}
\label{eq:XIII.4.3.1}
\tag{\thesubsection.1}
1\to\pi_1^{\LL}(X_{\bar{s}},a)\lto{u} \pi_1'(X,a)\lto{v}\pi_1(S,a)\to1
\end{equation*}
est exacte.
\end{proposition}

Compte tenu de~\ref{XIII.4.2}, il suffit de montrer l'injectivit\'e
du morphisme $u$, \ie de prouver que, pour tout rev\^etement
principal $\overline{Z}$ de~$X_{\bar{s}}$
de groupe un $\LL$-groupe~$C$, il existe un rev\^etement \'etale
$Z$ de~$X$ et un morphisme d'une composante connexe de~$Z|X_{\bar{s}}$
dans~$\overline{Z}$ (V~\ref{V.6.8}). Soient donc $\overline{Z}$ un
rev\^etement principal de~$X_{\bar{s}}$ de groupe un $\LL$-groupe~$C$ et~$\bar{z}$ sa classe dans
$\H^1(X_{\bar{s}},C_{X_{\bar{s}}})$. D'apr\`es \ref{XIII.1.5}~d), on
a un isomorphisme canonique
$$
\left(\R^1f_*C_X\right)_{\bar{s}}\isomto\H^1(X_{\bar{s}},C_{X_{\bar{s}}}),
$$
et d'apr\`es~\ref{XIII.1.16}, $\R^1f_*C_X$ est un faisceau
localement constant constructible. On peut donc trouver un
rev\^etement \'etale $S'$ de~$S$ tel que $\R^1f_*C_X|S'$ soit
constant. Si $\bar{s}\to S'$ est un point g\'eom\'etrique
au-dessus du point g\'eom\'etrique $\bar{s}\to S$, il existe un
\'el\'ement~$z$ de~$\H^0(S', \R^1f_*C_X)$ dont l'image dans
$\H^1(X_{\bar{s}},C_{X_{\bar{s}}})$ est $\bar{z}$. D'apr\`es le
lemme~\ref{XIII.4.3.1} ci-dessous, on peut trouver un rev\^etement
\'etale $S'_1$ de~$S'$ et un torseur $P$ sur~$X'_1$ de groupe~$C$
dont l'image dans $\H^0(S'_1,\R^1f_*C)$ soit \'egale \`a la
restriction de~$z$. Le torseur~$P$ est repr\'esentable par un
rev\^etement \'etale $Z$ de~$X'_1=X\times_SS'_1$ tel que
$Z\times_{X'_1}X_{\bar{s}}$ soit isomorphe \`a $\overline{Z}$. Si
l'on consid\`ere $Z$ comme un rev\^etement \'etale de~$X$, on a
alors un morphisme de~$Z\times_XX_{\bar{s}}$ dans $\overline{Z}$, ce
qui ach\`eve la d\'emonstration.

\begin{sublemme}
\label{XIII.4.3.1}
Soient $f\colon X\to S$ un morphisme $0$-acyclique et localement
$0$-acyclique, $g$ une section de~$f$. Soit $C$ un groupe constant
fini tel que $(C_X,f)$ soit cohomologiquement propre en dimension
$\leqslant0$ et que l'image directe par $f$ du champ des torseurs sous
$C_X$ soit constructible. Alors, pour toute section $z$ de
$\H^0(S,\R^1f_*C_X)$, on peut trouver un rev\^etement \'etale
$S_1$ de~$S$ et, si $X_1=X\times_SS_1$, un \'el\'ement de
$\H^1(X_1,C_{X_1})$, dont l'image par le morphisme canonique
$$
\H^1(X_1,C_{X_1})\to\H^0(S_1,\R^1f_*C_{X_1})
$$
soit \'egale \`a la restriction de~$z$ \`a
$\H^0(S_1,\R^1f_*C_{X_1})$.
\end{sublemme}

Pour tout sch\'ema $S'$ \'etale sur~$S$, on pose $X'=X\times_SS'$,
et on note $g'$ (\resp $F'$, etc.) l'image inverse de~$g$ (\resp $F$,
etc.) par le morphisme
$S'\to S$. Le pr\'efaisceau $G$ sur~$S$ d\'efini par
 

$$
G(S') = \left\{\text{classes mod. isomorphisme de
		torseurs $P$ sur~$X'$ de groupe $C_{X'}$, munis d'un
		isomorphisme $g'^*P\isomto C_{S'}$ } \right\}
$$
est alors un faisceau. Cela r\'esulte en effet par descente du fait
qu'un isomorphisme d'un torseur $P$ sur~$X'$ est bien
d\'etermin\'e par sa restriction \`a $g'(S')$. De plus, on a un
morphisme surjectif
$$
G\to \R^1f_*F.
$$
Soient $z$ un \'el\'ement de~$\H^0(S,\R^1f_*C_X)$ et $H$ le
sous-faisceau de~$G$ image r\'eciproque de~$z$. Il suffit de montrer
que $H$ est un faisceau localement constant constructible. Or, cette
propri\'et\'e \'etant locale sur~$S$, on peut supposer que $z$
provient d'un \'el\'ement de~$\H^1(X,C_X)$ repr\'esent\'e par
un torseur $P$ tel que $g^*P$ soit isomorphe \`a $C_S$. Se donner un
isomorphisme $i\colon g^*P\isomto C_S$ revient \`a se donner une
section globale de~$\SheafAut_{C_S}(g^*P)$ et deux isomorphismes $i$
et $i'$ d\'efinissent le m\^eme \'el\'ement de~$G(X)$ si et
seulement si $ii'^{-1}$ est l'image d'un \'el\'ement de
$\Aut_{C_X}(P)$. Si l'on consid\`ere l'injection canonique
$$
f_*\SheafAut_{C_X}(P)\to\SheafAut_{C_S}(g^*P)\simeq C,
$$
$H$ s'identifie donc au quotient de~$\SheafAut_{C_S}(g^*P)$ par
$f_*\SheafAut_{C_X}(P)$. D'apr\`es~\ref{XIII.1.16}
$f_*\SheafAut_{C_X}(P)$ est localement constant; il en est donc de
m\^eme de~$H$, ce qui ach\`eve la d\'emonstration.

\begin{exemples}
\label{XIII.4.4}
Notons que, si $S$ est un sch\'ema connexe, les hypoth\`eses
de~\ref{XIII.4.1} sont satisfaites lorsque le morphisme $f$ est propre
plat de pr\'esentation finie, \`a fibres g\'eom\'etriques
s\'eparables connexes, $\LL$ \'etant quelconque
(\cf X~\ref{X.1.3}). Les hypoth\`eses de~\ref{XIII.4.3} sont
satisfaites si de plus $f$ est lisse et a une section, si l'on
d\'esigne par $\LL$ l'ensemble des nombres premiers distincts des
caract\'eristiques r\'esiduelles de~$S$ (SGA~4 XV~2.1 et XVI~5.2).

Les hypoth\`eses de~\ref{XIII.4.1} sont aussi satisfaites si $S$ est connexe,
si l'on a un sch\'ema~$Z$ propre de pr\'esentation finie, plat sur~$S$, \`a
fibres g\'eom\'etriques s\'eparables connexes, tel que $X$ soit
le compl\'ementaire dans $Z$ d'un diviseur \`a croisements normaux
relativement \`a~$S$, $\LL$ \'etant l'ensemble des nombres
premiers distincts des caract\'eristiques r\'esiduelles de~$S$
\eqref{XIII.2.9}. Les hypoth\`eses de~\ref{XIII.4.3} sont
satisfaites si de plus $f$ est lisse et a une section.
\end{exemples}

\subsection{}
\label{XIII.4.5}
Reprenons les notations et les hypoth\`eses de~\ref{XIII.4.3}. Si
$\bar{s}$ est un point g\'eom\'etrique de~$S$ et $a=g(\bar{s})$,
la section $g$ permet de d\'efinir un morphisme
$$
w\colon \pi_1(S,a)\to \pi_1'(X,a),
$$
de sorte que $\pi_1'(X,a)$ s'identifie au produit semi-direct de
$\pi_1(S,a)$ par $\pi_1(X_{\bar{s}},a)$. Le groupe pro-fini
$\pi_1(S,a)$ op\`ere donc sur~$\pi_1(X_{\bar{s}}, a)$. Comme
$\pi_1^\LL(X_{\bar{s}},a)$ est limite projective stricte de groupes
invariants par l'action de~$\pi_1(S,a)$, la donn\'ee de
$\pi_1^\LL(X_{\bar{s}},a)$ muni de cette action est \'equivalente
\`a la donn\'ee d'un syst\`eme projectif strict de sch\'emas
en groupes \'etales finis sur~$S$ que l'on note
$$
\pi_1^\LL(X/S,g,\bar{s})\quad\text{ou
simplement}\quad\pi_1^\LL(X/S,g).
$$
\label{indnot:mk}
On a alors les propri\'et\'es suivantes:
\setcounter{subsubsection}{0}

\subsubsection{}
\label{XIII.4.5.1}
Pour tout sch\'ema en groupes \'etale fini $G$ sur~$S$, dont les
fibres sont des $\LL$-groupes, l'ensemble $E$ des classes de torseurs
$P$ sous l'image inverse $G_X$ de~$G$ sur~$X$, munis d'un isomorphisme
$g^*P\isomto G$, est canoniquement isomorphe \`a l'ensemble
$$
\Hom_S(\pi_1^\LL(X/S,g,\bar{s}),G) \quad\text{mod. automorphismes
intérieurs de~$G$.}
$$

\subsubsection{}
\label{XIII.4.5.2}
Pour tout sch\'ema en groupes \'etale fini $G$ sur~$S$ dont les
fibres sont des $\LL$-groupes, le faisceau $\R^1f_*G_X$ est
canoniquement isomorphe au faisceau associ\'e au pr\'efaisceau
$$
S'\mto \Hom_{S'}(\pi_1^\LL(X/S,g,\bar{s}),G)\
\text{mod. automorphismes intérieurs de~$G$.}
$$
($S'$ d\'esigne un sch\'ema \'etale sur~$S$).

\subsubsection{}
\label{XIII.4.5.3}
Soient
$S'$ un
$S$-sch\'ema
connexe, $\bar{s}$ un point g\'eom\'etrique de~$S'$, $X'$, $g'$
les images inverses respectives de~$X,g$ sur~$S'$. Alors
$\pi_1^\LL(X'/S',g',\bar{s})$ est canoniquement isomorphe \`a
l'image inverse de~$\pi_1^\LL(X/S,g,\bar{s})$ sur~$S'$. Pour tout
point g\'eom\'etrique $\xi$ de~$S$, la fibre
$\pi_1^\LL(X/S,g,\bar{s})_\xi$ est isomorphe \`a $\pi_1^\LL(X_\xi)$.

La donn\'ee de~$G$ est en effet \'equivalente \`a la donn\'ee
d'un $\LL$-groupe abstrait $\mathbf{G}$
sur lequel op\`ere $\pi_1(S,a)$, d'o\`u une action de
$\pi_1'(X,a)$ sur~$\mathbf{G}$. L'isomorphisme d\'efini
dans~\ref{XIII.4.5.1} s'obtient alors par restriction au sous-ensemble
$E$ \`a partir du morphisme canonique
\[
\H^1(\pi_1'(X,a),\mathbf{G})\to
\H^1(\pi_1^\LL(X_{\bar{s}},a),\mathbf{G}) =
\Hom(\pi_1^\LL(X_{\bar{s}},a),\mathbf{G})/\text{aut. int. }G,
\]
l'ensemble $E$ s'envoyant bijectivement sur le sous-ensemble des
morphismes de $\pi_1^\LL(X_{\bar{s}},a)$ dans $\mathbf{G}$ qui sont
compatibles avec l'action de
$\pi_1(S,a)$. L'assertion~\ref{XIII.4.5.3} r\'esulte alors de la
d\'efinition de~$\pi_1^\LL(X/S,g,\bar{s})$ compte-tenu de la suite
exacte d'homotopie \eqref{eq:XIII.4.3.1} et~\ref{XIII.4.5.2} se
d\'eduit de \ref{XIII.4.5.1} et~\ref{XIII.4.5.3}.

\begin{proposition}[Formule de K\"unneth]
\label{XIII.4.6}
\index{Kunneth (formule de)@K\"unneth (formule de)|hyperpage}Soient $k$ un corps s\'eparablement clos de caract\'eristique
$p\geq 0$, $X$ et $Y$ deux $k$-sch\'emas connexes, $a$ un point
g\'eom\'etrique de~$X$, $b$ un point g\'eom\'etrique de~$Y$,
$c$ un point g\'eom\'etrique de~$X\times_k Y$ au-dessus de~$a$ et
$b$. On suppose satisfaite l'une des deux conditions suivantes:
\begin{enumerate}
\item[a)] $X$ est de type fini sur~$k$ et les sch\'emas de type fini
sur une cl\^oture alg\'ebrique $\bar{k}$ de~$k$, de dimension
$\leq\dim X$, sont fortement d\'esingularisables. \textup{(SGA~5 I~3.1.5)}.
\item[b)] $X$ est quasi-compact et quasi-s\'epar\'e et tout
sch\'ema de type fini sur~$\bar{k}$ est fortement
d\'esingularisable.
\end{enumerate}

Alors, si $p'$ est l'ensemble des nombres premiers distincts de~$p$,
le morphisme
\begin{equation*}
\label{eq:XIII.4.6.0} \tag{\thesubsection.0} {}
\pi_1^{p'}(X\times_k Y,c)\to\pi_1^{p'}(X,a)\times\pi_1^{p'}(Y,b),
\end{equation*}
d\'eduit
des homomorphismes sur les groupes fondamentaux associ\'es aux
projections
$$
X\times_k Y\to X \quad \text{et} \quad X\times_k Y\to Y,
$$
est un isomorphisme.
\end{proposition}

On peut supposer $k$ alg\'ebriquement clos et $X$ r\'eduit
(SGA~4 VIII~1.1). Soient $Z=X\times_k Y$, $g\colon X\to k$ et $f\colon
Z\to Y$ les morphismes canoniques. Le morphisme $g$, donc aussi $f$,
est universellement localement $1$-asph\'erique pour $p'$
d'apr\`es~\ref{XIII.3.5}. Comme $X$ est connexe, $f$ est
$0$-acyclique (SGA~4 XV~1.16). D'autre part il r\'esulte
de~\ref{XIII.3.2} que, pour tout $p'$-groupe fini $C$, $(C_X,g)$ est
cohomologiquement propre relativement \`a $k$ en dimension $\leq
1$. Il en r\'esulte que $(C_Z,f)$ est cohomologiquement propre
relativement \`a $Y$ en dimension $\leq 1$ (\ref{XIII.1.5}~c)) et
que $f_*C_Z$ et $\R^1f_*C_Z$ sont des faisceaux constants. Par suite
$g$ satisfait \`a toutes les hypoth\`eses de~\ref{XIII.4.2}. On a
donc la suite exacte
$$
\pi_1^{p'}(X_b,c)\to\pi_1^{p'}(Z,c)\to\pi_1^{p'}(Y,b)\to 1.
$$
De plus le morphisme compos\'e
$$
\pi_1^{p'}(X_b,c)\to\pi_1^{p'}(Z,c)\to\pi_1^{p'}(Y,b)
$$
est un isomorphisme et l'on a donc la suite exacte
$$
1\to\pi_1^{p'}(X,a)\to\pi_1^{p'}(Z,c)\to\pi_1^{p'}(Y,b)\to 1.
$$
D'autre part le morphisme \eqref{eq:XIII.4.6.0} permet de d\'efinir
un morphisme de cette suite exacte dans la suite exacte
$$
1\to\pi_1^{p'}(X,a)\to
\pi_1^{p'}(X,a)\times\pi_1^{p'}(Y,b)\times\pi_1^{p'}(Y,b)\to 1,
$$
et il en r\'esulte que le morphisme \eqref{eq:XIII.4.6.0} est un
isomorphisme.

\subsection{}
\label{XIII.4.7} Soit
$$
\xymatrix@C=1.2cm{{X}\ar[d]_-f&{X_U}\ar[l]\ar[d]_-{f_U}&{X_{\bar{s}}}\ar[l]\ar[d]\\{S}&{U}\ar[l]&{\bar{s}}\ar[l]}
$$
un
diagramme dont les carr\'es sont cart\'esiens, o\`u $S$ est un
sch\'ema connexe par arcs (SGA~4 IX~2.12), $U$ un ouvert connexe de
$S$, $\bar{s}$ un point g\'eom\'etrique de~$U$. Soient $a$ un
point g\'eom\'etrique de~$X_{\bar{s}}$, $\LL$ un ensemble de
nombres premiers. Soit $g$ une section de~$f$ et supposons que les
conditions suivantes soient satisfaites:
\begin{enumerate}
\item[a)] Le morphisme $f$ est $0$-acyclique, localement $0$-acyclique,
et, pour tout rev\^etement \'etale $S'$ de~$S$ et tout
rev\^etement \'etale $E$ de~$X\times_S S'$ quotient d'un
rev\^etement galoisien de groupe un $\LL$-groupe, $(E,f_{(S')})$ est
cohomologiquement propre relativement \`a~$S'$ en dimension $\leq
0$.
\item[b)] Le morphisme $f_U$ est localement $1$-asph\'erique pour
$\LL$, et, pour tout faisceau de~$\LL$-groupe constant fini $F$ sur~$X_U$, $(F,f_U)$ est cohomologiquement propre en dimension $\leq 1$ et
les fibres de~$\R^1 f_{U*}F$ sont finies.
\end{enumerate}

On d\'eduit alors de~\ref{XIII.4.1} et~\ref{XIII.4.3} le diagramme
commutatif suivant dont les lignes sont exactes:
\begin{equation*}
\label{eq:XIII.4.7.0} \tag{\thesubsection.0} {}
\begin{array}{c}
\xymatrix{{1}\ar[r]&{\pi_1^{\LL}(X_{\bar{s}},a)}\ar[r]\ar@{=}[d]&{\pi'_1(X_U,a)}\ar[r]\ar[d]&{\pi_1(U,a)}\ar[r]\ar[d]&{1}\\
&{\pi_1^{\LL}(X_{\bar{s}},a)}\ar[r]&{\pi'_1(X,a)}\ar[r]&{\pi_1(S,a)}\ar[r]&{1}}
\end{array}
\end{equation*}

Gr\^ace \`a la section $g$, on a des morphismes
$$\pi_1(U,a)\to\pi'_1(X_U,a),\quad \pi_1(S,a)\to\pi'_1(X,a)\,\;$$
on en
d\'eduit un morphisme de la somme amalgam\'ee de~$\pi_1(S,a)$ et
$\pi'_1(X_U,a)$ au-dessus de~$\pi_1(U,a)$ dans $\pi'_1(X,a)$:
\begin{equation*}
\label{eq:XIII.4.7.1} \tag{\thesubsection.1} {}
\varphi\colon\pi=\pi_1(S,a)\coprod_{\pi_1(U,a)}\pi'_1(X_U,a)\to\pi'_1(X,a).
\end{equation*}
Supposons satisfaite la condition suivante:
\begin{enumerate}
\item[c)] Si $T=S-U$, on a $\prof \et_T(S)\geq 2$ (SGA~2 XIV~1.1).
\end{enumerate}
Alors le foncteur qui, \`a un rev\^etement \'etale de~$S$, fait
correspondre sa restriction \`a $U$ est pleinement fid\`ele
(SGA~2 XVI~1.4). Il en r\'esulte que le morphisme
$\pi_1(U,a)\to\pi_1(S,a)$ est surjectif (V~\ref{V.6.9}) et l'on
d\'eduit du
diagramme \eqref{eq:XIII.4.7.0} qu'il en est de m\^eme du morphisme
$\pi'_1(X_U,a)\to\pi'_1(X,a)$; a fortiori $\varphi$ est un
\'epimorphisme. Soit
\begin{equation*}
\label{eq:XIII.4.7.2} \tag{\thesubsection.2} {}
K=\Ker(\pi_1(U,a)\to\pi_1(S,a)).
\end{equation*}
Le groupe $\pi$ de \eqref{eq:XIII.4.7.1} s'identifie au quotient de
$\pi'_1(X_U,a)$ par le sous-groupe invariant ferm\'e engendr\'e
par l'image $L$ de~$K$ dans $\pi'_1(X_U,a)$. Consid\'erons
$\pi'_1(X_U,a)$ comme produit semi-direct de~$\pi_1(U,a)$ par
$\pi_1^\LL (X_{\bar{s}},a)$. Le groupe $K$ op\`ere alors par
automorphismes int\'erieurs sur~$\pi_1^\LL (X_{\bar{s}},a)$, et le
quotient $\pi=\pi'_1(X_U,a)/L$ s'identifie au produit semi-direct de
$\pi_1(U,a)/K=\pi_1(S,a)$ par le groupe $\pi_1^\LL(X_{\bar{s}},a)_K$
des coinvariants de~$\pi_1^{\LL}(X_{\bar{s}},a)$ sous~$K$. On a
finalement un \'epimorphisme
\begin{equation*}
\label{eq:XIII.4.7.3} \tag{\thesubsection.3} {}
\varphi\colon\pi=\pi_1^\LL(X_{\bar{s}},a)_K\cdot\pi_1(S,a)\to\pi'_1(X,a).
\end{equation*}
\label{indnot:ml}
La proposition qui suit donne des conditions sous-lesquelles le
morphisme $\varphi$ est un isomorphisme.

\setcounter{subsubsection}{3}

\begin{subproposition}
\label{XIII.4.7.4} Les notations sont celles de~\ref{XIII.4.7}. On
suppose que, en plus des conditions \textup{a), b), c)} les conditions
suivantes soient satisfaites:
\begin{enumerate}
\item[d)] Pour tout point $t$ de~$T=S-U$, le morphisme $f$ est
localement $1$-asph\'erique pour $\LL$ en $g(t)$.
\item[e)] Pour tout point $t\in T$, toute composante irr\'eductible
de la fibre $X_t$ contient $g(t)$ et, pour tout point $x$ de
$X_t-\{g(t)\}$ qui n'est pas maximal, on a
$$
\prof\hop_x(X)\geq 3 \quad\textup{(SGA~2 XIV~1.2)}
$$
et l'anneau $\cal{O}_{X,x}$ est noeth\'erien.
\end{enumerate}

Alors le morphisme \eqref{eq:XIII.4.7.3} est un isomorphisme.
\end{subproposition}

Comme on l'a dit pr\'ec\'edemment, le groupe $\pi$ s'identifie au
quotient de~$\pi'_1(X_U,a)$ par le sous-groupe invariant ferm\'e $L$
engendr\'e par l'image de~$K$ \eqref{eq:XIII.4.7.2} dans
$\pi'_1(X_U,a)$. Cela revient \`a dire que $\pi$ classe les
rev\^etements principaux $Z$ de~$X_U$ tels que $g^{-1}_U(Z)$ se
prolonge en un rev\^etement \'etale de~$S$, et qui induisent sur~$X_{\bar{s}}$ un rev\^etement qui s'obtient par extension du groupe
structural \`a partir d'un rev\^etement principal de groupe un
$\LL$-groupe. Pour prouver que $\varphi$ est un isomorphisme, il
suffit de montrer qu'un tel rev\^etement $Z$ se prolonge \`a $X$
tout entier.

Montrons d'abord que $Z$ se prolonge en un rev\^etement \'etale
d'un ouvert contenant $X_U$ et $g(S)$. Soit $W$ un sch\'ema
\'etale sur~$X$ dont l'image contient $X_U$ et $g(S)$, et posons
$W_U=W\times_S U$. Du fait que le morphisme $W\to S$ est $0$-acyclique
et de la relation $\mathrm{prof}.\,\et_T S\geq 2$, r\'esulte que l'on a
$$
\mathrm{prof}.\,\et_{(W-W_U)} W\geq 2
$$
(SGA~2 XIV~1.13); par suite, si $Z|W_U$ se prolonge en un
rev\^etement \'etale de~$W$, ce prolongement est unique \`a
isomorphisme unique pr\`es. Il en r\'esulte que le probl\`eme de
prolonger $Z$ \`a un voisinage de~$g(S)\cup X_U$ est local pour la
topologie \'etale au voisinage des points de~$g(T)$. Si $t$ est un
point de~$T$, on pose $x=g(t)$, et on note $\bar{X}$ (\resp $\bar{S}$)
le localis\'e strict de~$X$ en $\bar{x}$ (\resp de~$S$ en
$\bar{t}$), $\bar{U}=U\times_S \bar{S}$, $\bar{X}_U=\bar{X}\times_X
X_U$, $\bar{g}\colon\bar{S}\to\bar{X}$ le morphisme d\'eduit de
$g$. Il suffit de montrer que, pour tout point $t$ de~$T$, l'image
inverse $\bar{Z}$ de~$Z$ sur~$\bar{X}_U$ se prolonge \`a $\bar{X}$
ou, ce qui revient au m\^eme, est triviale. Or, par d\'efinition
de~$Z$, l'image inverse de~$\bar{Z}$ sur~$\bar{U}$ est
triviale. Pour prouver que $\bar{Z}$ est trivial, il suffit de montrer
qu'il est de la forme $\bar{f}^*_U E$, o\`u $E$ est un
rev\^etement principal de~$\bar{U}$; on aura en effet
$E\simeq\bar{g}^*_U\bar{f}^*_U E \simeq\bar{g}^*_U\bar{Z}$, d'o\`u
le r\'esultat puisque $\bar{g}^*_U\bar{Z}$ est trivial. Or le
morphisme $\bar{f}_U$ \'etant $0$-acyclique et localement $0$-acyclique,
il suffit, pour prouver que $\bar{Z}$ provient de~$\bar{U}$, de
montrer que, pour tout point g\'eom\'etrique alg\'ebrique sur un
point de~$\bar{U}$, que l'on peut supposer \^etre le point
$\bar{s}$, $\bar{Z}|X_{\bar{s}}$ est trivial (SGA~4 XV~1.15). Mais
$\bar{Z}|X_{\bar{s}}$ \'etant obtenu par extension du groupe
structural \`a partir d'un rev\^etement principal de groupe un
$\LL$-groupe, cela r\'esulte du fait que le morphisme $\bar{f}$ est
$1$-asph\'erique pour $\LL$.

On
a donc d\'emontr\'e qu'il existe un voisinage ouvert $V$ de
$g(S)\cup X_U$ tel que $Z$ se prolonge en un rev\^etement \'etale
$Z_V$ de~$V$. Montrons que $Z_V$ se prolonge \`a $X$ tout entier.
Il suffit de voir que, pour tout point $x$ de~$X-V$, on a
$$
\prof\hop_x X\geq 3.
$$
Or cela r\'esulte de l'hypoth\`ese e) et du fait qu'un point $x$
de~$X-V$ ne peut \^etre maximal dans sa fibre $X_t$ car, toute
composante irr\'eductible de~$X_t$ contenant $g(t)$, tout point
maximal de~$X_t$ appartient \`a $V$.

\begin{corollaire} \label{XIII.4.8}
Les hypoth\`eses sont celles de~\ref{XIII.4.7.4} mais on suppose de
plus que l'on a $\pi_1(S,a)=1$. Alors on a un isomorphisme
\begin{equation*}
\label{eq:XIII.4.8.*} \tag{$*$}
{}\pi_1^\LL(X,a)\isomto\pi_1^\LL(X_{\bar{s}},a)_K.
\end{equation*}
En particulier, si $\pi_1^\LL(X_{\bar{s}},a)$ est topologiquement de
pr\'esentation finie et si $K$ op\`ere sur~$\pi_1^\LL(X_{\bar{s}},a)$ par l'interm\'ediaire d'un groupe de type
fini, alors $\pi_1^\LL(X,a)$ est topologiquement de pr\'esentation
finie.
\end{corollaire}

L'isomorphisme \eqref{eq:XIII.4.8.*} a \'et\'e d\'emontr\'e
dans~\ref{XIII.4.7.4}. Supposons $\pi_1^\LL(X_{\bar{s}},a)$ quotient
du pro-$\LL$-groupe libre \`a $n$ g\'en\'erateurs
$L(x_1,\dots,x_n)$ par le sous-groupe invariant ferm\'e engendr\'e
par les \'el\'ements $y_1,\dots,y_p$ de~$L(x_1,\dots,x_n)$, et
supposons que $K$ agisse par l'interm\'ediaire d'un groupe
engendr\'e par des \'el\'ements $k_1,\dots,k_q$. Si pour tout
$i\in[1,n]$, $j\in[1,q]$, on note $z_{ij}$ un \'el\'ement de
$L(x_1,\dots,x_n)$ relevant l'\'el\'ement $(k_j\cdot
x_i)x_i^{-1}$, alors $\pi_1^\LL(X_{\bar{s}},a)_K$ est quotient de
$L(x_1,\dots,x_n)$ par le sous-groupe invariant ferm\'e engendr\'e
par les \'el\'ements $(y_i)_{i\in[1,p]}$,
$(z_{ij})_{i\in[1,n],j\in[1,q]}$.



\begin{remarques}
\label{rem:XIII.4.6}
a) Les conditions a) \`a e) de~\ref{XIII.4.7} sont
satisfaites lorsque $S$ est un sch\'ema normal connexe, $U$ un
ouvert dense r\'etrocompact de~$S$, $f$ un morphisme propre de
pr\'esentation finie, \`a fibres g\'eom\'etriquement connexes
et irr\'eductibles en tout point $t$ de~$T$, $f$ \'etant de plus
s\'eparable,
lisse aux points de~$X_U\cup g(T)$, $\LL$ \'etant l'ensemble des
nombres premiers distincts des caract\'eristiques r\'esiduelles de
$S$, et $X$ \'etant r\'egulier en tout point de~$X_t$. La
condition a) r\'esulte en effet de SGA~4 XV~4.1 et 1.4; les
conditions b) et d) r\'esultent de~\ref{XIII.1.4} et SGA~4 XV~2.1 et
XVI~5.2. Enfin e) r\'esulte de SGA~XIV~1.11.
\par\smallskip
b) Le corollaire~\ref{XIII.4.5} s'applique pour calculer le
groupe fondamental $\pi_1^{p'}(X)$ d'une surface $X$ propre et lisse
sur un corps s\'eparablement clos $k$ de caract\'eristique $p$
($p'$~d\'esignant l'ensemble des nombres premiers distincts de
$p$). La m\'ethode nous a \'et\'e communiqu\'ee par
J.P.\, \textsc{Murre}; elle consiste \`a se ramener, en faisant \'eclater
$X$, au cas o\`u l'on a une fibration $X\to\PP^1_k$ et un ouvert $U$
de~$\PP^1_k$ satisfaisant aux hypoth\`eses de~\ref{XIII.4.7} (voir
SGA~7 pour plus de d\'etails). La m\^eme m\'ethode peut \^etre
utilis\'ee plus g\'en\'eralement (\loccit) pour prouver que, si
$X$ est un $k$-sch\'ema connexe de type fini, et, si les sch\'emas
de type fini de dimension $\leq\dim X$ sur une cl\^oture
alg\'ebrique de~$k$ sont fortement d\'esingularisables
(SGA~5 I~3.1.5), alors $\pi_1^{p'}(X)$ est topologiquement de
pr\'esentation finie.
\end{remarques}

\section{Appendice I: Variations sur le lemme d'Abhyankar}
\label{XIII.5}
\index{lemme d'Abhyankar|hyperpage}Cet appendice contient diff\'erentes variantes du lemme d'Abhyankar.

\begin{proposition}
\label{XIII.5.1} Soient $X=\Spec A$ un sch\'ema local r\'egulier,
$D=\sum_{1\leq i\leq r} \divisor f_i$ un diviseur \`a croisements
normaux, o\`u les $f_i$ sont des \'el\'ements de l'id\'eal
maximal de~$A$ faisant partie d'un syst\`eme r\'egulier de
param\`etres. Soient $n_i$, $1\leq i\leq r$, des entiers $\geq 0$
et posons
$$
X'=X[T_1,\dots,T_r]/(T_1^{n_1}-f_1,\dots,T_r^{n_r}-f_r),
$$
$U'=U\times_X X'$. Alors $X'$ est r\'egulier et $U'$ est le
compl\'ementaire dans $X'$ du diviseur \`a croisements normaux
$\sum_{1\leq i\leq r}\divisor T_i$. Si les entiers $n_i$ sont premiers
\`a la caract\'eristique r\'esiduelle $p$ de~$X$, $U'$ est un
rev\^etement \'etale connexe de~$U$, mod\'er\'ement
ramifi\'e relativement \`a $D$ \textup{(\ref{XIII.2.3}~c))}.
\end{proposition}

En effet $X'$ est le spectre d'un anneau local $A'$ dont l'id\'eal
maximal est engendr\'e par $T_1,\dots,T_r$.
(On peut supposer $r=\dim(A)$.)
Comme $A'$ est fini et plat sur~$A$, donc de dimension $r$, $A'$ est
r\'egulier (EGA~$0_{\textup{IV}}$~17.1.1) et les $T_i$ forment un syst\`eme
r\'egulier de param\`etres de~$A'$. Supposons les $n_i$ premiers
\`a $p$. Comme tous les $f_i$ sont inversibles sur~$U$, le fait que
$U'$ soit \'etale sur~$U$ r\'esulte de I~\ref{I.7.4}. De plus $U'$
est mod\'er\'ement ramifi\'e relativement \`a $D$; soit en
effet $x_i$ le point g\'en\'erique de~$V(f_i)$, $\bar{R}$ le
localis\'e strict de~$R=\cal{O}_{X,x_i}$, $\bar{K}$ le corps des
fractions de~$\bar{R}$. Alors la $\bar{K}$-alg\`ebre qui
repr\'esente $U'|\bar{K}$ s'obtient \`a partir du corps
$\bar{K}[T_i]/(T_i^{n_i}-f_i)$ en faisant une extension non
ramifi\'ee; elle est donc mod\'er\'ement ramifi\'ee sur~$\bar{R}$.

\begin{proposition}[Lemme d'Abhyankar absolu]
\label{XIII.5.2}
\index{Abhyankar absolu (lemme d')|hyperpage}Soit $X$ un sch\'ema local r\'egulier,
$$
D=\sum_{1\leq i\leq r} \divisor f_i
$$
un diviseur \`a croisements normaux comme dans~\ref{XIII.5.1},
$Y=\Supp D$, $U=X-Y$. Soit~$V$ un rev\^etement \'etale de~$U$
mod\'er\'ement ramifi\'e relativement \`a $D$. Si $x_i$ est le
point g\'en\'erique du ferm\'e $V(f_i)$, $\cal{O}_{X,x_i}$ est
un anneau de valuation discr\`ete de corps des fractions $K_i$, et
on a $V|K_i=\Spec(\prod_{j\in J_i}L_j)$, o\`u les $L_j$ sont des
extensions finies s\'eparables de~$K_i$; notons $n_j$ l'ordre du
groupe d'inertie d'une extension galoisienne engendr\'ee par $L_j$
et $n_i$ le p.p.c.m. des $n_j$ quand $j$ parcourt $J_i$. Si l'on pose
$$
X'=X[T_1,\dots,T_r]/(T_1^{n_1}-f_1,\dots,T_r^{n_r}-f_r),
$$
$U'=U_{(X')}$, $V'=V_{(X')}$, etc., le rev\^etement \'etale $V'$
de~$U'$ se prolonge de mani\`ere unique \`a isomorphisme unique
pr\`es en un rev\^etement \'etale de~$X'$, et les $n_i$ sont
premiers \`a la caract\'eristique r\'esiduelle $p$ de~$X$.
\end{proposition}

L'unicit\'e r\'esulte du fait que $X'$ est normal
\eqref{XIII.5.1}; en effet
un rev\^etement \'etale de~$X'$ qui prolonge $V'$ est isomorphe au
normalis\'e de~$X'$ dans la fibre de~$V'$ au point g\'en\'erique
de~$X'$ \eqref{I.10.2}. Si $\bar{x}'$ est un point g\'eom\'etrique
de~$Y'$, on note $\bar{X}'$ le localis\'e strict de~$X'$ en
$\bar{x}'$, $\bar{V}'=V'_{(\bar{X}')}$, etc. Par descente, compte tenu
de l'unicit\'e, il suffit de montrer que, pour tout point
g\'eom\'etrique $\bar{x}'$ de~$Y'$, le rev\^etement \'etale
$\bar{V}'$ de~$\bar{U}'$ se prolonge \`a $\bar{X}'$. \'Etant donn\'e
qu'un rev\^etement \'etale d'un ouvert du sch\'ema r\'egulier~$X'$ qui contient tous les points $x'$ tels que l'on ait $\dim
\cal{O}_{X',x'}\leq 1$ se prolonge \`a $X$ tout entier
(SGA~2 XIV~1.11), on peut m\^eme se borner aux points $\bar{x}'$ qui
se projettent sur un point maximal de~$Y'$. Or, en un tel point
$\bar{x}'$, le fait que $\bar{V}'$ se prolonge en un rev\^etement
\'etale de~$\bar{X}'$ r\'esulte de~\ref{X.3.6}.

Montrons que les $n_i$ sont premiers \`a $p$. En effet, s'il n'en
\'etait pas ainsi, on aurait par exemple $p|n_1$. Quitte \`a
remplacer $X$ par
$$X[T_1,\dots,T_r]/(T_1^{n_1/p}-f_1,T_2^{n_2}-f_2,\dots,T_r^{n_r}-f_r),$$
on se ram\`ene au cas o\`u l'on~a $X'=X[T_1]/(T_1^p-f_1)$. Il suffit
de montrer que $V$ se prolonge en un rev\^etement \'etale de~$X$,
car on aura alors $n_1=1$ contrairement \`a l'hypoth\`ese. On peut
supposer pour cela que $X$ est strictement local. Soit $Z$ le
sous-sch\'ema ferm\'e de~$X$ d'\'equation $p=0$ et $Z_1=Z\cap
X_{f_1}$; $Z_1$ est un ouvert non vide de~$Z$. D'apr\`es ce qui
pr\'ec\`ede le rev\^etement \'etale $V'$ de~$U'$ se prolonge
en un rev\^etement \'etale $W'$ de~$X'$. Soient $W''_1$ et $W''_2$
les images inverses de~$W'$ par les deux projections $X''=X'\times_X
X'\rightrightarrows X'$, et montrons que l'isomorphisme de descente
$u\colon W''_1|U''\to W''_2|U''$ se prolonge en un $X''$-morphisme
$W''_1\to W''_2$ qui sera n\'ecessairement une donn\'ee de
descente sur~$W'$ relativement \`a $X'\to X$. Soit $Z''$
(\resp $Z''_1$) l'image inverse de~$Z$ (\resp $Z_1$) dans $X''$. Comme
le morphisme $Z''\to Z$ est radiciel, il existe un isomorphisme
$v\colon W''_1|Z''\to W''_2|Z''$ qui prolonge l'isomorphisme
$u|Z''_1$. Mais, comme $X$ est hens\'elien, on a une bijection
$$
\Hom_{X''}(W''_1,W''_2)\simeq\Hom_{Z''}(W''_1|Z'',W_2''|Z''),
$$
d'o\`u un morphisme $w\colon W''_1\to W''_2$ relevant $v$. Le
sous-sch\'ema \`a la fois ouvert
et ferm\'e de~$X''$ au-dessus duquel $u$ et $w$ co\"incident
contient $Z''_1$, donc \'egal \`a~$X''$, d'o\`u le fait que $V$
se prolonge \`a~$X$.

\setcounter{subsection}{3} \setcounter{subsubsection}{-1}

\subsubsection{}
\label{XIII.5.3.0} Reprenons les hypoth\`eses et les notations
de~\ref{XIII.5.2} en supposant de plus
$X$
strictement local. Alors
il r\'esulte de \loccit que tout rev\^etement \'etale connexe
de~$U$ mod\'er\'ement ramifi\'e relativement \`a $D$ est
quotient d'un rev\^etement mod\'er\'ement ramifi\'e de la
forme
$$
U'=U[T_1,\dots,T_r]/(T_1^{n_1}-f_1,\dots,T_r^{n_r}-f_r),
$$
o\`u les $n_i$ sont des entiers premiers \`a $p$. Soit
$\bbmu_n$ le groupe des racines $n$-i\`emes de l'unit\'e de
$U$. Le groupe des $U$-automorphismes de~$U'$ n'est autre que le
groupe $\bbmu_{n_1}\times\dots\times\bbmu_{n_r}$, une racine
$n_i$-i\`eme de l'unit\'e $\xi_i$ op\'erant sur~$U'$ en
transformant $T_i$ en $\xi_i T_i$. On a donc l'\'enonc\'e suivant:

\setcounter{subsection}{2}

\begin{corollaire}
\label{XIII.5.3} Soit $X$ un sch\'ema strictement local r\'egulier
de caract\'eristique r\'esiduelle $p\geq 0$, $D=\sum_{1\leq i\leq
n} \divisor f_i$ un diviseur \`a croisements normaux sur~$X$,
$U=X-\Supp D$. Posons
$$
\tilde U=\varprojlim_{(n_i)}
U[T_1,\dots,T_r]/(T_1^{n_1}-f_1,\dots,T_r^{n_r}-f_r),
$$
la limite projective \'etant prise suivant l'ensemble ordonn\'e
filtrant (pour la relation de divisibilit\'e) des familles d'entiers
$n_i>0$, premiers \`a $p$. Alors $\tilde U$ est un rev\^etement
universel mod\'er\'ement ramifi\'e de~$U$. Par suite le groupe
fondamental mod\'er\'ement ramifi\'e de~$U$ est
$$
\pi_1^{\tame}(U)\simeq \prod_{\ell\neq p} \ZZ_\ell [1]^r
\qquad \text{(isomorphisme canonique)}
$$
o\`u l'on a pos\'e
$\ZZ_\ell[1]=\varprojlim_{n>0}\bbmu_{\ell^n}$. Le groupe
$\pi_1^{\tame}(U)$ est non canoniquement isomorphe \`a $\prod_{\ell\neq p}
\ZZ_\ell^r$.
\label{indnot:mm}\end{corollaire}

\begin{proposition}
\label{XIII.5.4} Soient $f\colon X\to S$ un morphisme de
sch\'emas, $D=\sum_{1\leq i\leq r} \divisor f_i$ un diviseur \`a
croisements normaux relativement \`a~$S$ \eqref{XIII.2.1}, o\`u,
pour chaque point $x$ de~$Y=\Supp D$, si $I(x)\subset[1,r]$ est
l'ensemble des $i$ tels que l'on ait $f_i(x)=0$, le sous-sch\'ema
$V((f_i)_{i\in I(x)})$ est lisse sur~$S$ de codimension
$\card I(x)$
dans $X$. Soit $U=X-Y$. Soient $x$ un point de
$Y$, $X_1=\Spec\cal{O}_{X,x}$, $U_1=U\times_X X_1$, $n_i$, $i\in I(x)$
des entiers et
$$
X'=X[T_i]_{i\in I(x)}/(T_i^{n_i}-f_i).
$$
Alors, si $x'$ est le point de~$X'$ au-dessus de~$x$, $X'$ est lisse
sur~$S$ en $x'$. Si les entiers~$n_i$ sont premiers \`a la
caract\'eristique $p$ de~$\kres(x)$, $U'_1=U_1\times_X X'$ est un
rev\^etement \'etale connexe de~$U_1$ mod\'er\'ement
ramifi\'e sur~$X_1$ relativement \`a~$S_1$ \eqref{XIII.2.1.1}.
\end{proposition}

Si $s=f(x)$, la fibre g\'eom\'etrique $X'_{\bar{s}}$ est
r\'eguli\`ere au point $x'$ \eqref{XIII.5.1}; comme $X'$ est plat
sur~$S$ au voisinage de~$Y$, cela prouve que $X'$ est lisse sur~$S$ en
$x'$ (EGA~IV~12.1.6). Si les entiers $n_i$ sont premiers \`a $p$,
$U'_1$ est un rev\^etement \'etale de~$U_1$ \eqref{I.7.4}; il est
mod\'er\'ement ramifi\'e sur~$X_1$ relativement \`a~$S$ car il
en est ainsi sur les fibres g\'eom\'etriques en chaque point de~$S$
\eqref{XIII.5.1}. Enfin le fait que $U'_1$ soit connexe r\'esulte de
(SGA~4 XVI~3.2).

\begin{proposition}[Lemme d'Abhyankar relatif]
\label{XIII.5.5}
\index{Abhyankar relatif (lemme d')|hyperpage}Soient $X$ un $S$-sch\'ema, $D$ un
diviseur \`a croisements normaux relativement \`a~$S$, comme
dans~\eqref{XIII.5.4}. Soient $Y=X-\Supp D$,
$U=X- Y$, $x$ un point de~$Y$, $X_{1}$ le localis\'e
strict de~$X$ en un point g\'eom\'etrique au-dessus de~$x$,
$U_{1}=U\times_{X}X_{1}$, $V_{1}$ un rev\^etement \'etale de
$U_{1}$. On suppose que, pour tout point maximal $s$ de~$S$,
$V_{1\overline{s}}$ est mod\'er\'ement ramifi\'e sur~$X_{1\overline{s}}$ relativement \`a $\overline{s}$. Alors on peut
trouver des entiers $n_{i}$ premiers \`a la caract\'eristique $p$
de~$\kres(x)$, avec $i\in I(x)$, tels que, si l'on pose
$$
X'_{1}=X_{1}[T_{i}]_{i\in I(x)}\big/ (T_{i}^{n_{i}}-f_{i}) \quoi,
$$
$U'_{1}=U_{1}\times_{X_{1}}X'_{1}$, etc., le rev\^etement \'etale
$V'_{1}$ de~$U'_{1}$ se prolonge de mani\`ere unique \`a
isomorphisme unique pr\`es en un rev\^etement \'etale de
$X'_{1}$. En particulier
$V_{1}$ est mod\'er\'ement ramifi\'e sur~$X_{1}$ relativement
\`a~$S$.
\end{proposition}

On peut supposer $S$ local noeth\'erien de point ferm\'e~$f(x)$.
Pour chaque point maximal $s$ de~$S$ et pour chaque $i\in I(x)$, soit
$x_{i}$ le point g\'en\'erique du ferm\'e $V(f_{i})$ de la fibre
$X_{1\overline{s}}$. L'anneau local $(\cal{O}_{X_{1},x_{i}})_{\red}$
est un anneau de valuation discr\`ete de corps de fractions $K_{i}$
et l'on a
$V_{1}|K_{i}=\Spec(\prod_{j\in I(x_{i})}L_{j})$,
o\`u $L_{j}$ est une extension finie s\'eparable de~$K$; on note
$n_{j}$ l'ordre du groupe d'inertie d'une extension galoisienne
engendr\'ee par $L_{j}$ et $n_{i}$ le p.p.c.m. des $n_{j}$ quand
$s$ parcourt les points maximaux de~$S$ et
$j\in I(x_{i})$.

Les $n_{i}$ \'etant ainsi choisis, nous allons montrer que $V'_{1}$
se prolonge de fa\c con unique en un rev\^etement \'etale de
$X'_{1}$. L'unicit\'e r\'esulte du fait que, $X'$ \'etant lisse
sur~$S$ aux points de~$Y$, on a $\prof\et_{Y'_{1}}(X'_{1})\geqslant 2$
(SGA~4 XVI~3.2 ou SGA~2 XIV~1.19). Soient $x'_{1}$ un point de
$Y'_{1}$, $\overline{x}'_{1}$ un point g\'eom\'etrique au-dessus
de~$x'_{1}$, et notons $\overline{X'_{1}}$ le localis\'e strict de
$X'_{1}$ en $\overline{x}'_{1}$,
$\overline{U'_{1}}=U'_{1(\overline{X'_{1}})}$,~etc. Par descente,
compte tenu de l'unicit\'e, il suffit de montrer que
$\overline{V}'_{1}$ se prolonge \`a $\overline{X'_{1}}$. De plus on
peut se borner \`a prendre pour $x'_{1}$ les point maximaux de
$Y'_{1}$; en effet on aura alors un prolongement de~$V'_{1}$ sur un
ouvert $W'_{1}$ de~$X'_{1}$ contenant les points maximaux de~$Y'_{1}$;
or, si $Z'_{1}=X'_{1}- W'_{1}$, on a
$\codim(Z'_{1s},X'_{1s})\geqslant 2$ si $s$ est un point maximal de
$S$ et $\codim(Z'_{1s},X'_{1s})\geqslant 1$, $\prof\et_{s}(S)\geqslant
1$ si $s$ est un point de~$S$ qui n'est pas maximal; le fait que
$V'_{1}$ se prolonge \`a $X'_{1}$ tout entier r\'esulte alors
de~(SGA~2 XIV~1.20). Mais, en un point g\'eom\'etrique
$\overline{x}'_{1}$ au-dessus d'un point maximal de~$Y'_{1}$,
$X'_{1\red}$ est le spectre d'un anneau de valuation discr\`ete, et
le fait que $\overline{V'_{1}}$ se prolonge \`a $\overline{X'_{1}}$
r\'esulte alors de~X~\ref{X.3.6}.

Montrons que les $n_{i}$ sont premiers \`a $p$. En effet, s'il n'en
\'etait pas ainsi, on aurait un indice $i_{0}\in I(x)$ tel que $p$
divise $n_{i_{0}}$. Quitte \`a remplacer $X$ par
$X_{1}[T_{i_{0}},T_{i}]_{i\in I(x)}\big/
(T_{i}^{n_{i_{0}}/p}-f_{i_{0}},T_{i}^{n_{i}}-f_{i})$, on se ram\`ene
au cas o\`u $X'_{1}=X_{1}[T]/T^{p}-f_{i_{0}}$. D'apr\`es ce qui
pr\'ec\`ede le rev\^etement
\'etale $V'_{1}$ de~$U'_{1}$ se prolonge en un rev\^etement
\'etale $E'_{1}$ de~$X'_{1}$. Soit $\eta$ le point ferm\'e de
$S$; comme le morphisme $X'_{1\eta}\to X_{1\eta}$ est radiciel,
$V_{1\eta}$ se prolonge en un rev\^etement \'etale $E_{1\eta}$ de
$X_{1\eta}$. On en d\'eduit alors, comme dans~\ref{XIII.5.2}, que
$E'_{1}$ est muni d'une donn\'ee de descente relativement au
morphisme $X'_{1}\to X_{1}$, qui prolonge la donn\'ee de descente
naturelle que l'on a sur~$E'_{1}|U'_{1}$. Il en r\'esulte que
$V_{1}$ se prolonge \`a $X_{1}$; mais ceci entra\^ine que l'on a
$n_{i_{0}}=1$ contrairement \`a l'hypoth\`ese $n_{i_{0}}=p$.

\begin{corollaire}
\label{XIII.5.6}
Soient $X$ un $S$-sch\'ema, $D=\sum_{1\leqslant i\leqslant
r}\divisor f_{i}$ un diviseur \`a croisements normaux relativement
\`a~$S$, comme dans~\ref{XIII.5.4}. Soient $\overline{x}$ un point
g\'eom\'etrique de~$X$, $\overline{X}$ le localis\'e strict de
$X$ en $\overline{x}$, $\overline{Y}=Y_{(\overline{X})}$,
$\overline{U}=\overline{X}-\overline{Y}$ et
$$
\widetilde{U}= \varprojlim_{(n_{i})}
\overline{U}[T_{i}]_{i\in I(x)}\big/ (T_{i}^{n_{i}}-f_{i}) \quoi,
$$
la limite projective \'etant prise suivant l'ensemble filtrant des
familles d'entiers $n_{i}>0$, premiers \`a la caract\'eristique
$p$ de~$\kres(x)$. Alors $\widetilde{U}$ est un rev\^etement universel
mod\'er\'ement ramifi\'e de~$\overline{U}$ relativement \`a~$S$. Par suite le groupe fondamental mod\'er\'ement ramifi\'e
de~$\overline{U}$ est
$$
\pi_{1}^{\tame}(\overline{U})\simeq \prod_{\ell\neq
p}\ZZ_{\ell}[1]^{I(x)} \qquad \textrm{(isomorphisme canonique)\quoi.}
$$
Le groupe $\pi_{1}^{\tame}(\overline{U})$ est isomorphe non
canoniquement \`a $\prod_{\ell\neq p}\ZZ_{\ell}^{I(x)}$.
\end{corollaire}

\begin{subremarque}
\label{XIII.5.6.1}
Soient $X$ un $S$-sch\'ema, $D=\sum_{1\leqslant i\leqslant
r}\divisor f_{i}$ un diviseur \`a croisements normaux relativement
\`a~$S$, comme dans~\ref{XIII.5.4}, $U=X-\Supp D$.
Pour toute partie $I\subset[1,r]$ soit
$$
X_{I}= \Big(\tbigcap_{i\in I}V(f_{i})\Big) \cap \Big(\tbigcap_{i\in\complement
I}X_{f_{i}}\Big)\quoi .
$$
Soit $p$ un entier premier ou nul et soit $Z$ un sous-ensemble de
$X_{I}$ dont tous les points sont de caract\'eristique~$p$. Soit
$$
\widetilde{U}_{I}= \varprojlim_{(n_{i})}
U[T_{i}]_{i\in I}
\big/ (T_{i}^{n_{i}}-f_{i}) \quoi,
$$
o\`u
la limite projective est prise suivant l'ensemble filtrant des
familles d'entiers $n_{i}>0$, premiers \`a $p$. Alors, pour tout
point g\'eom\'etrique $\overline{x}$ de~$Z$, l'image inverse de
$\widetilde{U}_{I}$ sur~$\overline{U}$ s'identifie au rev\^etement
universel mod\'er\'ement ramifi\'e de~$\overline{U}$.
\end{subremarque}

\begin{corollaire}
\label{XIII.5.7}
Les notations sont celles de~\ref{XIII.5.6}. Soient $\overline{S}$ le
localis\'e strict de~$S$ en~$\overline{x}$,
$$
\overline{g}\colon\overline{U}\to\overline{S} \quad \text{et} \quad
\tilde{g}\colon\widetilde{U}\to\widetilde{S}
$$
les morphismes canoniques. Alors les morphismes $\overline{g}$ et
$\tilde g$ sont $0$-acycliques \textup{(SGA~4 XV~1.3)}. Soient $G$ un faisceau
en groupes constructible sur~$\overline{S}$, $F=\overline{g}^*G$, $P$
un torseur sous $F$. Alors, pour que $P$ soit mod\'er\'ement
ramifi\'e sur~$\overline{X}$ relativement \`a $\overline{S}$, il
faut et il suffit que son image inverse $\widetilde{P}$ sur~$\widetilde{U}$ soit triviale.
\end{corollaire}

En effet, pour tout sch\'ema
$\overline{X'}=\overline{X}[T_{i}]_{i\in I(x)}\big/
(T_{i}^{n_{i}}-f_{i})$, o\`u les $n_{i}$ sont des entiers~$>0$
premiers \`a $p$, le morphisme
$\overline{f}'\colon\overline{X'}\to\overline{S}$ est $0$-acyclique.
Les fibres g\'eom\'etriques de~$\overline{f}'$ aux diff\'erents
points de~$\overline{S}$ sont donc connexes et m\^eme
irr\'eductibles. Il en est donc de m\^eme des fibres
g\'eom\'etriques des morphismes
$\overline{g}'\colon\overline{U'}\to\overline{S}$, ce qui prouve que
les $\overline{g}'$, donc aussi $\tilde g$, sont
$0$-acycliques (SGA~4 XV~1.16).

Il est clair qu'un torseur $P$ sur~$\overline{U}$ de groupe $F$, dont
l'image inverse sur~$\widetilde{U}$ est triviale, est
mod\'er\'ement ramifi\'e sur~$\overline{X}$ relativement \`a
$\overline{S}$. Montrons que r\'eciproquement, si $P$ est
mod\'er\'ement ramifi\'e sur~$\overline{X}$ relativement \`a
$\overline{S}$, son image inverse sur~$\widetilde{U}$ est triviale.

Il r\'esulte de (SGA~4 IX~2.14~(ii)) que l'on peut trouver un
morphisme fini $n\colon S_{1}\to\overline{S}$, un faisceau en groupes
constant $C$ sur~$S_{1}$, un monomorphisme $G\to n_*C$.
Consid\'erons le diagramme commutatif suivant form\'e de
carr\'es cart\'esiens:
$$
\xymatrix{{\widetilde{U}_{1}}\ar[r]\ar[d]_{r}&{U_{1}}\ar[r]^{g_{1}}\ar[d]_{q}&{S_{1}}\ar[d]_{n}\\{\widetilde{U}}\ar[r]&{\overline{U}}\ar[r]^{\overline{g}}&{\overline{S} \rlap{\quoi.}}\\ }
$$
Soient
$C_{1}$ (\resp $\widetilde{C}_{1}$) l'image inverse de~$C$ sur~$U_{1}$
(\resp sur~$\widetilde{U}_{1}$). On a un diagramme commutatif, dans
lequel $i$ et $j$ sont des isomorphismes (SGA~4 VIII~5.8):
\begin{equation*}
\label{eq:XIII.5.7.*}
\tag{$*$}
\begin{array}{c}
\xymatrix{{\H^{1}(\overline{U},q_*C_{1})}\ar[r]^{i}\ar[d]&{\H^{1}(U_{1},C_{1})}\ar[d]\\{\H^{1}(\widetilde{U},r_*\widetilde{C}_{1})}\ar[r]^{j}&{\H^{1}(\widetilde{U}_{1},\widetilde{C}_{1}) \rlap{\quoi.}}\\ }
\end{array}
\end{equation*}
Soit $Q$ le torseur sous $q_*C_{1}$ d\'eduit de~$P$ par
l'extension du groupe structural $F\to q_*C_{1}$.
D'apr\`es~\ref{XIII.2.1.4}, $Q$ est mod\'er\'ement ramifi\'e
sur~$\overline{X}$ relativement \`a $\overline{S}$. Au torseur $Q$
correspond, gr\^ace \`a $i$, un torseur $Q_{1}$ sous $C_{1}$, et
il est clair que $Q_{1}$ est mod\'er\'ement ramifi\'e sur~$X_{1}=\overline{X}\times_{\overline{S}}S_{1}$ relativement \`a~$S_{1}$. Il r\'esulte donc de~\ref{XIII.5.6} que l'image inverse
$\widetilde{Q}_{1}$ de~$Q_{1}$ sur~$\widetilde{U}_{1}$ est triviale,
et le diagramme~\eqref{eq:XIII.5.7.*} montre alors que l'image inverse
$\widetilde{Q}$ de~$Q$ sur~$\widetilde{U}$ est triviale.

Consid\'erons le diagramme commutatif suivant, dont la deuxi\`eme
ligne est exacte (SGA~4 XII~3.1):
\begin{equation*}
\label{eq:XIII.5.7.**}
\tag{$**$}
\begin{array}{c}
\xymatrix{{\H^{0}(\overline{S},n_*C/G)}\ar[r]\ar[d]_{k}&{\H^{1}(\overline{S},G)\rlap{$~=~1$}}\ar[d]\\{\H^{0}(\widetilde{U},r_*\widetilde{C}_{1}/\widetilde{F})}\ar[r]&{\H^{1}(\widetilde{U},\widetilde{F})}\ar[r]&{\H^{1}(\widetilde{U},r_*\widetilde{C}_{1}) \rlap{\quoi.}}\\ }
\end{array}
\end{equation*}
Comme le morphisme $\widetilde{U}\to\overline{S}$ est $0$-acyclique,
$k$ est un isomorphisme. Le fait que $\widetilde{P}$ soit trivial
r\'esulte alors de~\eqref{eq:XIII.5.7.**}.

\section[Appendice II: finitude pour les images
directes des champs]{Appendice II: th\'eor\`eme de finitude pour les images
directes des champs}
\label{XIII.6}

\begin{proposition}
\label{XIII.6.1}
Soient $S$ un sch\'ema localement noeth\'erien, $f\colon X\to S$
un morphisme. Si $S'$ est un $S$-sch\'ema, on note $X'$
(\resp $f'$, etc.) l'image inverse de~$X$ (\resp $f$, etc.) par le
morphisme $S'\to S$. Supposons que, pour tout
sch\'ema $S'$ \'etale sur~$S$, pour tout faisceau d'ensembles
constructible $F$ sur~$X'$, $f'_*F$ soit constructible, et que, pour
tout faisceau en groupes constructible $F$ sur~$X'$, $\R^{1}f'_*F$
soit constructible. Soit $\Phi$ un champ $1$-constructible sur~$X$~\eqref{XIII.0}. Alors $f_*\Phi$ est $1$-constructible.
\end{proposition}

Pour tout sch\'ema $S'$ \'etale sur~$S$ et pour tout objet $x$ de
$(f_*\Phi)_{S'}$, on a un isomorphisme
$$
\SheafAut_{S'}(x)\simeq f'_*\SheafAut_{X'}(x) \quoi,
$$
o\`u, dans le deuxi\`eme membre de l'\'egalit\'e, $x$ est
consid\'er\'e comme objet de~$\Phi_{X'}$. Les hypoth\`eses
faites entra\^inent donc que $f_*\Phi$ est constructible.
Soit $S\Phi$ le faisceau des sous-gerbes maximales de
$\Phi$~\cite[III~2.1.7]{XIII.1}. Comme $f_*(S\Phi)$ est constructible, on
peut lui appliquer (SGA~4 IX~2.7), et le fait que $f_*\Phi$ soit
$1$-constructible r\'esulte alors du lemme qui suit.

\begin{sublemme}
\label{XIII.6.1.1}
Soient $S$ un sch\'ema localement noeth\'erien, $f\colon X\to S$
un morphisme, $\Phi$ un champ sur~$X$. On suppose donn\'e un faisceau
sur~$S$, repr\'esentable par un $S$-sch\'ema \'etale de type
fini $T$, un morphisme surjectif
$$
a\colon T\to f_*(S\Phi )
$$
et un objet $p$ de la fibre $\Phi_{X_T}$ (o\`u $X_T=X\times_ST$),
d\'efinissant dans $f_*(S\Phi )(T)= S\Phi (X_T)$ un \'el\'ement
\'egal \`a l'image $q$ par $a$ de la section identique de
$T(T)$. Soit $f_T\colon X_T\to T$ le morphisme canonique et supposons
que le faisceau $\R^1f_{T*}(\SheafAut_{X_T}(p))$ soit constructible;
alors il en est de m\^eme de~$S(f_*\Phi )$.
\end{sublemme}

Le morphisme canonique $f^*f_*\Phi\to\Phi$ donne un morphisme
$$
S(f^*f_*\Phi )\simeq f^*(S(f_*\Phi ))\to S\Phi ,
$$
d'o\`u un morphisme canonique
$$
\varphi\colon S(f_*\Phi )\to f_*(S\Phi ).
$$
Soient
$F=S(f_*\Phi )$ et $G$ l'image de~$F$ par $\varphi$; d'apr\`es
(SGA~4 IX,~2.9) $G$ est un faisceau constructible.

Il suffit de montrer que, pour tout point $s$ de~$S$, il existe un
ouvert non vide~$U$ de~$s$ tel que $F|U$ soit localement constant
constructible. Soient $s\in S$, $\overline{s}$ un point
g\'eom\'etrique au-dessus de~$s$, $\overline{q}_1,\dots,\overline{q}_n$ les \'el\'ements de~$G_{\overline{s}}$. Par
d\'efinition de~$T$, il existe des $S$-morphismes
$h_i\colon\overline{s}\to S'$ tels que l'on ait
$\overline{q}_i=h_i^*(q)$. Soient $S'$ le produit fibr\'e sur~$S$ de
$n$ sch\'emas isomorphes \`a $T$, $\overline{s}\to S'$ le produit
fibr\'e des $h_i$, $X'=X\times_SS'$, $q_i$ (\resp $p_i$) l'image
inverse de~$q$ (\resp $p$) par la $i$-\`eme projection de~$S'$ sur~$T$. Si $\Psi_i$ est la sous-gerbe maximale de~$\Phi|X'$
engendr\'ee par $p_i$, le faisceau
$F_i=\R^1f_*'(\SheafAut_{X'}(p_i))$ n'est autre que le faisceau
$S(f_*'\Psi_i)$ des sous-gerbes maximales de~$f_*'\Psi_i$. En
particulier l'injection canonique $\Psi_i\to\Phi |X'$ donne un
morphisme
$$
\alpha_i\colon F_i\to F|S'.
$$

Nous allons montrer que $\alpha_i$ est une bijection de~$F_i$ sur
l'image inverse de~$q_i$ dans $F|S'$. Pour tout sch\'ema $S''$
\'etale sur~$S'$ toute section $y$ de~$F_i(S'')$ a pour image
$q_i|S''$ dans $F(S'')$, car, localement pour la topologie \'etale
sur~$S''$, $y$ est d\'efini par un objet~$x$ de~$\Phi_{X''}$ qui est
isomorphe \`a $p_i|X''$. R\'eciproquement, si $y\in F(S'')$ a pour
image $q_i|S''$ dans $F(S'')$, localement pour la topologie \'etale
sur~$S''$, $y$ est d\'efini par un objet~$x$ de~$\Phi_{X''}$ qui est
isomorphe \`a $p_i$; par suite $x$ est un objet de~$\Psi_{iX''}$ et
par suite $y\in F_i(S'')$.

La d\'emonstration s'ach\`eve en utilisant~\ref{XIII.6.1.2}
ci-dessous. On peut en effet trouver un voisinage ouvert $U'$ de~$s$
tel que $q_1|U',\dots,q_n|U'$ soient des sections de~$G(U')$ et
engendrent ce faisceau. Comme les $F_i|U'$ et $G|U'$ sont
constructibles, il en est de m\^eme de~$F|U'$
d'apr\`es~\ref{XIII.6.1.2}; quitte \`a remplacer $U$ par un ouvert
plus petit, $F|U$ est localement constant, ce qui ach\`eve la
d\'emonstration.

\begin{sublemme}
\label{XIII.6.1.2}
Soient
$S$ un sch\'ema localement noeth\'erien, $F\to G$ un
morphisme surjectif de faisceaux en groupes sur~$S$. Soient $q_i$ une
famille finie de sections de~$G$ sur~$X$ qui engendrent $G$, et, pour
chaque $i$, soit $F_i$ le sous-faisceau de~$F$ image inverse de
$q_i$. Alors, si $G$ et les $F_i$ sont constructibles, il en est de
m\^eme de~$F$.
\end{sublemme}

Pour prouver que $F$ est constructible, il suffit de montrer que, pour
tout point~$s$ de~$S$, il existe un voisinage ouvert $U$ de~$s$ tel
que $F|U$ soit localement constant constructible. Soit donc $s$ un
point de~$S$. Comme les faisceaux $F_i$ et $G$ sont constructibles, on
peut trouver un voisinage ouvert $U$ de~$s$ tel que $F_i|U$ et $G|U$
soient localement constants. Montrons alors que $F|U$ est localement
constant. D'apr\`es (SGA~4 IX~2.13~(i)), il suffit de voir que, si
$\overline{s}$ est un point g\'eom\'etrique au-dessus de~$s$,
$\tilde{s}$ un point g\'eom\'etrique de~$U$ et
$\overline{s}\to\tilde{s}$ un morphisme de sp\'ecialisation, le
morphisme canonique
$$
F_{\tilde{s}}\to F_{\overline{s}}
$$
est bijectif.

On consid\`ere les diagrammes commutatifs
$$
\xymatrix{{(F_i)_{\tilde{s}}}\ar[r]^-{\tilde{a}_i}\ar[d]^{\wr}&{F_{\tilde{s}}}\ar[r]^-{\tilde{a}}\ar[d]_b&{G_{\tilde{s}}}\ar[d]^{\wr}\\{(F_i)_{\overline{s}}}\ar[r]^-{\overline{a}_i}&{F_{\overline{s}}}\ar[r]^-{\overline{a}}&{G_{\overline{s}}}}
$$
Soit $\overline{q}_i$ (\resp $\tilde{q}_i$) l'image inverse de~$q_i$
dans $G_{\overline{s}}$ (\resp $G_{\tilde{s}}$); les morphismes
$\overline{a}$ et $\tilde{a}$ sont surjectifs, et le morphisme
$\overline{a}_i$ (\resp $\tilde{a}_i$) induit une bijection de
$(F_i)_{\overline{s}}$ (\resp $(F_i)_{\tilde{s}}$) sur~$\overline{a}^{-1}(q_i)$ (\resp sur~$\tilde{a}^{-1}(q_i)$). Il r\'esulte
donc du diagramme ci-dessus que $b$ est un isomorphisme.

\begin{corollaire}
\label{XIII.6.2}
Soient
$S$ un sch\'ema localement noeth\'erien, $f\colon X\to S$
un morphisme propre. Soit $\Phi$ un champ $1$-constructible sur~$X$,
alors $f_*\Phi$ est un champ \hbox{$1$-constructible}.
\end{corollaire}

La d\'emonstration de~\ref{XIII.6.1} prouve aussi le r\'esultat
suivant, compte tenu de \ref{XIII.2.4}~2).

\begin{corollaire}
\label{XIII.6.3}
Soient $S$ un sch\'ema localement noeth\'erien, $f\colon X\to S$ un
morphisme, $D$ un diviseur sur~$X$ \`a croisements normaux
relativement \`a~$S$ \eqref{XIII.2.1}, $Y=\Supp D$, $U=X-Y$,
$i\colon U\to X$ l'immersion canonique. Soit $\Phi$ un champ sur~$U$
donn\'e, localement pour la topologie \'etale sur~$X$ et $S$,
comme image inverse d'un champ $\Psi$ $1$-constructible sur~$S$. Alors
le champ $i_*^{\tame}\Phi$ est $1$-constructible.
\end{corollaire}


\begin{thebibliography}{0}{XIII.7}

\bibitem[1]{XIII.1} J.\, \textsc{Giraud}, Cohomologie non ab\'elienne de
degr\'e~$2$, th\`ese, Paris (1966).

\bibitem[2]{XIII.2} J.\, \textsc{Giraud},
Cohomologie non ab\'elienne, Springer, Berlin-New York (1971).

\bibitem[3]{XIII.3} 
H.\,\textsc{Seifert}, W.\,\textsc{Threlfall}, Lehrbuch der Topologie, Chelsea, New
York (1934).


\bibitem[4]{XIII.4}
J-P.\,\textsc{Serre}, Cohomologie galoisienne, Lecture Notes
\no 5, Springer, Berlin (1964).

\bibitem[5]{XIII.5}
J-P.\,\textsc{Serre}, Corps locaux, Hermann, Paris (1962).
\end{thebibliography}

